% Source presentations for register row N-c-018.

\subsection{Three modal-\(\mu\) source presentations}
\label{sec:p09-presentation-n-c-018}

\subsubsection{Kozen's modal calculus and aconjunctive fragment}
\label{sec:p09-n-c-018-kozen}

Kozen defines two languages. The finitary language \(L_\mu\) has
propositional constants \(P\) (including \(0,1\)), propositional variables
\(X\), program constants \(a\), and the grammar
\[
 p::=X\mid P\mid p\lor q\mid\neg p\mid
       \langle a\rangle p\mid\mu X.p,
\]
where every free occurrence of \(X\) in the body of \(\mu X.p\) lies beneath
an even number of negations.  The infinitary companion \(L_{\mu+}\) adds
\(\alpha X.p\) for an ordinal \(\alpha\), with the same positivity
condition on \(X\). It denotes a specified transfinite approximation
stage rather than the least fixed point. The finite input language
remains \(L_\mu\); ordinal-stage formulas are used in the model
construction.

\paragraph{Environments and denotations.}
A model is \(M=(S,I)\), where \(I(P)\subseteq S\),
\(I(a)\subseteq S\times S\), \(I(0)=\varnothing\), and \(I(1)=S\).
No finiteness or total-transition assumption is imposed here.
Let \(\mathsf{Var}\) be the set of propositional variables.
An environment \(\eta:\mathsf{Var}\to\mathcal P(S)\) assigns
their values; \(\eta[X:=A]\) changes only
the value of \(X\). The basic denotations are
\[
\begin{aligned}
 \llbracket X\rrbracket_\eta&=\eta(X),&
 \llbracket P\rrbracket_\eta&=I(P),\\
 \llbracket p\lor q\rrbracket_\eta
   &=\llbracket p\rrbracket_\eta\cup\llbracket q\rrbracket_\eta,&
 \llbracket\neg p\rrbracket_\eta
   &=S\setminus\llbracket p\rrbracket_\eta,\\
 \llbracket\langle a\rangle p\rrbracket_\eta
   &=\{s\in S:\exists t\in\llbracket p\rrbracket_\eta\
                         (s,t)\in I(a)\}.
\end{aligned}
\]
The derived box consequently satisfies
\[
 [a]p:=\neg\langle a\rangle\neg p,\qquad
 \llbracket[a]p\rrbracket_\eta
  =\{s:\forall t\,((s,t)\in I(a)\Rightarrow
                          t\in\llbracket p\rrbracket_\eta)\}.
\]
At a state without \(a\)-successors every box is true and
every diamond is false. Truth is
\(M,s,\eta\models p\) exactly when
\(s\in\llbracket p\rrbracket_\eta\). A closed formula has
no free propositional variables, and its denotation is
independent of \(\eta\).

\paragraph{Least and greatest fixed points.}
For a monotone map \(F:\mathcal P(S)\to\mathcal P(S)\), put
\[
\begin{aligned}
 L_F&=\bigcap\{A\subseteq S:F(A)\subseteq A\},\\
 U_F&=\bigcup\{A\subseteq S:A\subseteq F(A)\}.
\end{aligned}
\]
These are respectively the least and greatest fixed points.
For every pre-fixed \(A\), monotonicity gives
\(F(L_F)\subseteq F(A)\subseteq A\), hence
\(F(L_F)\subseteq L_F\). It follows that
\(F(F(L_F))\subseteq F(L_F)\), so \(F(L_F)\) is itself
pre-fixed. The defining intersection then gives
\(L_F\subseteq F(L_F)\), proving equality.
Dually, every post-fixed \(A\) satisfies
\(A\subseteq F(A)\subseteq F(U_F)\), giving
\(U_F\subseteq F(U_F)\). Its image is post-fixed by
monotonicity, so \(F(U_F)\subseteq U_F\). The extremal
claims follow directly from the intersection and union.

For the body \(p\) of a binder, hold the other variable
values fixed and form
\[
 F_\eta(A):=\llbracket p\rrbracket_{\eta[X:=A]}.
\]
Positivity makes this map monotone, and the semantic
clause is
\[
 \llbracket\mu X.p\rrbracket_\eta=L_{F_\eta}.
\]
If \(c(A)=S\setminus A\), the map \(cFc\) is monotone
and its pre-fixed sets are exactly the complements of
the post-fixed sets of \(F\). Consequently
\[
 U_F=c(L_{cFc}).
\]
This proves the greatest-fixed-point meaning of
\[
 \nu X.p:=\neg\mu X.\neg\bigl(p[\neg X/X]\bigr),
 \qquad
 \llbracket\nu X.p\rrbracket_\eta=U_{F_\eta}.
\]
All displayed substitutions replace free occurrences
and avoid capture \cite[pp.~334--336]{Kozen1983}.

\paragraph{Polarity, parameters, and substitution.}
The monotonicity argument also covers nested binders.
First note that if \(F(A)\subseteq G(A)\) for every
\(A\subseteq S\), then
\[
 L_F\subseteq L_G,\qquad U_F\subseteq U_G.
\]
Indeed, \(L_G\) is pre-fixed for \(F\), and \(U_F\)
is post-fixed for \(G\).
Structural induction now proves two polarity assertions:
if all free occurrences of \(X\) in \(p\) are positive,
increasing \(\eta(X)\) can only increase its denotation;
if all are negative, it can only decrease it.
Constants and variables give the base cases. Union and
the diamond preserve inclusion; complement reverses it
and exchanges the two assertions. At a binder for a
different variable \(Y\), the body induction compares
the induced operators pointwise at every \(Y\)-value.
The fixed-point comparison just proved gives the required
inclusion between their least or greatest fixed points.
For an ordinal-stage binder, the same pointwise
operator comparison gives inclusion between equally
indexed approximants: their initial sets agree;
the successor step uses the body comparison and
monotonicity; the limit step takes unions.
This is a transfinite induction on the stage ordinal.
A binder for \(X\) removes dependence on its outside
environment value. This completes the simultaneous
polarity induction and justifies every admitted binder.

For a formula \(q\) free for \(X\) in \(p\), the
substitution lemma is
\[
 \llbracket p[q/X]\rrbracket_\eta
   =\llbracket p\rrbracket_{
         \eta[X:=\llbracket q\rrbracket_\eta]}.
\]
Being free for \(X\) means that no free variable of an
inserted \(q\) becomes captured by an enclosing binder.
Rename bound variables first when necessary.
The lemma is proved by structural induction. Its
Boolean and modal clauses commute with interpretation.
At a binder for \(Y\ne X\), freshness ensures that
the value of \(q\) does not change when \(Y\)'s
approximating set changes. The body induction therefore
identifies the two induced operators at every set;
their fixed points are equal. For \(\alpha Y.p\),
transfinite induction instead identifies their
approximants at every stage up to \(\alpha\).
At a binder for \(X\),
substitution stops and the outside value of \(X\)
has no effect. These are all the cases.

\paragraph{Transfinite approximants and stabilization.}
For a monotone \(F\), define
\[
\begin{aligned}
 A_0&=\varnothing,& A_{\alpha+1}&=F(A_\alpha),&
 A_\lambda&=\bigcup_{\beta<\lambda}A_\beta,\\
 B_0&=S,& B_{\alpha+1}&=F(B_\alpha),&
 B_\lambda&=\bigcap_{\beta<\lambda}B_\beta
                  \quad(\lambda\text{ a limit ordinal}).
\end{aligned}
\]
The first chain is increasing and the second decreasing.
For example, at an increasing limit stage, each
\(A_{\beta+1}\) lies in \(F(A_\lambda)\); taking their
union gives \(A_\lambda\subseteq F(A_\lambda)\).
The successor step follows by monotonicity, and
complement duality gives the descending case.

Both chains stabilize at some ordinal. Otherwise
their stages would be pairwise distinct subsets of
the set \(S\) for arbitrarily large ordinals. Taking
an ordinal that cannot inject into \(\mathcal P(S)\)
rules this out.
Once \(A_{\kappa+1}=A_\kappa\), every later stage is
unchanged and \(A_\kappa\) is fixed. Transfinite
induction puts every \(A_\alpha\) inside every
pre-fixed set, so \(A_\kappa=L_F\). Dually the
stabilized descending chain is \(U_F\).
Thus the notation \(\bigcup_\alpha A_\alpha\)
for the least fixed point abbreviates a stabilized
set-valued union; no proper-class union is needed.
If \(S\) is finite, there are at most \(|S|\) strict
changes in either chain.

The least \(\kappa\) with
\(A_{\kappa+1}=A_\kappa\) is the closure ordinal of
this particular operator. For the binder body \(p\),
\(\llbracket\alpha X.p\rrbracket_\eta=A_\alpha\);
the fixed-point formula \(\mu X.p\) denotes the
stabilized value. Closure ordinal is an ordinal
measuring iteration, whereas \(\operatorname{CL}(p_0)\)
below is a finite set of formulas.

\paragraph{Closed source formulas and active variables.}
Fix a closed input \(p_0\) in positive normal form: negations apply
only to proposition constants, and binders have been renamed so that
no variable is bound twice. Write \(\operatorname{Sub}(p_0)\) for its
raw syntactic subformulas, retaining their positions and binder links.
Kozen orders containment with the containing formula on the left:
\[
 p\leq q\quad\Longleftrightarrow\quad q\in\operatorname{Sub}(p),
 \qquad
 p<q\quad\Longleftrightarrow\quad
 q\in\operatorname{Sub}(p)\setminus\{p\}.
\]
Abbreviate the unique binding formula for \(X\) by \(\sigma X\),
where \(\sigma\in\{\mu,\nu\}\); thus \(\sigma X\) abbreviates
an entire formula \(\sigma X.u\), not just its binder symbol.
Definition~3.2.1 lists all variables whose binders properly contain \(p\):
\[
 V_p=(Y_1,\ldots,Y_n),\qquad
 \sigma Y_1<\cdots<\sigma Y_n<p.
\]
The list runs from outermost to innermost and includes variables
that are neither free nor active in \(p\). For an ordered subsequence
\(\bar X=(X_1,\ldots,X_k)\) and formulas
\(\bar q=(q_1,\ldots,q_k)\), use the book's
replacement-before-variable notation:
\[
 p[\bar q/\bar X]
   :=p[q_k/X_k][q_{k-1}/X_{k-1}]\cdots[q_1/X_1].
\]
The rightmost-listed, innermost variable is replaced first;
these are sequential capture-avoiding substitutions, not a simultaneous
substitution. The empty list does nothing. The source writes
\(p[\bar X/\bar q]\) for this same operation.
Definition~3.2.2 is therefore
\[
 \sigma V_p=(\sigma Y_1,\ldots,\sigma Y_n),\qquad
 e(p)=p[\sigma V_p/V_p],\qquad
 \operatorname{CL}(p_0)=\{e(p):p\in\operatorname{Sub}(p_0)\}.
\]

Every \(e(p)\) is closed. Indeed, every free variable of a raw
subformula has an enclosing binder. Replacing an inner variable by
its binding formula introduces only variables bound farther outside;
the later substitutions close those parameters. Consequently
\[
 |\operatorname{CL}(p_0)|
   \leq|\operatorname{Sub}(p_0)|\leq|p_0|.
\]
Proposition~3.2.3 characterizes this image as the least set \(D\)
of closed formulas satisfying all seven conditions
\[
\begin{array}{rll}
 (i)&p_0\in D,\\
 (ii)&\neg P\in D&\Longrightarrow P\in D,\\
 (iii)&p\lor q\in D&\Longrightarrow p,q\in D,\\
 (iv)&p\land q\in D&\Longrightarrow p,q\in D,\\
 (v)&\langle a\rangle p\in D&\Longrightarrow p\in D,\\
 (vi)&[a]p\in D&\Longrightarrow p\in D,\\
 (vii)&\sigma X.u\in D&\Longrightarrow u[\sigma X.u/X]\in D.
\end{array}
\]
Here is the closure and leastness proof. The map \(e\) fixes
closed subformulas and commutes with Boolean and modal constructors;
in particular \(e(\neg P)=\neg P\). For a binder
\(b=\sigma X.u\), abbreviate its ordered outer substitution by
\(S=[\sigma V_b/V_b]\). Then
\[
 e(b)=\sigma X.(uS),\qquad e(X)=e(b),\qquad
 e(u)=u[b/X]S=(uS)[e(b)/X].
\]
To check the middle identity, substitutions for binders between
\(b\) and the raw occurrence \(X\) leave that variable alone;
the substitution for \(X\) inserts \(b\), and the remaining
outer substitutions give \(e(b)\). The body identity follows
because its enclosing list appends \(X\) to \(V_b\).
The globally distinct binder names make each indicated replacement
free for its position. These identities prove that the image obeys
(i)--(vii). Conversely, if \(D\) obeys those conditions, descend
the raw syntax tree from \(p_0=e(p_0)\). The Boolean and modal
component rules put each child's closed image in \(D\); (vii)
and the body identity do the same at a binder; (ii) handles a
negated constant, and a variable is a leaf. Thus every \(e(p)\)
belongs to \(D\), proving leastness, including vacuous binders.
This finite vocabulary is not closure under arbitrary formula formation
and is not identified with the PDL closure. The tableau below uses
the raw subformulas and binder links, not these closed images
\cite[defs.~3.2.1--3.2.2, prop.~3.2.3, pp.~336--337]{Kozen1983}.

Let \(\operatorname{FV}(p)\) denote the free propositional variables
of \(p\). For a variable \(Y\) and raw subformula \(p\), let
\(W_{Y,p}\) be the subsequence of \(V_p\) comprising the
variables \(Z\) with \(\sigma Y<\sigma Z<p\).
Definition~3.3.1 says exactly
\[
 Y\in A_p\quad\Longleftrightarrow\quad
 \sigma Y<p\ \text{ and }\
 Y\in\operatorname{FV}\bigl(p[\sigma W_{Y,p}/W_{Y,p}]\bigr).
\]
Such a \(Y\) is \emph{active in \(p\)}. The ordered list
\(A_p\) is a subsequence of \(V_p\); write
\(A_{\mu p}\) and \(A_{\nu p}\) for its respective
\(\mu\)- and \(\nu\)-variable subsequences. All retain the
outer-to-inner order. For a raw occurrence \(X\),
\[
 A_X=A_{\sigma X}\mathbin{\cdot}(X),
\]
where the dot appends an entry. Activity is a syntactic relation
between a bound variable and a source subformula, not a semantic
value or mere free occurrence.

Lemma~3.3.3 states its transitivity:
\[
 Y\in A_p\ \text{ and }\ X\in A_{\sigma Y}
 \quad\Longrightarrow\quad X\in A_p.
\]
For a proof, list all intervening binders in the order
\[
 \sigma X<\sigma X_1<\cdots<\sigma X_r<\sigma Y
   <\sigma Y_1<\cdots<\sigma Y_s<p,
\]
and put \(\bar X=(X_1,\ldots,X_r)\),
\(\bar Y=(Y_1,\ldots,Y_s)\). Sequential substitution gives
\[
\begin{split}
 &p[\sigma\bar Y/\bar Y]
   [\sigma Y[\sigma\bar X/\bar X]/Y]
   [\sigma\bar X/\bar X]\\
 &\hspace{8mm}=
 p[\sigma\bar Y/\bar Y][\sigma Y/Y]
   [\sigma\bar X/\bar X].
\end{split}
\]
The first factor on the left has a free \(Y\), and the
formula inserted there has a free \(X\), by the two hypotheses.
The remaining substitutions neither bind nor replace \(X\).
The right side therefore has a free \(X\) after precisely all
intervening binders have been substituted, as required.
We will also use the resulting prefix fact: if \(X\in A_p\),
the portion of \(A_p\) through \(X\) is exactly \(A_X\).
For an enclosing \(Y\), one direction is the lemma. Conversely,
any free \(Y\) exposed below \(\sigma X\) comes from its
body or an intervening binder inserted into that body, so the
same outer substitutions expose \(Y\) in \(\sigma X\).
Thus retaining an active \(X\) retains this entire preceding
active-variable prefix \cite[pp.~337--338]{Kozen1983}.

For example, in the source formula
\[
 \mu X.\nu Y.
 \bigl(X\land\mu Z.(\langle a\rangle Y\lor[b]Z)\bigr),
 \tag{3.3.2}
\]
\(X\) is not free in \(\langle a\rangle Y\).
The intervening \(Z\)-substitution leaves this formula unchanged,
but replacing \(Y\) by its \(\nu Y\)-binding formula exposes
a free \(X\). Hence \(X\) is active there, and also in the
right conjunct of the displayed conjunction.

Definition~3.3.4 now supplies the exact restriction announced in the
abstract.  A positive-normal formula is \emph{aconjunctive in a
\(\mu\)-variable \(X\)} when, at every subformula \(p\land q\), \(X\) is
active in at most one of \(p,q\); it is \emph{aconjunctive} when this holds
for every \(\mu\)-variable.  Aconjunctivity is a syntactic property of a
formula, not a model condition and not a prohibition on conjunction.  It is
nonredundant with positivity and continuity: positivity provides monotonicity,
whereas aconjunctivity prevents one active least-fixed-point obligation from
being copied into both conjuncts.  Formula (3.3.2) is excluded because \(X\)
is active in both sides of its displayed conjunction, while
\[
 \mu X.(P\lor\langle a\rangle X)
\]
is permitted.  The distinction is used at the tableau conjunction rule:
printed p.~344 says that the \(\mu\)-counter lists then split disjointly;
without the restriction a counter would have to be duplicated
\cite[def.~3.3.4, pp.~338--339, 344]{Kozen1983}.

\paragraph{A closure ordinal beyond \(\omega\).}
A monotone \(F\) is \(\omega\)-continuous when, for every
increasing sequence \(A_0\subseteq A_1\subseteq\cdots\),
\[
 F\left(\bigcup_{n<\omega}A_n\right)
       =\bigcup_{n<\omega}F(A_n).
\]
For its increasing approximant chain this identity gives
\(F(A_\omega)=\bigcup_n A_{n+1}=A_\omega\), so the
least fixed point is reached by stage \(\omega\).
Positivity and aconjunctivity do not imply this continuity.

An explicit version of Kozen's Figure~1 has states
\[
 S=\{r\}\cup\{(n,j):n\in\mathbb N,\ 0\leq j\leq n\},
\]
edges \(r\xrightarrow{a}(n,n)\) for every \(n\), and
\((n,j+1)\xrightarrow{a}(n,j)\) for \(j<n\).
There are no other \(a\)-edges. The operator for
\(\mu X.[a]X\) is
\[
 F(A)=\{s:\text{every }a\text{-successor of }s\text{ lies in }A\}.
\]
Induction gives
\[
 A_k=\{(n,j)\in S:j<k\}\qquad(k<\omega).
\]
The base stage is empty; at each successor stage
the terminal states and exactly the next heights
enter. The root stays out at every finite stage
because some chain head has not yet entered.
Consequently
\[
 A_\omega=S\setminus\{r\},\qquad
 A_{\omega+1}=S=A_{\omega+2}.
\]
The closure ordinal is exactly \(\omega+1\).
Moreover, \(r\) lies in \(F(\bigcup_k A_k)\) but
in none of the \(F(A_k)\), directly disproving
\(\omega\)-continuity. The formula has no conjunction
and is aconjunctive \cite[fig.~1, pp.~339--340]{Kozen1983}.

More generally, membership in \(\mu X.[a]X\) carries
an ordinal rank: take the first approximant containing
the state. This stage is a successor ordinal, and
every \(a\)-successor has strictly smaller rank.
An infinite \(a\)-path would give an infinite
descending sequence of ordinals, which is impossible.
Conversely, outside the fixed point each state has
an \(a\)-successor outside it. Dependent choice then
produces an infinite path. Under this usual set-theoretic
principle, the formula denotes exactly the states
from which no infinite \(a\)-path begins.
In the displayed countable model the same conclusion
is immediate: every path chooses one finite chain
and then decreases its integer height until it stops.

\paragraph{The PDL comparison.}
Choosing \(X\) fresh for \(p\), the familiar PDL cases are
\[
 \langle a^*\rangle p
   \equiv\mu X.(p\lor\langle a\rangle X),
 \qquad
 [a^*]p
   \equiv\nu X.(p\land[a]X).
\]
For the diamond identity, finite iteration from the empty
set collects the states with an \(a\)-path of bounded
finite length to \(p\). Its union is exactly the
reflexive-transitive reachability interpretation of
\(\langle a^*\rangle p\).
For the box identity, fix an environment and put
\(P_\eta=\llbracket p\rrbracket_\eta\). Let \(W\)
consist of the states whose every finite \(a\)-reachable
state lies in \(P_\eta\). It satisfies
\[
 W=P_\eta\cap\{s:\forall t\,((s,t)\in I(a)\Rightarrow t\in W)\}.
\]
The length-zero path checks \(p\) now, and a positive
path begins with an \(a\)-edge. Every post-fixed set of this operator
stays inside \(p\) along each finite path by induction
on path length, so it is contained in \(W\).
Thus \(W\) is the greatest fixed point.

Proposition~4.1 states that \(\mu X.[a]X\) is not
equivalent to any PDL formula. Suppose \(\varphi\)
were such an equivalent. The root of the displayed
well-founded model satisfies \(\varphi\).
Filter this model through the finite Fischer--Ladner
closure of \(\varphi\). The PDL filtration theorem
preserves \(\varphi\) at the quotient root, and its
transition relation contains the images of the
original \(a\)-edges. Original paths from the root
have arbitrarily large finite length. Their images
in the finite quotient therefore contain a reachable
repeated state, giving an \(a\)-cycle. Repeating that
cycle yields an infinite path from the quotient root.
It cannot satisfy \(\mu X.[a]X\), contradicting the
assumed equivalence. This proves the strict expressive
separation \cite[prop.~4.1]{Kozen1983}.

For the single-model comparison in Proposition~4.2,
number the dots of Figure~2 by \(\mathbb N\) from left
to right. Work with one program \(a\), set
\[
 I(a)=\{(n+1,n):n\geq0\}
       \cup\{(2k+3,2k+1):k\geq0\},
\]
and interpret every ordinary proposition constant
as false, with \(I(0)=\varnothing\) and \(I(1)=\mathbb N\).
This specifies the otherwise unlabelled frame as a
complete model. There are no self-loops, and \(0\)
has no successor.

The operator of \(\mu X.[a]\langle a\rangle X\) is
\[
 F(B)=\{s:\forall t\,((s,t)\in I(a)\Rightarrow
                    \exists u\,((t,u)\in I(a)\land u\in B))\}.
\]
Its least approximants are
\[
 A_j=\{2r:0\leq r<j\}\quad(j<\omega).
\]
To verify the successor step, \(0\) always enters
vacuously. An even state \(2r>0\) has just the
successor \(2r-1\). That odd state can reach
\(2r-2\) and possibly the preceding odd state;
only \(2r-2\) can belong to an even approximant.
Thus \(2r\) enters exactly when \(2r-2\) has entered.
The odd state \(1\) fails because its successor \(0\)
is dead. Every larger odd state has an even successor
whose only successor is odd, so it also fails
against an even approximant. These observations
prove the formula for \(A_j\) by induction.
The union \(A_\omega\) is the even set and is fixed,
giving
\[
 \llbracket\mu X.[a]\langle a\rangle X\rrbracket
       =\{0,2,4,\ldots\}.
\]

In this model, Proposition~4.2 states that every PDL
formula, including programs with tests and converse,
defines a finite or cofinite subset of \(\mathbb N\)
\cite[prop.~4.2 and fig.~2, pp.~340--341]{Kozen1983}.
The restriction on proposition interpretations matters:
an arbitrary primitive proposition labelled by the
even set would itself violate that conclusion.

\paragraph{Finite/cofinite sets under regular programs.}
\label{sec:p09-kozen-finite-cofinite}
We prove the claim including arbitrary finite nesting
of program iteration, tests, and converse. Write \(b\)
for the converse of \(a\), not for an independently
chosen relation. In an auxiliary frame on all of
\(\mathbb Z\), let \(a\) have every downward unit
edge and every downward skip of length two between
odd integers; \(b\) reverses these edges. A realization
of a word in \(a,b\) is a finite walk spelling its
letters in order. Every unit step is available,
and any walk may be translated by an even integer.

\emph{Word lemma.}
Suppose a nonempty word \(w\) has a realization that
ends strictly below its start with the same parity.
There is an integer \(C_w\geq0\) such that, for
every integer lower boundary \(B\geq0\) and every target
\(y\geq B+C_w\), all sufficiently large integers
can reach \(y\) by a power of \(w\) without ever
going below \(B\).

First, both parity types communicate under powers
of \(w\). If \(|w|\) is odd, the all-unit walk
switches parity. If \(|w|\) is even, one of the
first two steps of an all-unit walk has an odd
source. Replace that unit step by the available
skip and keep the others unit steps. The endpoint
parity switches. Thus fixed walks spelling powers
of \(w\) can join either parity type to the other.

The doubled word \(ww\) contains adjacent identical
letters. Otherwise \(w\) would be an even alternating
word, \((ab)^k\) or \((ba)^k\). Direct calculation
shows that \(ab\) is the union of complete relations
on the pairs \(\{2j,2j+1\}\), and \(ba\) is the
corresponding union on \(\{2j-1,2j\}\).
Their positive powers do not enlarge them, so neither
admits a strictly descending same-parity walk.
This contradicts the hypothesis.

Put \(d=\#b(w)-\#a(w)\). An all-unit realization
of \(ww\) has displacement \(2d\). Choose its
starting parity \(t\) so that the source of the
adjacent equal letters is odd. Replacing those two
unit steps by two skips is possible, because the
first skip keeps the state odd. The displacement
changes by \(2\) or \(-2\).
In units of two, the two \(t\)-to-\(t\) walks
therefore have consecutive integer displacements.

The assumed negative walk has cell displacement
\(-c\), \(c\geq1\), at some parity \(u\).
Join \(t\) to \(u\) and back by fixed parity
walks. Inserting sufficiently many copies of the
negative walk gives a \(t\)-to-\(t\) walk with
negative cell displacement \(h\).
Append the same sufficiently large number of
copies of this walk to each of the consecutive
walks just constructed. The resulting two
displacements are \(-r,-(r+1)\) for some \(r\geq1\).
All these walks still spell powers of \(w\).

Every integer \(L\geq r(r-1)\) is a nonnegative
combination of \(r,r+1\): writing \(L=qr+s\),
\(0\leq s<r\), gives
\[
 L=(q-s)r+s(r+1),\qquad q\geq r-1\geq s.
\]
The two fixed descending walks therefore connect
every sufficiently high state of parity \(t\)
to any prescribed lower state of parity \(t\).
Fixed input and output parity walks cover the
other start and target parities.

The construction also has a uniform lower-bound
estimate. Choose \(K_0\geq0\) so that neither of
the two fixed descending walks goes more than
\(K_0\) below its own endpoint. A concatenation
ending at \(z\) stays at least \(z-K_0\), since
its successive endpoints strictly decrease.
For each target parity \(p\), fix an output walk
from \(t\) to \(p\), with displacement \(\delta_p\)
and minimum prefix displacement \(m_p\), including
the empty prefix. To end at \(y\), it starts at
\(z=y-\delta_p\) and stays at least
\(y-\delta_p+m_p\). Choose
\[
 C_w\geq\max_{p\in\{0,1\}}
             \{0,\delta_p+K_0,\delta_p-m_p\}.
\]
Then \(y\geq B+C_w\) keeps the descending walks
and output walk above \(B\). Each fixed input
walk stays above \(B\) from sufficiently high
starts and leaves a large enough even difference
to \(z\) for the displayed integer combination.
This proves the word lemma in the original
half-line frame as well.

\emph{Guarded-automaton lemma.}
Consider a finite automaton with control set \(Q\),
initial control \(q_0\), final controls, and edges
labelled \(a,b\), or \(D?\), where \(D\subseteq
\mathbb N\) is finite or cofinite. A configuration
is \((q,n)\). Action edges move \(n\) through the
frame; a test changes only control and requires
\(n\in D\). An epsilon edge is the everywhere-true
test. Acceptance is a finite execution reaching
a final control. Its set of accepted starting
integers is finite or cofinite.

Choose an integer \(B\geq1\) beyond all exceptions in the
finitely many tests. Above \(B\), each test is
always true or always false. Delete the false
tests and treat the others as epsilon edges,
forming the high-tail control graph.
If it has a path from \(q_0\) to a final control,
fix such a path and realize every action by a
unit step. All sufficiently high starts keep
this finite execution above \(B\), so the
accepted set is cofinite.

Otherwise every accepting execution from a high
start must visit below \(B\). If accepted starts
are bounded, the set is finite. Suppose they
are unbounded. Project configurations above
\(B\) to a weighted graph on
\(Q\times\{0,1\}\), where the second coordinate
is parity. Unit \(a,b\) edges have weights
\(-1,+1\) and switch parity. Skip edges have
weights \(-2,+2\) from odd to odd; epsilon
edges have weight zero and preserve parity.

For each negative simple cycle of this finite
graph, select one base vertex and its word \(w\)
of non-epsilon labels. The cycle supplies a
negative same-parity realization, so the word
lemma applies. There are finitely many such
cycles. Choose an integer \(H\geq B+C_w\) for all their
selected words, and \(H\geq B\); if there are
none, set \(H=B\).
Put \(s=2|Q|\), and choose an accepted start
\[
 x>H+1+2(s-1).
\]
Stop an accepting execution at its first height
below \(H+2\). This prefix stays at least \(H\)
and ends at \(H\) or \(H+1\), since a downward
edge changes height by at most two. Its weighted
projection has displacement below \(-2(s-1)\).

A finite walk decomposes into a simple path
and simple cycles by deleting repeated-vertex
segments and decomposing the removed closed walks.
The simple path contributes at least
\(-2(s-1)\). Hence a negative simple cycle
occurs in this decomposition. All its vertices
occurred on the original prefix, even if the
cycle is no longer a contiguous part of that
prefix. Let \(q\) be its selected base control
and \(y\geq H\) the height at one original
occurrence of that base vertex. Retain the
original accepting suffix from \((q,y)\).

The cycle's control sequence begins and ends
at \(q\); all its tests are true above \(B\).
Every above-\(B\) realization of a power of
its word \(w\) can therefore follow the same
control sequence, inserting its tests at the
prescribed points. It need not follow the
cycle's original parity choices.
Retain also the control prefix from \(q_0\)
to \(q\), but realize its actions by unit
steps. From every sufficiently large new
start \(z\), this prefix stays above \(B\)
and ends at \((q,z+d_0)\), for a fixed
displacement \(d_0\).
The word lemma connects that height to
\(y\geq B+C_w\), still above \(B\).
Appending the retained accepting suffix
proves that every sufficiently large \(z\)
is accepted. The unbounded accepted set is
therefore cofinite, proving the automaton lemma.

\emph{Induction on PDL formulas.}
Finite/cofinite sets are closed under complement
and finite Boolean operations. The proposition
interpretations have this property.
For \(\langle\pi\rangle\varphi\), every formula
tested inside \(\pi\), and the target \(\varphi\),
is a proper formula subexpression. Their
denotations are finite or cofinite by induction.
Push converse to atomic programs using
\[
\begin{aligned}
 (\alpha\cup\beta)^{-1}
   &=\alpha^{-1}\cup\beta^{-1},&
 (\alpha;\beta)^{-1}
   &=\beta^{-1};\alpha^{-1},\\
 (\alpha^*)^{-1}&=(\alpha^{-1})^*,&
 (\psi?)^{-1}&=\psi?,&
 (\alpha^{-1})^{-1}&=\alpha.
\end{aligned}
\]
The result is a regular expression over \(a,b\)
and finitely many tests.
Construct its finite automaton by using one
edge for an atom, epsilon branching for union,
an epsilon join for composition, and an epsilon
zero-iteration route and return edge for star.
Induction on these operations identifies
accepting automaton executions with program
executions, including zero iterations.
Append a test for \(\llbracket\varphi\rrbracket\)
at the exit, leading to a fresh final control.
The accepted starts are exactly
\(\llbracket\langle\pi\rangle\varphi\rrbracket\),
so the automaton lemma applies.
Boxes follow by complement duality.
This proves the finite/cofinite property for
the full stated PDL syntax. The even set is
neither finite nor cofinite, completing the
single-model expressive separation.

\paragraph{The source calculus.}
\label{sec:p09-kozen-calculus}
The deductive system uses equations and the abbreviation
\(p\leq q\) for \(p\lor q=q\).  Besides equational logic and Boolean algebra,
its modal and fixed-point principles are
\[
\begin{array}{ll}
\langle a\rangle X\lor\langle a\rangle Y
   =\langle a\rangle(X\lor Y),&
\langle a\rangle X\land[a]Y
   \leq\langle a\rangle(X\land Y),\\
\langle a\rangle0=0,&
p[\mu X.p/X]\leq\mu X.p,\\
p[Y/X]\leq Y\ \Longrightarrow\ \mu X.p\leq Y.&
\end{array}
\]
The last rule has the side condition that \(Y\) does not occur in
\(p\); it is Park fixed-point induction. Substitution instances
are admitted only when free occurrences remain uncaptured
and the resulting binders retain their positivity conditions.
A formula is \emph{consistent}
when its equality to \(0\) is not derivable.  This is a proof-theoretic
property of a formula, differs from semantic satisfiability until
Theorem~6.3.1 connects them, and is used to choose the nonrejecting tableau
subtree.  Proposition~5.7 records bound-variable change, monotonicity,
unfolding, vacuous binding, and a strengthened induction consequence
\cite[pp.~341--342]{Kozen1983}.

\paragraph{Soundness of the calculus.}
\label{sec:p09-kozen-calculus-soundness}
An equation is valid when its two denotations agree in every
model and environment. An inequality denotes inclusion.
The Boolean axioms hold in the powerset algebra.
Diamond distributes over union because a successor in
a union lies in at least one component, and there is
no successor satisfying the empty predicate.
For the mixed modal inequality, a state satisfying
\(\langle a\rangle p\land[a]q\) has a particular
\(a\)-successor satisfying \(p\), and every such
successor satisfies \(q\). That same witness satisfies
\(p\land q\), proving the inequality.

For the fixed-point axiom, substitution identifies the
left denotation with \(F(L_F)=L_F\); inclusion follows.
For Park induction, fix a model and environment and
write \(B=\eta(Y)\). If the premise holds as a
state-set inclusion, then \(F(B)\subseteq B\).
The least-pre-fixed-point characterization gives
\(L_F\subseteq B\), which is exactly the conclusion.
The premise is an inclusion of predicates on the whole
state space, not merely an implication known at one
selected state.

Reflexivity, symmetry, transitivity, and replacement
of equal expressions preserve semantic equality.
In particular, replacement inside a binder is sound
because validity holds at every value of its bound
variable; the induced operators agree pointwise,
so their fixed points agree. Capture-avoiding
substitution is sound by the substitution lemma.
Induction on a derivation therefore proves validity
of every derived equation.
If a model has a state satisfying \(p\), an equation
\(p=0\) cannot be derivable, since it would make
the denotation of \(p\) empty in that model.
This proves the model-to-consistency implication
used in Theorem~6.3.1. Soundness applies to the
whole well-formed calculus; aconjunctivity restricts
the completeness construction, not this argument.

\paragraph{Derived fixed-point laws.}
Write \(p(t)\) for capture-avoiding substitution of \(t\)
for the designated variable \(X\) in a positive formula
context \(p(X)\). Other free variables are parameters.
Let \(\operatorname{FV}(q)\) denote the free propositional
variables of \(q\), and rename binders away from them
before substitution. Park induction is used in its
formula-instance form
\[
 \frac{p(b)\leq b}{\mu X.p(X)\leq b},
 \qquad X\notin\operatorname{FV}(b).
\]
Its premise is a derivable inequality. The following
derivations give the six laws of Proposition~5.7 with
their variable conditions explicit
\cite[prop.~5.7, pp.~341--342]{Kozen1983}.

\emph{Renaming.}
If \(Y\) is fresh for \(p\), then
\[
 \mu X.p=\mu Y.(p[Y/X]).
\]
Let the two sides be \(a,b\).
The pre-fixed-point axiom for \(b\) gives \(p(b)\leq b\);
induction for the \(X\)-binder yields \(a\leq b\).
Conversely the axiom for \(a\) makes it a pre-fixed
instance of the renamed \(Y\)-body, so induction gives
\(b\leq a\). Boolean antisymmetry proves equality.

\emph{Comparison of bodies.}
If \(p(X)\leq q(X)\) is derivable and both contexts
are positive in \(X\), then
\[
 \sigma X.p(X)\leq\sigma X.q(X),
 \qquad \sigma\in\{\mu,\nu\}.
\]
For \(\mu\), put \(b=\mu X.q(X)\).
Substitution in the premise and the pre-fixed-point
axiom give \(p(b)\leq q(b)\leq b\); induction
yields \(\mu X.p(X)\leq b\).
For \(\nu\), Boolean complement reverses the premise
between the two dual least-fixed-point bodies.
Apply the just-proved \(\mu\)-case and complement
the conclusion. This gives the displayed \(\nu\)-case.

\emph{Monotonicity of a positive context.}
If \(q\leq r\) is derivable, then
\[
 p(q)\leq p(r).
\]
Prove this by structural induction on the context,
together with the reversed inequality for negative
contexts. The Boolean rules preserve or reverse
order according to polarity. Diamond monotonicity
follows from its union axiom: substituting into
\(q\lor r=r\) gives
\(\langle a\rangle q\lor\langle a\rangle r
 =\langle a\rangle r\). Box monotonicity follows by
duality. At a binder for a fresh \(Y\), the body
induction gives an inequality between its two
substituted bodies; comparison of bodies gives
the same order between their fixed points.
A binder for the designated variable ends its
free occurrences. These cases complete the induction.

\emph{Unfolding.}
For \(a=\mu X.p(X)\), the axiom gives \(p(a)\leq a\).
Context monotonicity then gives \(p(p(a))\leq p(a)\).
Using \(p(a)\) as the candidate in Park induction
yields \(a\leq p(a)\). Thus
\[
 p(\mu X.p(X))=\mu X.p(X).
\]
Complement duality gives the same identity for \(\nu\).
The substitutions are free for their positions;
bound-variable renaming ensures this before the
calculation is used.

\emph{Vacuous binding.}
If \(X\notin\operatorname{FV}(q)\), then
\[
 \mu X.q=q.
\]
The axiom gives \(q\leq\mu X.q\), while the reflexive
inequality \(q\leq q\) is the premise for the reverse
Park-induction inequality. The dual \(\nu\)-identity
follows as well.

\emph{Strengthened induction.}
Let \(q\) be independent of the bound variable \(X\).
Then
\[
 p\bigl(\mu X.(q\land p(X))\bigr)\leq q
 \quad\Longrightarrow\quad
 \mu X.p(X)\leq q.
\]
Put \(a=\mu X.(q\land p(X))\). Unfolding gives
\(a=q\land p(a)\). Under the premise \(p(a)\leq q\),
this becomes \(a=p(a)\). In particular \(a\) is a
pre-fixed formula for \(p\), so induction gives
\(\mu X.p(X)\leq a\leq q\).
Independence of \(q\) is used here as a parameter
condition: it remains unchanged during the \(X\)-iteration.

The last law supplies a useful consistency step without
assuming completeness. If \(q\land\mu X.p(X)\) is
consistent, then so is
\[
 q\land p\bigl(\mu X.(\neg q\land p(X))\bigr).
\]
If the latter were derivably zero, Boolean algebra
would give
\(p(\mu X.(\neg q\land p(X)))\leq\neg q\).
Strengthened induction with parameter \(\neg q\)
would yield \(\mu X.p(X)\leq\neg q\), making
the former conjunction derivably zero as well.
This contradicts its consistency.

The distinction between induction and a local
implication can be seen in a two-state model.
Let \(I(a)=\{(s,t)\}\), \(I(P)=\{t\}\), and
\(\eta(Y)=\varnothing\), and put
\(p(X)=P\lor\langle a\rangle X\).
At \(s\), the implication \(p(Y)\to Y\) is true
because its antecedent is false. But the least
approximants for \(p\) are
\(\varnothing,\{t\},\{s,t\},\{s,t\}\), so
\(\mu X.p(X)\to Y\) is false there.
Rule (5.6) requires the predicate inclusion
\(p(Y)\leq Y\), which fails at \(t\); it does
not license that invalid local inference.

\paragraph{The restricted tableau and its theorem.}
Fix an aconjunctive closed positive-normal \(p_0\), with
distinct names for its bound variables. Normalize
\(\neg0\) to \(1\) and \(\neg1\) to \(0\) first.
Keep \(\operatorname{Sub}(p_0)\) as the original
subformula occurrences with their unique binder links.
Let \(\mathcal Q\) be its quotient by equality of
raw formula syntax and binder links. Thus identical
raw formulas with the same links are one formula key
for the duplicate-formula rule, and
\(|\mathcal Q|\leq|\operatorname{Sub}(p_0)|\leq|p_0|\).
In particular, \(X\) is the same raw key whenever
a counter for \(X\) is incremented. These keys
are distinct from occurrence-indexed game controls.
A tableau node \(s\) carries
\(\Gamma_s\subseteq\mathcal Q\), with root
\(\Gamma_{r_0}=\{p_0\}\).
Members of \(\Gamma_s\) need not be closed in isolation:
the binder rules expose their original bodies and variables.
This raw formula family is distinct from
\(\operatorname{CL}(p_0)\), whose members have been closed
by substitution.

Choose a fixed representative occurrence for each key
when writing its enclosing list \(V_p\).
The active list and the ordered closing translations
are independent of this choice. To verify this,
start with the free variables of the common raw
formula and follow an inner-to-outer substitution.
A substitution has an effect only when its variable
is currently free. It then introduces only free
variables of that variable's fixed binding formula,
all bound farther outside its unique binder.
Those outside binders enclose both representative
occurrences. Repeating this argument shows that
every effective replacement and its order are
determined by the raw formula and its binder links.
Additional enclosing variables at one occurrence
never become free and have no effect.
The same argument applies when a least binder is
annotated by a closed formula or assigned an
ordinal stage: neither operation introduces
additional free variables. Hence active lists,
ordinary closing, and the annotated closing
translations used below descend to \(\mathcal Q\),
up to hygienic renaming.
For example, the two \(X\)-occurrences in
\(\mu X.(X\lor\nu Y.X)\) have enclosing lists
\((X)\) and \((X,Y)\), but the \(Y\)-substitution
does nothing at the latter.

The five rules are
\cite[pp.~342--343]{Kozen1983}:
\begin{enumerate}
\item replace \(p\land q\) by both components in one successor;
\item branch on \(p\lor q\), putting one component in each successor;
\item replace \(\sigma X.p\) by its body \(p\) in one successor,
      retaining the original binder link of each occurrence of \(X\);
\item replace a bound variable \(X\) by the body of its unique binder in one
      successor; and
\item for every \(\langle b\rangle p\in\Gamma_s\), create a \(b\)-successor
      containing \(p\) and every \(q\) for which
      \([b]q\in\Gamma_s\).
\end{enumerate}
The disjunction rule is existential in the alternating procedure; all modal
successors are universal obligations.

A counter \(c\) is an identity with a fixed associated variable
\(X(c)\); its value at node \(s\) is \(c(s)\).
Distinct counters may have the same associated variable.
For \(A_p=(X_1,\ldots,X_k)\), its local list is
\[
 C(s,p)=(c_1,\ldots,c_k),\qquad X(c_i)=X_i.
\]
Write \(C_\mu(s,p)\) and \(C_\nu(s,p)\) for its
\(\mu\)- and \(\nu\)-counter subsequences. Put
\(m=|p_0|\) and \(N=2^m\). Denote by \(\mathsf A_T\)
the same procedure with any specified finite counter cap \(T\).
The estimates first displayed below use the source choice \(T=N\);
the completeness construction will use a separately specified
context-counted cap \(M\).
For \(T=N\), before rejection the values lie in
\(\{0,\ldots,N\}\) for \(\mu\)-counters and
\(\{0,1\}\) for \(\nu\)-counters. The global priority list
\[
 G(s)\in\operatorname{Shuffle}
       \bigl(C_\mu(s,p):p\in\Gamma_s\bigr),\qquad
 C(r_0,p_0)=G(r_0)=()
\]
contains each live \(\mu\)-identity once and preserves each
local order; its leftmost entry has highest priority.
The notation \(\operatorname{Shuffle}\) means precisely those
interleavings. These dynamic proof-search priorities are not the
later parity-game priorities.

For an ordered variable list \(A\), let \(L\upharpoonright A\)
retain, in order, exactly the entries of \(L\) whose associated
variables belong to \(A\). Unchanged values are inherited.
The following assignments from \(s\) to a successor \(t\)
are provisional until duplicate-formula collisions and their
consequent resets have been resolved:

Fix an order of the raw input subformulas and an order
of the output slots of each rule. A provisional list's
origin is its source-formula index together with its
output-slot index; inherited lists have a designated
inheritance slot. Process collision groups in destination
formula order and their lists in origin order.
In a \(\nu\)-only collision, the earlier-origin list
is retained. The \(\mu\)-priority comparison below
still determines the winner when \(\mu\)-entries occur.
Apply the prescribed deletion and dependent resets
after each disposal. This convention makes collision
processing deterministic and invariant under renaming
counter identities.
\begin{enumerate}
\item \emph{Binder rule (6.1.6).}
For \(\sigma X.u\) replaced by its raw body \(u\), if
\(X\in\operatorname{FV}(u)\), create a fresh identity \(c\)
with \(X(c)=X\), \(c(t)=0\), and set
\[
 C(t,u)=C(s,\sigma X.u)\mathbin{\cdot}(c).
\]
For a \(\mu\)-binder also set
\(G(t)=G(s)\mathbin{\cdot}(c)\); a new \(\nu\)-counter
is not put in \(G\). If \(X\notin\operatorname{FV}(u)\),
create nothing: \(C(t,u)=C(s,\sigma X.u)\) and
\(G(t)=G(s)\).

\item \emph{Disjunction rule (6.1.7).}
For the successor choosing \(p\) from \(p\lor q\), set
\[
 C(t,p)=C(s,p\lor q)\upharpoonright A_p.
\]
For the other successor use the analogous restriction to \(A_q\).
Delete from that successor's \(G\) every \(\mu\)-counter
removed by the restriction, and apply the deletion resets below.

\item \emph{Conjunction rule (6.1.8).}
For the single successor of \(p\land q\), set
\[
 C(t,p)=C(s,p\land q)\upharpoonright A_p,\qquad
 C(t,q)=C(s,p\land q)\upharpoonright A_q.
\]
Before collisions \(G\) is unchanged. Activity in a conjunction
is the union of activity in its components; aconjunctivity makes
the two \(\mu\)-sublists disjoint, so they partition the old
\(\mu\)-sublist. A \(\nu\)-identity may be shared by both
lists, with one value wherever it occurs. Thus the disjointness
invariant needed for \(G\) concerns \(\mu\)-identities;
shared \(\nu\)-flags retain one common value.

\item \emph{Variable rule (6.1.9).}
If the raw variable \(X\) has binder \(\sigma X.u\), set
\[
 C(t,u)=C(s,X),\qquad c(t)=c(s)+1,
 \qquad X(c)=X.
\]
Here \(c\) is the last counter of \(C(s,X)\).
For a \(\mu\)-variable reset to zero every \(\mu\)-counter
to its right in \(G\). On every local successor list, an
increment or reset of a \(\mu\)-counter resets every
\(\nu\)-counter after it on that list to zero.
For a \(\nu\)-variable this rule is applicable only if
\(c(s)=0\); its new value is then \(1\).
This zero guard controls applicability, not rejection of the
greatest-fixed-point obligation.

\item \emph{Modal rule (6.1.10).}
For the successor demanded by \(\langle b\rangle p\),
whose provisional label is \(\{p,q_1,\ldots,q_n\}\)
with all relevant \([b]q_i\in\Gamma_s\), set
\[
 C(t,p)=C(s,\langle b\rangle p),\qquad
 C(t,q_i)=C(s,[b]q_i)\quad(1\leq i\leq n).
\]
Restrict \(G\) to the identities occurring in these inherited
lists, apply deletion resets, and reset all surviving
\(\nu\)-counters to zero. Each diamond has its own successor;
every box with that same program contributes to it.

\item \emph{Untouched-formula rule (6.1.11).}
If a nonmodal rule acts on \(q\ne p\), while
\(p\in\Gamma_s\), then
\[
 C(t,p)=C(s,p).
\]
Inherited counters still undergo the global and dependent resets
required by (6.1.9) and (6.1.13).

\item \emph{Collision rule (6.1.12).}
If two provisional lists describe the same raw formula, their
associated-variable sequences coincide. When they contain
\(\mu\)-counters, retain the list whose first \(\mu\)-entry
occurs earlier in \(G\); discard the other list and remove its
\(\mu\)-identities from \(G\). The retained list keeps its
values, subject to the prescribed resets: there is no
componentwise maximum in this case. When neither list contains a
\(\mu\)-counter, write them as \(C=(c_1,\ldots,c_k)\)
and \(C'=(c'_1,\ldots,c'_k)\), retain \(C\), and set
\[
 c_i(t)=\max\{c_i(t),c'_i(t)\}\quad(1\leq i\leq k).
\]
Resolve additional collisions in the same way. A discarded list
does not destroy a \(\nu\)-identity still shared by another
retained list. The source writes \(C'\longrightarrow C\)
for discarded-to-retained lists and \(c'\longrightarrow c\)
for corresponding entries with \(X(c')=X(c)\). Such an arrow
need imply neither equality of identities nor equality of values.

\item \emph{Deletion and dependent resets (6.1.13).}
If deleting an earlier \(\mu\)-counter changes a surviving
\(\mu\)-counter's position in \(G\), reset the survivor to
zero. Every increment or reset of a \(\mu\)-counter also
resets to zero all \(\nu\)-counters after it on each local list
containing it. These rules apply to branch restriction, modal
restriction, and collision disposal; a reset acts on the identity
wherever it is shared. New \(\mu\)-identities are appended to
\(G\), and deletion preserves the survivors' relative order.
\end{enumerate}
These are the eight counter assignments, not additional structural
formula rules \cite[eqs.~(6.1.6)--(6.1.13), pp.~343--344]{Kozen1983}.

The procedure rejects a completed configuration if its label
contains \(0\) or both \(P,\neg P\), or if a live
\(\mu\)-counter exceeds \(N\). A rule's provisional lists,
collisions, and resets are resolved before testing its
successor configuration.
Normalization of \(\neg0,\neg1\) and explicit
rejection of \(0\) retain their Boolean meanings.
Otherwise apply an available nonmodal rule, respecting the
\(\nu\)-zero guard. When none is available, perform the modal
action if there are diamond demands; with no applicable rule or
diamond demand, accept. A flag-blocked \(\nu\)-variable is
not an applicable rule, and boxes alone impose no successor demand.

At most \(m\) formulas, each with at most \(m\) local entries,
are present. In particular
\[
 K:=\max_s|G(s)|\leq m^2.
\]
A counter value requires \(O(m)\) bits, since it ranges from
\(0\) through \(2^m\), inclusively. Formula indices, variable
indices, identity references, sharing information, and the order of
\(G\) require no more asymptotic space. A current configuration
therefore has an \(O(m^3)\)-bit encoding
\cite[pp.~344--345]{Kozen1983}.

\emph{Finite increments and resets of a \(\mu\)-identity.}
Fix a live \(\mu\)-counter \(c\). New identities are born
to its right in \(G\), so its rank can decrease only when
earlier identities are deleted, at most \(K-1\) times.
Partition its lifetime into at most \(K\) constant-rank
intervals. In any one interval the prefix of \(G\) ending at
\(c\) has fixed identities. An increment anywhere in that
prefix strictly increases its lexicographic value vector:
higher-priority values stay fixed, the incremented value increases,
and lower-priority values are reset. No deletion before \(c\)
occurs inside the interval. Each digit lies in
\(\{0,\ldots,N\}\), so the prefix has at most
\((N+1)^K\) value vectors.
Every increment or reset of \(c\) inside the interval accompanies
one of these strict increases. Boundary rank changes contribute
only finitely many further resets. Thus each identity has finitely
many increments and resets, bounded in order by
\[
 O\bigl(K((N+1)^K+1)\bigr)=2^{O(m^3)},
\]
also allowing a last rejecting increment. If \(K=0\) there is
no \(\mu\)-identity to count.

This lifetime argument concerns an individual identity. The source's
Definition~6.2.5 instead also follows collision replacements:
if \(c_1,\ldots,c_r\) occupy successive lifetime intervals
\([s_i,t_i]\), with each \(s_{i+1}\) immediately succeeding
\(t_i\) and \(c_i\longrightarrow c_{i+1}\) at the collision,
then
\[
 a(c_1,s_1,t_r)=\sum_{i=1}^r a_i,
\]
where \(a_i\) counts increments of \(c_i\) on its interval.
A collision chain may therefore move between different
identities with different current values. The saturation
argument below tracks births and individual lifetimes.

\emph{Finite nonmodal saturation.}
Here is a proof of the nonmodal conclusion of Lemma~6.2.7 using
the individual lifetimes just established. Suppose a path continues
forever using only Boolean, binder, and raw-variable rules, without
rejection or the modal or accepting action. Every rule except a
raw-variable rule strictly decreases the sum of the raw formula
sizes in the label. Therefore some variable is unfolded infinitely
often. Choose such an \(X\) whose binder is outermost among
those unfolded infinitely often: no proper enclosing variable is
unfolded infinitely often.

Count occurrences of the raw binding formula \(\sigma X\)
inside all formulas in the current label. Boolean decomposition
does not increase this count, and opening \(\sigma X\)
strictly decreases it. A raw-variable rule can increase it only
for a variable with a proper enclosing binder: the body of
\(\sigma X\), and the body of any binder inside it, contain
no occurrence of that enclosing raw binding formula. Bodies of
unrelated binders cannot contain it either. By the choice of
\(X\), there are only finitely many of these replenishing
enclosing-variable steps. Each inserts a finite body. Hence the
raw \(\sigma X\)-binder is opened only finitely often, and
only finitely many \(X\)-counter identities are born. This
count uses no fairness assumption on the order of rule choices.

If \(X\) is a \(\mu\)-variable, each of those finitely many
identities has finitely many increments by the lifetime argument.
Collisions select existing identities, and aconjunctivity never
duplicates a \(\mu\)-identity. Thus \(X\) cannot be unfolded
infinitely often.

If \(X\) is a \(\nu\)-variable, every \(\mu\)-counter
preceding its counter on a local list belongs to a proper
enclosing variable. The same binder-occurrence count shows that
each such enclosing variable has finitely many counter births:
every variable enclosing it also encloses \(X\), and so has
only finitely many raw-variable unfoldings. There are finitely
many enclosing variables, hence finitely many possible preceding
\(\mu\)-identities. Each has finitely many increments and
resets by the lifetime argument. Consequently the dependent-reset
rules cause only finitely many resets of \(X\)-counters, and
there are no modal resets on this path. A collision does not
decrease the retained \(\nu\)-identity's value: it either
keeps the selected list's value or takes a maximum in a
\(\nu\)-only list. There are finitely many \(X\)-identities,
and the zero guard permits at most one increment of each between
successive resets. Sharing an identity at conjunction creates
neither a new identity nor a new zero value. Again there are only
finitely many \(X\)-unfoldings, a contradiction.

Every nonmodal segment therefore reaches rejection, a modal action,
or acceptance after finitely many steps. This proves saturation
between modal actions; it does not make an arbitrary unbounded
sequence of modal actions a finite computation, and it does not
replace the separate simulation and model-construction arguments
\cite[lemmas~6.2.6--6.2.7, pp.~345--346]{Kozen1983}.

\paragraph{Deciding nonrejection on the configuration graph.}
Infinite modal continuations are allowed. The relevant
outcome is whether disjunction choices can avoid
rejection against every modal-successor choice,
not whether every branch eventually halts.
The finite configuration graph makes that outcome
a decidable greatest-fixed-point condition
\cite[\S6.2 and thm.~6.3.1]{Kozen1983}.

Fix an order for selecting an applicable administrative
rule, using the origin-based collision convention
already specified. Its structural successors and counter
updates remain those above. A configuration records
the raw label, counter variables and values, local
lists and sharing, and global order. Counter names
themselves carry no information beyond these records.
Identify configurations differing only by a consistent
renaming of live counter identities. Every rule
respects this identification, including creation of
a fresh counter: its new identity is distinct from
all currently live ones.
The \(O(m^3)\)-bit bound therefore gives at most
\[
 |\mathcal C|=2^{O(m^3)}
\]
configuration types, including a rejecting outcome
for counter overflow.

Partition this finite graph into existential
disjunction configurations \(\mathcal C_\exists\)
and universal configurations \(\mathcal C_\forall\);
put deterministic one-successor steps in the
second class. A modal configuration has all of
its required diamond successors. Give an accepting
terminal a self-loop and a rejecting terminal
a self-loop, solely for evaluating the graph.
The former is safe and the latter belongs to
\(\mathsf{Bad}\). Every configuration now has a
nonempty finite successor set
\(\operatorname{Succ}(c)\).

For \(W\subseteq\mathcal C\), define
\[
\begin{split}
 \mathcal H(W)=(\mathcal C\setminus\mathsf{Bad})\cap
 \bigl(&\{c\in\mathcal C_\exists:
                   \operatorname{Succ}(c)\cap W\ne\varnothing\}\\
       {}\cup&
       \{c\in\mathcal C_\forall:
                   \operatorname{Succ}(c)\subseteq W\}\bigr).
\end{split}
\]
Then \(\nu W.\mathcal H(W)\) is exactly the
nonrejecting region. At a configuration in this
fixed point, choose a successor in it at each
existential node. Every universal successor also
stays in it, so every play avoids rejection.
The choice depends only on the current configuration.

Conversely, iterate from \(W_0=\mathcal C\) by
\(W_{i+1}=\mathcal H(W_i)\). A configuration
removed at stage \(i+1\), unless already bad,
has all successors removed by stage \(i\) if
it is existential, or has one such successor
if it is universal. Thus universal choices
can force the removal rank to decrease
regardless of existential choices, until a
bad configuration is reached. This proves
the claimed characterization in both directions.
It also reproduces rejection propagation:
one rejecting successor suffices at a universal
node, while both disjunctive successors must
reject at an existential node.

The descending sequence has at most
\(|\mathcal C|\) strict changes. Configurations,
their at most \(O(m)\) successors, and all
counter arithmetic can be generated from their
finite descriptions. Even direct enumeration
and fixed-point iteration therefore take
\(2^{O(m^3)}\) time. This is a finite
simulation of nonrejection, not an extra rule
asserting that a repeated tableau label may
simply be accepted: counter values, identity
sharing, and priorities are part of the state.

Finite nonmodal saturation also bounds the
nonmodal portions of this graph. A cycle made
only of nonterminal, nonmodal steps could be
repeated, lifting each fresh-counter creation
to a fresh identity, and would give an infinite
nonmodal computation. The saturation proof
excludes this. Such a segment has at most
\(|\mathcal C|-1\) edges before its next
outcome. Its counter increments, including
any collision-chain changes of identity,
are therefore bounded by \(2^{O(m^3)}\).

For example, with \(p_0=\nu X.\langle a\rangle X\)
and its one flag \(x\), there is the recurring
sequence
\[
\begin{aligned}
 (\{\langle a\rangle X\},x=0)
   &\longrightarrow(\{X\},x=0)\\
   &\longrightarrow(\{\langle a\rangle X\},x=1)\\
   &\longrightarrow(\{X\},x=0).
\end{aligned}
\]
The modal step resets the flag and permits another
variable expansion. No rejection or terminal
acceptance occurs, but the configurations remain
in the nonrejecting region. A one-state
\(a\)-loop is the corresponding satisfying model.

For a minimal execution, take \(p_0=\mu X.X\).
The root has no active-variable counters. Exposing the
body produces label \(\{X\}\) and a new \(\mu\)-counter
initially zero. Each variable expansion returns to
the same label and increments that counter. With no
other rule or reset, it eventually exceeds the finite
bound and rejects. This agrees with
\(\llbracket\mu X.X\rrbracket=\varnothing\).
For \(p_0=\nu X.X\), body exposure instead creates
a zero-valued \(\nu\)-flag. One variable expansion
returns to \(\{X\}\) with flag \(1\). That expansion
is now inapplicable, and there is no modal demand
or clash, so the branch accepts. This agrees with
\(\llbracket\nu X.X\rrbracket=S\). The two executions
distinguish a persistently unfolded least-fixed-point
obligation from a greatest-fixed-point flag.

\paragraph{From nonrejection to a model.}
\label{sec:p09-kozen-nonrejection-model}
Assume now that the initial configuration belongs to the
nonrejecting region. We prove a finite directed-model
consequence of that hypothesis, using ordinal-labelled
tableau obligations. No consistency hypothesis is used in
this implication. The model extraction and its six
saturation conditions follow the construction on
pp.~349--353 \cite{Kozen1983}; the finite representation
and its backedge-depth convention will be specified below.

First extend the individual-identity event bound to
collision continuations. Let \(L\in\mathbb N\) be a
uniform bound on one \(\mu\)-identity's increments
along a nonrejecting path, supplied by the lifetime
argument above; it may be chosen with
\(L=2^{O(m^3)}\). For a tracked \(\mu\)-counter
\(c\), also track the first \(\mu\)-counter \(d\)
in its local list. The active-prefix property retains
\(d\) whenever it retains \(c\). At a collision the
discarded list is replaced by a list whose first
\(\mu\)-counter occurs strictly earlier in \(G\).
Between collisions that first counter never moves
right: births are appended, and deletion preserves
the survivors' order. Thus its position strictly
decreases at every collision replacement. Starting
at a position at most \(K\), a continuation passes
through at most \(K\) counter identities.

Count an increment on the transition performing it,
before resolving collisions, even when its provisional
list is then discarded. This can add at most one final
increment per replaced identity to the lifetime count.
Consequently every finite continuation satisfies
\[
 a(c,t,u)\leq B:=K(L+1)=2^{O(m^3)}.
\]
If \(K=0\), there is no \(\mu\)-counter and set
\(B=0\). This is a uniform bound over all descendant
endpoints and all choices of modal successors, not
merely a bound on the current numerical value.

Choose one nonrejecting successor at every disjunction,
retain every required modal successor, and retain the
one successor at each administrative step. Starting
from the initial node \(r_0\), these choices give a
rooted nonrejecting tableau tree \(T'\). Its leaves
are accepting configurations. Here the artificial
terminal self-loops used for the safety computation
are not tableau edges. Nonmodal saturation makes
every segment between modal actions finite; after
the disjunction selection it is also unbranched.

Let \(S\) consist of the modal-expansion nodes and
the leaves of \(T'\). For a tableau node \(t\), let
\(U(t)\) be the path containing \(t\) and its
preceding nodes back to, but not including, its most
recent strict ancestor in \(S\); if there is no such
ancestor, include \(r_0\). For \(s\in S\), this is
its incoming nonmodal segment, including its endpoint.
Each outgoing modal edge enters a unique such segment,
whose endpoint exists by nonmodal saturation. The
first member \(\rho\) of \(S\) on the initial
segment exists, is unique, and satisfies
\(r_0\in U(\rho)\); it need not equal \(r_0\).

Define a transition model \(M=(S,I)\).
Set \(I(0)=\varnothing\) and \(I(1)=S\).
For program constants \(a\) and ordinary
proposition constants \(P\), put
\[
\begin{aligned}
 (s,t)\in I(a)
   &\quad\Longleftrightarrow\quad
   \text{an actual \(a\)-labelled tableau edge from \(s\)
          enters \(U(t)\)},\\
 I(P)&=\{s\in S:P\in\Gamma_s\}.
\end{aligned}
\]
An accepting leaf has no outgoing model edges.
Boxes do not create edges by themselves: each edge
is justified by a particular diamond demand.

For a live \(\mu\)-counter \(c\) at a tableau
node \(t\), define its future-increment rank
\[
 \alpha(t,c)=\sup\{a(c,t,u):
       u\text{ is a descendant endpoint of a finite
       counter continuation in }T'\}.
\]
The zero-length continuation contributes zero. Follow
the corresponding local-list entry through inheritance
and restriction and, at a collision, through the
discarded-to-retained arrow. The event convention above
counts a raw-variable increment even when that step
also changes the retained identity. The bound gives
\(\alpha(t,c)\in\{0,\ldots,B\}\). A reset of a
counter's present value is not a future increment and
does not identify that value with \(\alpha(t,c)\).

For \(p\in\Gamma_t\), write
\(V_p=(X_1,\ldots,X_n)\) and write the binder
of \(X_i\) as \(\sigma_iX_i.u_i\).
For \(X_i\in A_{\mu p}\), let \(c_i\) be
the unique \(X_i\)-entry of \(C(t,p)\);
the index \(i\) refers to \(V_p\), not to
the position in its shorter active sublist.
Define
\[
 q_i(t)=
 \begin{cases}
   \alpha(t,c_i)X_i.u_i,
       &X_i\in A_{\mu p},\quad X(c_i)=X_i,\\
   \sigma_iX_i.u_i,&X_i\notin A_{\mu p},
 \end{cases}
 \qquad
 e''(t,p)=p[\bar q(t)/\bar X].
\]
The replacements use the book's ordered convention
\(p[q_n(t)/X_n]\cdots[q_1(t)/X_1]\): inner
parameters are replaced first, and the later outer
replacements close the formulas inserted earlier.
Thus \(e''(t,p)\) is closed. This substitutes a
finite ordinal stage for an active \(\mu\)-parameter;
it does not change the ordinary binder kind at a raw
binding formula being opened. Put
\[
 \Theta'_t=\{e''(t,p):p\in\Gamma_t\},
 \qquad
 \Theta_s=\bigcup_{t\in U(s)}\Theta'_t
       \quad(s\in S).
\]

The rank recurrences are obtained by partitioning
finite continuations at their first step. Along a
nonmodal successor a continued counter has the same
rank unless its own raw variable is expanded; in
that case its rank is one plus its successor rank,
following a collision arrow when necessary. A
diamond-list counter continues only along that
diamond's required edge. A box-list counter continues
along every matching diamond edge, and its rank is
the supremum of the successor ranks, with empty
supremum zero. Aconjunctivity prevents the same
\(\mu\)-identity from also belonging to an
unrelated diamond or box list. Finally, retaining an
active variable retains its entire active prefix.
Outer \(\mu\)-parameters therefore have equal
ranks on a nonmodal continuation that does not
expand them; a list collision transfers the
corresponding prefix entries together.

\paragraph{Saturation of the ordinal-closed labels.}
For formulas in \(L_{\mu+}\), write \(v\preceq u\)
when \(v\) is obtained from \(u\) solely by decreasing
ordinal indices at some ordinal-stage binders.
All other constructors, binder kinds, and variable
links are unchanged, up to capture-avoiding renaming.
In particular, this relation does not replace an
ordinary \(\mu\)-binder by a finite stage. It is
reflexive and transitive, is preserved by hygienic
substitution and positive formula formation, and
implies semantic inclusion by the approximant and
polarity arguments already proved.

For \(s\in S\), the following are the six
saturation conditions. In the fifth condition
\(\kappa\) is an ordinal or the symbol \(\mu\),
and every ordinal is declared smaller than that
symbol for this comparison:
\[
\begin{aligned}
 p\land q\in\Theta_s
   &\Longrightarrow p,q\in\Theta_s,\\
 p\lor q\in\Theta_s
   &\Longrightarrow p\in\Theta_s\text{ or }q\in\Theta_s,\\
 \langle a\rangle p\in\Theta_s
   &\Longrightarrow
      \exists t\ ((s,t)\in I(a)\land p\in\Theta_t),\\
 [a]p\in\Theta_s
   &\Longrightarrow
      \forall t\ ((s,t)\in I(a)\Rightarrow
         \exists p'\preceq p\ (p'\in\Theta_t)),\\
 \kappa X.u\in\Theta_s
   &\Longrightarrow
      \exists\beta\in\{0,\ldots,B\}\
        (\beta<\kappa\land u[\beta X.u/X]\in\Theta_s),\\
 \nu X.u\in\Theta_s
   &\Longrightarrow u[\nu X.u/X]\in\Theta_s.
\end{aligned}
\]
These are conditions (6.3.16)--(6.3.21), with the
finite range of the ranks made explicit
\cite[pp.~350--351]{Kozen1983}.

For the first two conditions, follow the raw
conjunction or disjunction through its finite
nonmodal segment, using collision continuations if
its list is absorbed by an existing copy. It must
eventually be decomposed, since otherwise the endpoint
would still have an applicable nonmodal rule. The
outside ranks are unchanged on that continuation.
The translated children are therefore exactly the
conjuncts, or the selected disjunct, of the formula
being followed. Their entries belong to the same
segment union \(\Theta_s\).

A diamond similarly persists to the modal endpoint.
Its required edge exists, and its counter continuation
has the same ranks along that edge. Hence its exact
translated body belongs to the next segment union.
For a box, every matching edge receives its raw
body. Each rank on one such edge is at most its
supremum before branching, so its translated body
is obtained only by lowering ordinal indices.
This proves the diamond and box conditions, including
the vacuous box case when there is no matching edge.
Literals persist to the segment endpoint. Consequently
\(P\in\Theta_s\) iff \(P\in\Gamma_s\), no
\(\Theta_s\) contains both \(P\) and \(\neg P\),
and \(0\notin\Theta_s\), by nonrejection and
the normalization of Boolean constants.

For the fifth condition, first consider a formula
whose top-level origin is a raw ordinary
\(\mu\)-binder. At its opening the outside
ordered substitutions are unchanged. In the
nonvacuous case, take the new counter, or its
corresponding retained entry after a collision,
with successor rank \(\beta\leq B\). The
translated body is then exactly
\(u[\beta X.u/X]\), with \(\beta<\mu\).
For a finite-stage formula \(\alpha X.u\), its
top-level origin is a raw occurrence of the
\(\mu\)-variable \(X\) with rank \(\alpha\).
It cannot remain unexpanded at the saturated
endpoint. At its eventual expansion, with continued
counter \(c'\), the rank recurrence gives
\[
 \alpha=\alpha(t,c)=1+\alpha(t^+,c'),
 \qquad \beta=\alpha(t^+,c')<\alpha.
\]
The outside active parameters are unchanged, so the
translated raw body is \(u[\beta X.u/X]\) in the
same segment. Counting the increment before a
collision is what makes this equality valid also
when \(c'\) has a different identity. In particular
a top-level zero-stage formula cannot occur in
\(\Theta_s\).

For the greatest-fixed-point condition, an ordinary
raw \(\nu\)-binder is treated by its opening, and
a raw \(\nu\)-variable with flag zero by its
permitted expansion. If its flag is one at the
saturated endpoint, trace the provenance of that
one backward inside the same segment. At inheritance
follow the old flag; at a \(\nu\)-only maximum
collision follow an input attaining the value one,
even if that input's identity was discarded; at
sharing follow the input from which the common flag
arose. A reset or modal step cannot be crossed while
tracing a one. The segment is finite and births have
value zero, so this trace reaches a genuine
\(\nu\)-variable expansion within the segment.

The earlier expansion has exactly the present
outside \(\mu\)-parameters. In a retained
\(\mu\)-containing list, an increment or reset
of any preceding \(\mu\)-counter would reset
the flag, so none occurs on the provenance interval;
the rank recurrence preserves these parameters.
A \(\nu\)-only collision has no preceding
\(\mu\)-parameters to compare. Nor can sharing
a \(\nu\)-flag across conjunction children
share a preceding active \(\mu\)-counter:
activity transitivity would make that variable
active in both children, contrary to aconjunctivity.
The shared case therefore also has no preceding
\(\mu\)-parameters. Thus the earlier genuine
expansion placed exactly \(u[\nu X.u/X]\) in
this \(\Theta_s\), proving the sixth condition.

These binder arguments include vacuous binders.
If \(X\notin\operatorname{FV}(u)\), opening the
binder creates no \(X\)-counter. For \(\mu\)
choose \(\beta=0\); then
\(u[0X.u/X]=u\). For \(\nu\), likewise
\(u[\nu X.u/X]=u\). The translated raw body
already supplies the required formula, without
assuming a nonexistent counter.

\paragraph{The general truth lemma.}
For a closed formula \(q\), define
\[
 q^\Theta=\{s\in S:
           \exists q'\preceq q\ (q'\in\Theta_s)\}.
\]
In particular,
\(q'\preceq q\Longrightarrow(q')^\Theta\subseteq q^\Theta\).
Let \(p\) be positive-normal in its free variables
\(\bar Y=(Y_1,\ldots,Y_n)\), and let
\(\bar q=(q_1,\ldots,q_n)\) be closed positive-normal
\(L_{\mu+}\)-formulas. Write \(p^M(\bar A)\)
for its denotation in \(M\) with \(Y_i\) assigned
\(A_i\), and put
\(\bar q^\Theta=(q_1^\Theta,\ldots,q_n^\Theta)\).
We prove
\[
 \bigl(p[\bar q/\bar Y]\bigr)^\Theta
       \subseteq p^M(\bar q^\Theta).
\]
This is the parameterized truth assertion (6.3.22)
\cite[pp.~351--353]{Kozen1983}. All substitutions
are capture-avoiding; since the parameter formulas
are closed, their sequential order has no semantic
effect. The proof is structural induction, treating
all ordinal stages of a binder body by a further
transfinite induction before its ordinary
least-fixed-point case.

For an ordinary proposition constant,
\(P^\Theta=I(P)\). The absence of literal clashes
gives \((\neg P)^\Theta\subseteq S\setminus I(P)\).
Also \(0^\Theta=\varnothing\) and
\(1^\Theta\subseteq S\). For a free variable
\(Y_i\), substitution gives exactly
\(q_i^\Theta\), its assigned semantic value.

For a Boolean constructor, take the particular
index-decreased formula witnessing membership in
the left side. Lowering indices preserves that
outer constructor. The first saturation condition
gives both conjunct witnesses, while the second
gives a witness for one disjunct. Therefore, writing
\(u_{\bar q}=u[\bar q/\bar Y]\) and similarly
for \(v\),
\[
\begin{aligned}
 (u_{\bar q}\land v_{\bar q})^\Theta
     &\subseteq u_{\bar q}^\Theta\cap v_{\bar q}^\Theta
       \subseteq u^M(\bar q^\Theta)\cap v^M(\bar q^\Theta),\\
 (u_{\bar q}\lor v_{\bar q})^\Theta
     &\subseteq u_{\bar q}^\Theta\cup v_{\bar q}^\Theta
       \subseteq u^M(\bar q^\Theta)\cup v^M(\bar q^\Theta).
\end{aligned}
\]
The last inclusions are the body induction hypotheses,
and their right sides are the required Boolean denotations.

For a diamond, apply its saturation condition to
the witnessing diamond. Some successor has a body
witness below \(u_{\bar q}\), so it lies in
\(u_{\bar q}^\Theta\) and hence, by induction,
in \(u^M(\bar q^\Theta)\). This is precisely the
diamond semantics. For a box, every matching
successor receives a formula below the witnessing
body and therefore, by transitivity, below
\(u_{\bar q}\). Every such successor consequently
lies in \(u^M(\bar q^\Theta)\), which is the
box semantics. The assertion includes a state
with no matching successors.

For the binder cases, let its body be
\(u(X,\bar Y)\), renaming its bound variable
away from the parameters. Set
\[
 U=u[\bar q/\bar Y],\qquad
 r_\beta=\beta X.U,\qquad
 F(A)=u^M(A,\bar q^\Theta).
\]
The operator \(F\) is monotone. Assume first
that the binder is \(\alpha X.u\) for an
ordinal \(\alpha\). Applying the fifth
saturation condition to each individual witness gives
\[
 (\alpha X.U)^\Theta
       \subseteq\bigcup_{\beta<\alpha}
                  \bigl(U[r_\beta/X]\bigr)^\Theta.
\]
The union is necessary: the witness \(\beta\)
may depend on the state. More explicitly, a witness
can have outer index \(\gamma\leq\alpha\) and
a body \(U'\preceq U\). Its unfolding uses some
\(\beta<\gamma\), and
\[
 U'[\beta X.U'/X]\preceq U[\beta X.U/X].
\]
This gives the indicated membership without assuming
that the undecreased formula belongs to \(\Theta_s\).
Structural induction on the body, with the additional
closed parameter \(r_\beta\), gives
\[
 \bigl(U[r_\beta/X]\bigr)^\Theta
       \subseteq F(r_\beta^\Theta).
\]
Transfinite induction on \(\alpha\), monotonicity,
and the increasing approximant chain now give,
for every \(\beta<\alpha\),
\[
 F(r_\beta^\Theta)
       \subseteq F(F^\beta(\varnothing))
       =F^{\beta+1}(\varnothing)
       \subseteq F^\alpha(\varnothing).
\]
Taking the union proves the ordinal-stage case.
For \(\alpha=0\) the witness union is empty,
as required. This induction proves the stage
assertion for every ordinal for the fixed body.

For the ordinary binder \(\mu X.u\), each
witness instead supplies some finite
\(\beta\leq B\), and hence
\[
 (\mu X.U)^\Theta
       \subseteq\bigcup_{\beta=0}^{B}
                  \bigl(U[r_\beta/X]\bigr)^\Theta
       \subseteq\bigcup_{\beta=0}^{B}
                  F^{\beta+1}(\varnothing)
       \subseteq\mu F.
\]
Here the ordinal-stage case has already been proved,
so no circular appeal to the \(\mu\)-case occurs.
There is still no single unfolding witness asserted
uniformly for all states; \(B\) only bounds the
set of possible witnesses.

Finally put \(r=\nu X.U\) and \(A=r^\Theta\).
The sixth condition, applied to an index-decreased
witness and then using substitution compatibility
of \(\preceq\), gives
\[
 A\subseteq\bigl(U[r/X]\bigr)^\Theta
       \subseteq u^M(A,\bar q^\Theta)=F(A).
\]
The second inclusion is structural induction on
the body with the additional closed argument
\(r\). Thus \(A\) is post-fixed for \(F\),
and the greatest-fixed-point characterization yields
\(A\subseteq\nu F\). This completes all binder
cases and the truth lemma.

The initial enclosing-variable list is empty, so
\(e''(r_0,p_0)=p_0\). Since
\(r_0\in U(\rho)\), we have
\(p_0\in\Theta_\rho\). Applying the truth
lemma with no parameters proves
\[
 \rho\in p_0^\Theta\subseteq\llbracket p_0\rrbracket_M,
 \qquad M,\rho\models p_0.
\]
This conclusion used only the selected nonrejecting
tree and its counter bounds.

\paragraph{A finite directed model and an explicit backedge depth.}
Choose the disjunctive successor by the positional
strategy on the finite nonrejecting configuration
region already constructed. Keep the same fixed
administrative scheduling and origin-based collision
convention. Configurations are normalized up to a
consistent renaming of counter identities; the rules
and all available successors respect that normalization.
Let \(\mathcal E\subseteq\mathcal C\) be the
reachable entry configurations of nonmodal segments:
the initial configuration \(c_0\) and the
configurations immediately following actual modal
edges. From \(c\in\mathcal E\), the chosen
nonmodal path has a unique modal or accepting
endpoint \(h(c)\), reached after finitely many steps.

Define the finite directed model \(D\) with
state set \(\mathcal E\) by placing an
\(a\)-edge \(c\to d\) for each actual
\(a\)-diamond-demand edge from \(h(c)\)
whose successor configuration is \(d\), and by
making each ordinary \(P\) true at \(c\) exactly when
\(P\in\Gamma_{h(c)}\). Interpret Boolean constants
at every state in the usual way. If \(h(c)\)
is accepting, \(c\) has no program edges; the
endpoint's artificial safety-loop
is not used. There are at most
\[
 |\mathcal E|\leq|\mathcal C|=2^{O(m^3)}
\]
states, each with at most \(m\) outgoing demand
edges before any duplicate edges are identified.

Unfold the same configuration strategy into the
selected tableau tree \(T'\) and apply the model
construction just proved. Map each endpoint
\(s\in S\) to the entry configuration of its
segment \(U(s)\). Its endpoint has normalized
configuration \(h(c)\), so this map preserves
proposition labels. Every tree-model edge projects
to an edge of \(D\). Conversely, every edge of
\(D\) out of an entry \(c\) comes from an
actual demand at \(h(c)\); every tree occurrence
with that entry has the corresponding demand edge
and a successor segment with the prescribed entry.
Thus the projection has both edge-matching
properties and is a bisimulation. The bisimulation
invariance proved in
Section~\ref{sec:p09-mu-bisimulation-invariance}
therefore transfers \(M,\rho\models p_0\) to
\(D,c_0\models p_0\). This is a finite directed
model obtained from full configurations, not a
filtration of raw labels.

For a precise rooted-tree representation of \(D\),
unfold paths from \(c_0\) only while their
projected states remain distinct. At a node
projecting to \(c\), handle each graph edge
\(c\xrightarrow{a}d\) as follows. If \(d\)
already occurs on that root-to-node path, put an
\(a\)-edge to its unique earlier occurrence,
allowing the current node itself when \(d=c\).
Otherwise create the corresponding child projecting
to \(d\). Copy proposition labels from \(D\).
Every path of underlying tree edges has distinct
projected states, so its depth, counting tree
edges and not backedges, is at most
\[
 |\mathcal E|-1\leq|\mathcal C|-1=2^{O(m^3)}.
\]
The branching is finite and this depth is finite,
so the representation is finite. Its projection
to \(D\) again preserves labels and has both
edge-matching properties: every representation
edge projects to a graph edge, and every graph
edge is represented at every node projecting to
its source. Bisimulation invariance therefore
preserves \(p_0\) at its root.

This representation is an ordinary rooted tree
with backedges to ancestors or to the current
node itself. It permits nested or crossing
backedges. The asserted depth counts only its
underlying tree edges; no additional restriction
on cycle nesting, crossing, or backedge targets
is part of this proved finite directed-model
consequence.

\paragraph{A context-counted completeness construction.}
\label{sec:p09-kozen-context-completeness}
For the same closed aconjunctive input class, take
\[
 M=2^m(m^2+1)^{m^2}.
\]
We prove that \(\mathsf A_M\) does not reject
any consistent input. This changes only the
finite overflow cap: the raw labels, all
eight counter assignments, collision ordering,
resets, scheduling, and completed-configuration
rejection convention are unchanged.
In particular, the annotations used below
are proof objects, not extra data stored
by the procedure. The argument uses the
annotated closing substitution and strengthened
induction of \cite[pp.~347--349]{Kozen1983},
with the context count established explicitly here.

\paragraph{Closed annotations and local consistency.}
Attach a closed formula \(r(s,c)\) to each
live \(\mu\)-identity \(c\). An annotation is
a formula value, not a reference to the current
annotation of another identity. Its initial value
is \(1\); inheritance preserves it; every
numerical reset sets it to \(1\); deletion
discards it. The only other change will
be specified at a \(\mu\)-variable expansion.
For \(V_p=(X_1,\ldots,X_\ell)\), write
the binder of \(X_i\) as \(\sigma_iX_i.b_i\).
Define
\[
 q_i(s,p)=
 \begin{cases}
  \mu X_i.(r(s,c_i)\land b_i),
       &X_i\in A_{\mu p},\\
  \sigma_iX_i.b_i,&X_i\notin A_{\mu p},
 \end{cases}
 \qquad
 e'(s,p)=p[\bar q(s,p)/\bar X],
\]
where \(c_i\) is the unique \(X_i\)-entry
in the local list \(C(s,p)\).
Use the previously defined inner-to-outer
ordered substitution. All outside parameters
are thereby closed. Let
\[
 E_s=\bigwedge_{p\in\Gamma_s}e'(s,p);
\]
an empty conjunction is \(1\).
Use the fixed enumeration of raw formulas when
forming these conjunctions and the contexts below.
Since \(r(s,c)\leq1\), fixed-point
monotonicity and positivity imply
\(e'(s,p)\leq e(p)\).
Replacing an annotation by \(1\) weakens
every affected translated formula.
The \(\nu\)-flags do not enter this translation.
All inequalities in this proof are
derivable formula inequalities in the
calculus, not consequences of an assumed
semantic completeness theorem.

Suppose \(E_s\) is consistent and the
raw variable \(X\), with binder
\(\mu X.b\) and counter \(c\), is expanded.
Put \(R=r(s,c)\).
For each \emph{inherited old} occurrence
\(p\in\Gamma_s\setminus\{X\}\), retain
its own local list, keep annotations
of identities higher than \(c\),
and replace lower annotations by \(1\).
Let its resulting translation be
\(\widehat e_s(p)\), and set
\[
 D_s=\bigwedge_{p\in\Gamma_s\setminus\{X\}}
                                      \widehat e_s(p).
\]
The identity \(c\) belongs only to
the old \(X\)-list, by \(\mu\)-disjointness.
Thus \(D_s\) is independent of its
new annotation. If \(b\) is already
in the old label, its factor in
\(D_s\) uses its inherited old list,
not the incoming body's list.

Close the outside parameters of
the \(X\)-occurrence with their
inherited annotations and call its
body \(b'(X)\). Both \(R,D_s\)
are closed and
\[
 E_s\leq D_s\land\mu X.(R\land b'(X)).
\]
The right side is therefore consistent.
Apply the strengthened-induction
consistency consequence proved above
with \(q=D_s\) and
\(p(X)=R\land b'(X)\).
After dropping the outer conjunct \(R\),
it gives consistency of
\[
 L_s=
 D_s\land
 b'\bigl(\mu X.(R\land\neg D_s\land b'(X))\bigr).
\]
Set the incoming annotation of \(c\)
to \(R\land\neg D_s\).
Before resolving collisions, retain
both the incoming translated body
and every inherited old translated
formula. Their conjunction is \(L_s\),
including both decorated occurrences
when their raw formulas coincide.
Collision disposal removes one of
these conjuncts, whichever local
list its priority rule retains.
Every subsequent deletion reset
replaces an annotation by \(1\).
Both operations weaken the temporary
conjunction, so
\[
 L_s\leq E_{s^+}.
\]
The completed successor is consistent.
If \(c\) is itself reset or discarded
during collision disposal, its new
constraint is reset or discarded
at that same operation.

For the remaining rules, first keep
the full temporary decorated family
and then apply the same weakening
argument for collisions and resets.
A conjunction is the conjunction
of its two translated children;
aconjunctivity prevents copying
a \(\mu\)-identity into both.
For a disjunction, the disjunction
of the two temporary successor
conjunctions equals the translated
parent conjunction. They cannot
both be inconsistent.
A binder opening with new
annotation \(1\), or a permitted
\(\nu\)-variable step, uses the
corresponding fixed-point equation.
For a diamond demand, normal modal
reasoning gives
\[
 \langle a\rangle u\land
       \bigwedge_i[a]v_i
 \leq
 \langle a\rangle\bigl(u\land\bigwedge_i v_i\bigr),
\]
where \(u,v_i\) are the translated
demand and box bodies.
If the temporary successor
conjunction were derivably \(0\),
then the parent would be derivably
below \(\langle a\rangle0=0\).
Thus every required modal successor
is consistent. A raw \(0\) or
opposite proposition literals
already makes \(E_s\) inconsistent.

Starting from \(E_{r_0}=p_0\), these
arguments construct a selected
annotated tableau tree with consistent
labels, retaining one consistent
disjunctive successor and all required
successors of other rules.
At this stage the tree is generated
without an overflow cutoff. We next
prove that none of its completed
configurations exceeds \(M\).

\paragraph{Context signatures across modal steps.}
Fix a surviving \(\mu\)-identity \(c\)
from its creation or last reset
until a later point before its
next reset. Its variable \(X(c)=X\)
is fixed. Let \(H_c\) be the identities
higher than \(c\) at the beginning
of this interval and let \(h=|H_c|\).
These identities, their order,
and their closed annotations remain
fixed throughout the interval:
births are appended below \(c\);
an earlier deletion or increment
would reset \(c\); and an earlier
reset is caused by a still earlier
event that would also reset \(c\).
Their holder formulas may change.
The argument does not require
unchanged local lists between nodes.

Use the raw-key family \(\mathcal Q\)
defined above and put
\(k_p=|A_{\mu p}|\).
At a predecessor \(s\) of an
increment of \(c\), its context
signature records, for every
\(p\in\mathcal Q\setminus\{X\}\),
whether \(p\) is absent or present.
For a present \(p\), it records
at each of its \(k_p\) active
\(\mu\)-positions the occupying
identity in \(H_c\), if there is
one, and otherwise a marker \(*\).
The marker means annotation \(1\),
as prescribed in \(D_s\).
No \(\nu\)-flag, lower identity,
or counter value is recorded.

Equal signatures give equal
formulas \(D_s\), up to a fixed
hygienic renaming of bound variables.
Indeed, the raw formula fixes the
ordered substitution positions.
A recorded higher identity selects
the same fixed closed annotation;
\(*\) selects the ordinary binding
formula, since \(1\land b_i=b_i\).
Inactive and \(\nu\)-positions have
their fixed original replacements.
These choices determine the complete
ordered substitution, including
its nested binding formulas, and
the presence bits determine its
conjuncts. Closed annotations are
unaffected by the later outside
substitutions. This comparison
works across arbitrary intervening
modal steps.

Suppose two increment predecessors
\(u<v\) for this unreset surviving
counter have equal signatures.
At \(u\) its annotation acquires
\(\neg D_u\), which is still a
conjunct of \(r(v,c)\).
At \(v\) we have
\[
 E_v\leq D_v=D_u.
\]
The translated raw \(X\)-conjunct
at \(v\) is a least fixed point
whose body is conjoined with
\(r(v,c)\). Unfolding therefore gives
\[
 E_v\leq r(v,c)\leq\neg D_u.
\]
Consequently \(E_v\leq D_u\land\neg D_u=0\),
contradicting consistency.
Thus the context signatures at
these increment predecessors
are all distinct.
No implication has been transported
from one model state to another;
the two endpoint contexts were
compared as closed formulas.

\paragraph{Counting excludes overflow.}
There are at most \(K\leq m^2\)
live \(\mu\)-identities and
\(\sum_{p\in\mathcal Q}k_p\leq m^2\).
For a given raw \(p\ne X\),
there is one absence record and
at most \((h+1)^{k_p}\) presence
records. Hence the number of
possible signatures in this
reset-free interval is at most
\[
\begin{aligned}
 \prod_{p\in\mathcal Q\setminus\{X\}}
               \bigl(1+(h+1)^{k_p}\bigr)
 &\leq 2^{|\mathcal Q|-1}
                    (h+1)^{\sum_{p\ne X}k_p}\\
 &\leq 2^{m-1}(m^2+1)^{m^2}<M.
\end{aligned}
\]
Variable compatibility and identity
disjointness can only reduce this
overcount. The fixed annotations
may be large formulas, but they
are fixed parameters, not extra
signature choices.

A completed value \(M+1\) would
require \(M+1\) increments since
the counter's creation or last
reset, all surviving without a
reset or disposal of that identity.
Two of their predecessors would
have the same signature, contradicting
the preceding argument.
Thus no selected node has a live
counter above \(M\), and its
consistency excludes literal rejection.
This selected tree has exactly
the closure required for membership
in the nonrejecting region of
\(\mathsf A_M\): one surviving
choice at a disjunction and all
required choices at every other
step. Equivalently, induction on
a finite rejection rank shows
that none of its nodes can reject.
No global halting assertion is used.

As a useful sharper case, if all
annotations in \(H_c\) equal \(1\),
then
\[
 D_s=\bigwedge_{p\in\Gamma_s\setminus\{X\}}e(p)
\]
depends only on the raw label.
There are at most \(2^{m-1}\)
such contexts. In particular,
this sharper count applies to
the highest-priority counter,
whose higher family is empty.
The general bound above also
accounts for nontrivial higher
annotations and their changing holders.

\paragraph{Completeness and the resulting finite bounds.}
The model direction already proved
is valid for every finite cap \(T\).
To see the dependence explicitly,
the constant-rank prefix has at
most \((T+1)^K\) value vectors.
One identity's increment count
has a uniform bound
\[
 L_T=O\bigl(K((T+1)^K+1)\bigr),
\]
and collision continuations have
bound \(B_T=K(L_T+1)\).
These finite bounds give nonmodal
saturation, the bounded future
ranks, all six saturation conditions,
and the truth lemma exactly as
before. The finite-configuration
quotient and its bisimulation
argument likewise use only the
specified finite counter range.

For \(T=M\),
\[
 \log_2 M=m+m^2\log_2(m^2+1)
          =O(m^2\log(m+1)).
\]
The values of at most \(m^2\) live
\(\mu\)-counters, with \(\nu\)-flags,
raw labels, local lists,
sharing information and priority
order, therefore admit an
\(O(m^4\log(m+1))\)-bit encoding.
There are at most
\[
 2^{O(m^4\log(m+1))}
\]
normalized configurations.
Their at most \(O(m)\) successors
are effectively generated from
the same rule data.
Descending iteration of the finite
safe-configuration operator gives
a decision in
\(2^{O(m^4\log(m+1))}\) time.
The extracted directed model has
at most that many states; its
ancestor-or-self backedge
representation has underlying
tree depth at most one less
than the directed model's number of states.

We have consequently proved, for
every closed aconjunctive positive-normal
input \(p_0\), with the binder-renaming
and Boolean-constant conventions fixed
above, the equivalence
\[
\begin{array}{ll}
 \text{(i)}&p_0\text{ is consistent in the calculus};\\
 \text{(ii)}&\mathsf A_M\text{ does not reject }p_0;\\
 \text{(iii)}&p_0\text{ has a finite directed model}.
\end{array}
\]
The new context argument proves
(i)\(\Rightarrow\)(ii); the
rank construction proves
(ii)\(\Rightarrow\)(iii);
soundness proves
(iii)\(\Rightarrow\)(i).
The finite-model and exponential-time
bounds just displayed accompany
these implications.
The cap \(M\) and its estimates
belong to this context-counted
procedure, distinct from the
earlier \(N=2^m\) presentation.
The input class has not changed.
The following theorem covers the
full calculus.  The later sections
develop automata equivalence, parity
games, and alternation hierarchies.

\paragraph{Completeness of the full calculus.}
\label{sec:p09-kozen-full-completeness}
Walukiewicz proves completeness for the entire finitary
\(L_\mu\) language, with arbitrary finite nesting of its well-formed
fixed points \cite[\S6, Theorem~25 of the author version]{Walukiewicz2000}.
The fact and theorem numbers below refer to that version.

\begin{hbtheorem}{thm:p09-kozen-full-completeness}{cited=Walukiewicz2000}{Weak completeness of Kozen's equational calculus}
For every finitary \(L_\mu\)-formula \(p\), validity at every state of
every Kozen model, under every environment, is equivalent to
\(\vdash p=1\) in the calculus of
Section~\ref{sec:p09-kozen-calculus}.
Equivalently, \(p\) is unsatisfiable if and only if
\(\vdash p=0\).
The formula need not be aconjunctive or initially guarded.
Models have nonempty state sets and arbitrary relations for their
program constants; fixed points are taken in the full powerset algebra.
\end{hbtheorem}

The source states Theorem~25 as provability of the negation of every
unsatisfiable formula.  Boolean complementation turns that statement
into \(\vdash p=0\).  Applying it to \(\neg p\) gives the validity
form \(\vdash p=1\).  Soundness in
Section~\ref{sec:p09-kozen-calculus-soundness} supplies the converse
implications.

The source's system K1--K6 has the same Boolean equations, modal
principles, pre-fixed-point axiom, and Park rule displayed here
\cite[\S4]{Walukiewicz2000}.
Its convention \(p\leq q\) abbreviates \(p\land q=p\).
This is equivalent in Boolean equational logic to the present
\(p\lor q=q\): from either equation the other follows by absorption
and substitution.  For a formula candidate \(b\) in Park induction,
rename the bound variable \(X\) fresh for \(b\) and take \(Y\) fresh
for both the body and \(b\).  The capture-avoiding rule instance
\(Y:=b\) is
\[
 p[b/X]\leq b\quad\Longrightarrow\quad\mu X.p\leq b.
\]
It is applied to its proved premise as a whole rule instance.
These conventions identify the source theorem with the registered
equational calculus \texttt{sheet:hilbert-mu}.

The proof in the cited paper includes the required deductive
normalization step.  Fact~17 gives provable equivalence to its
positive, guarded, well-named presentation with cover notation.
In this positive form, negations apply only to proposition letters
and free variables.
For a positive guarded formula \(\varphi\), Theorem~26 constructs a
semantically equivalent disjunctive formula \(\widehat\varphi\)
together with
\[
 \vdash\varphi\leq\widehat\varphi.
\]
The direction of this proved implication matters: a derivation of
\(\widehat\varphi=0\) then yields \(\varphi=0\).
Theorem~18 supplies a derivation of the negation of every unsatisfiable
disjunctive formula.  Combining it with the displayed provable
implication and Fact~17 gives the stated completeness result
\cite[Fact~17 and Theorems~18, 25--26]{Walukiewicz2000}.
The theorem above cites that proof.  The semantic normal-form
construction later in this presentation establishes its own
same-model equivalence and finite-model consequences.

\paragraph{Equations, consistency, and finite premises.}
For arbitrary formulas \(p,q\), the theorem and Boolean algebra give
\[
\begin{aligned}
 \bigl(\forall\mathcal M,\eta\quad
       \llbracket p\rrbracket_\eta
          \subseteq\llbracket q\rrbracket_\eta\bigr)
    &\Longleftrightarrow\ \vdash p\leq q,\\
 \bigl(\forall\mathcal M,\eta\quad
       \llbracket p\rrbracket_\eta
          =\llbracket q\rrbracket_\eta\bigr)
    &\Longleftrightarrow\ \vdash p=q.
\end{aligned}
\]
For the first equivalence, inclusion is validity of
\(\neg p\lor q\), and the equation
\(\neg p\lor q=1\) is Boolean-equivalent to \(p\lor q=q\).
For the second, apply the first in both directions and use
antisymmetry.  In particular, every formula consistent in the
previously defined sense has a model and a satisfying state.

For a finite set \(\Gamma\) of formulas, let
\(\Gamma\models_{\mathrm{loc}}q\) mean that every state satisfying
all of \(\Gamma\), in any model and environment, satisfies \(q\).
Then
\[
 \Gamma\models_{\mathrm{loc}}q
 \quad\Longleftrightarrow\quad
 \vdash\left(\bigwedge\Gamma\right)\leq q,
 \qquad \bigwedge\varnothing=1.
\]
This follows by evaluating the finite conjunction at the same state
and applying the first equivalence.  The right-hand side is
theoremhood of one inequality in the equational calculus.

\paragraph{Why the completeness statement is weak.}
Let \(p\) be an ordinary proposition, \(a\) one program constant,
and put
\[
 r=\mu X.(p\lor\langle a\rangle X),\qquad
 \Gamma=\{r\}\cup\{[a]^n\neg p:n\in\mathbb N\}.
\]
Here \([a]^0u=u\) and \([a]^{n+1}u=[a]([a]^nu)\).
The reachability calculation above shows that \(r\) holds exactly
when some finite \(a\)-path reaches a \(p\)-state.
For a finite subset of \(\Gamma\), choose \(N\) greater than every
displayed exponent in that subset.  A chain
\(0\to1\to\cdots\to N\), with \(p\) true only at \(N\),
satisfies the subset at \(0\): its \(p\)-state is reachable, while
each required earlier depth satisfies \(\neg p\).
If the finite subset has no such exponent, any \(N\geq1\) works.

The whole set \(\Gamma\) has no satisfying state.  The formula \(r\)
requires a finite path of some length \(n\) to \(p\), whereas
\([a]^n\neg p\) forbids \(p\) at the endpoint of every path of that
length.  This proves failure of compactness.
A derivation in any sound finitary consequence calculus uses only
finitely many premises.  It therefore cannot derive contradiction
from this \(\Gamma\), since that finite subset has a model.
Consequently weak completeness for individual valid formulas and
the displayed finite-premise result do not extend to strong
completeness for arbitrary local premise sets.

\subsubsection{Emerson--Jutla's tree-automata and determinacy presentation}
\label{sec:p09-presentation-n-c-018-emerson-jutla}

Emerson--Jutla connect fixed-point formulas, automata on infinite binary
trees, and games with recurring-color objectives. Their automata distinguish
the choice of a transition, the choice of an input direction, and the choice
of a continuing obligation. We first define those layers, then construct
the formula automaton and derive its strategy property
\cite[\S\S2--4]{EmersonJutla1991}.

\paragraph{Nondeterministic and alternating binary-tree games.}
For a finite alphabet \(\Sigma\), an input is a full infinite binary tree
\[
 t:\{0,1\}^*\longrightarrow\Sigma.
\]
Definition~2.1 gives a nondeterministic tree automaton as
\[
 \mathcal N=(\Sigma,OR,AND,\psi,\delta,L,s_0,\Omega),
\]
where \(OR\) is a finite state set, \(AND\) is a set of transition-choice
nodes,
\[
 \psi\subseteq OR\times AND,\qquad
 \delta:AND\times\{0,1\}\longrightarrow OR,\qquad
 L:AND\longrightarrow\Sigma,
\]
\(s_0\in OR\), and \(\Omega\subseteq OR^\omega\) is the acceptance condition.
At a tree node, Player~I chooses an \(AND\)-node related by \(\psi\) to the
current \(OR\)-state and labelled by the current input symbol; Player~II chooses
direction \(0\) or \(1\).  A legal Player-I strategy is winning exactly when
every input branch produces an \(OR\)-state sequence in \(\Omega\).  Thus
acceptance is existential in the strategy and universal in the input branch.

For pairs \((GREEN_i,RED_i)\), the source's Streett condition requires, for
every \(i\), that infinitely many visits to \(GREEN_i\) imply infinitely many
visits to \(RED_i\).  Its \emph{Pairs} condition is the Boolean complement of
that Streett condition.  These names denote acceptance families on infinite
\(OR\)-state sequences; they are not names for the transition graph.

Definition~2.4 inserts a universal-choice layer:
\[
 \mathcal A=
 (\Sigma,OR,DAND,AAND,\psi,\delta,\chi,L,s_0,\Phi),
\]
with
\[
 \psi\subseteq OR\times DAND,\qquad
 \delta:DAND\times\{0,1\}\longrightarrow AAND,\qquad
 \chi\subseteq AAND\times OR,\qquad
 L:DAND\longrightarrow\Sigma.
\]
The \(DAND\)-nodes specify the two input directions, whereas the
\(AAND\)-nodes specify universal nondeterministic choices.  Player~II now
chooses both a direction and an \(OR\)-successor of the resulting
\(AAND\)-node.  Consequently, one input branch and one Player-I strategy can
generate many \(OR\)-state sequences.

The source calls one finite generated \(OR\)-state sequence a \emph{thread} and
an infinite sequence all of whose finite prefixes are so generated an
\emph{ithread}.  This distinction is necessary because an input branch no
longer determines one automaton path.  Its standing is a sequence generated by
the automaton game, not a relation on the input tree; a thread is finite,
whereas an ithread is infinite.  The distinction is used both in
Definition~2.6, which requires every ithread on every branch to satisfy
\(\Phi\), and in the construction of Theorem~2.2.

Definition~2.7 calls a Player-I strategy \emph{history-free} when its choice at
an input node depends only on that node and the current automaton state, not on
the preceding universal choices.  An alternating automaton is history-free
when every accepted tree has such an accepting strategy.
The finite memory construction below supplies the strategy information
needed for general Muller and Streett acceptance
\cite[\S2]{EmersonJutla1991}.

\paragraph{Total presentations and their parameters.}
\label{sec:p09-emerson-jutla-nondeterminization}
We give the conversion to
nondeterministic Pairs automata
with explicit states and
bounds. A finite history
follows the transition
relations and current input
labels. Strategies select
legal moves at their
nonterminal owned histories.
A player with no legal
next move loses that
finite play. Winning
requires every compatible
maximal play to have
the prescribed winning
finite outcome or accepted
infinite control sequence.

Use a total presentation:
every \(OR\)-state and
input letter has an
enabled \(DAND\) choice,
and every directional
\(AAND\)-target has an
\(OR\)-successor.
All plays in this
presentation are infinite,
so its acceptance is
exactly the all-ithreads
condition.
A partial presentation
can be completed with
accepting and rejecting
absorbing controls.
Missing Player-I choices
lead in both directions
to the rejecting control;
empty universal successor
sets lead to the
accepting control.
Both controls have
transitions at every
letter.

If the original data
have \(n_0\) controls,
\(m_0\) directional choices,
and \(k_0\) Streett pairs,
the completed data satisfy
\[
 n\leq n_0+2,\qquad
 m\leq m_0+(n_0+2)|\Sigma|,\qquad
 k=\max\{k_0,1\}.
\]
When the pair list
is empty, first add
the vacuous pair
\((\varnothing,\varnothing)\).
Place the rejecting
control in one green
set and no red set,
and the accepting
control in no green
set. Original infinite
plays keep their old
condition, while the
two added tails have
the required outcomes.
The bounds below use
these total-presentation
parameters \(n,m,k\),
with \(k\geq1\).
If \(\Sigma\) is empty,
there are no full
input trees; one
nondeterministic control
with no transition
entries suffices.

\paragraph{A record for Streett responses.}
\label{sec:p09-streett-response-record}
We first prove a
finite word-monitor lemma.
Let an underlying state
stream have \(h\geq0\)
Streett pairs
\((G_i,R_i)\), indexed
by \(K=\{1,\ldots,h\}\).
Keep a permutation \(\pi\)
of these pair names.
At current state \(z\),
let
\[
 \operatorname{Resp}(z)=\{i\in K:z\in R_i\},\qquad
 \operatorname{Act}(z)=\{i\in K:z\in G_i\cup R_i\}.
\]
Positions in \(\pi\)
are numbered \(1,\ldots,h\).
If no name is active,
emit the even default
\(2h+2\).
Otherwise let \(a\)
be the first active
position and \(i=\pi(a)\).
Emit \(2a\) if
\(i\in\operatorname{Resp}(z)\),
and \(2a+1\) otherwise.
Denote this output by
\(\operatorname{out}(\pi,z)\).
A response takes precedence
when that name is
both triggered and
responded to.
Compute the output from
the old permutation,
then update it by
\[
 \operatorname{upd}(\pi,z)=
 \bigl(\pi\upharpoonright(K\setminus\operatorname{Resp}(z))\bigr)
 \mathbin{\cdot}
 \bigl(\pi\upharpoonright\operatorname{Resp}(z)\bigr).
\]
Both restrictions keep
their previous relative
order. This defines
one deterministic update
even for simultaneous
responses.

To prove correctness,
let \(F\) be the
names responding only
finitely often.
After their last responses
and one further response
of every other name,
the names in \(F\)
form a permanently
fixed prefix.
If a Streett implication
fails, take the first
name in that prefix
triggered infinitely often,
at position \(b\).
Earlier names eventually
have no events.
The name at \(b\)
never responds again,
so \(2b+1\) recurs
infinitely and every
smaller priority occurs
only finitely.
The output rejects.

If every implication
holds, the entire
fixed prefix eventually
has no events.
When no names remain
outside it, the even
default occurs forever.
Otherwise let the prefix
length be \(f\).
The oldest remaining
name stays at position
\(f+1\) until its
next response, which
must occur.
That response emits
\(2(f+1)\).
This happens after
every proposed final
such event, so the
even priority recurs
infinitely; no smaller
one does. The argument
also covers simultaneous
responses, since the
old positions determine
the output.
For \(h=0\), the
empty record always
emits \(2\), matching
the vacuous condition.

For a finite nonempty state
set \(Z\), initial state
\(z_0\), and total deterministic
map \(\eta:Z\times\Gamma\to Z\),
use the full product
\[
 Z'=Z\times\operatorname{Perm}(K)\times\{0,\ldots,2h+2\},
 \qquad p(z,\pi,c)=c.
\]
Choose \(\pi_0\in\operatorname{Perm}(K)\).
Store the last emitted color
using the transition
\[
 (z,\pi,c)\xrightarrow{a}
 \bigl(\eta(z,a),\operatorname{upd}(\pi,z),
                         \operatorname{out}(\pi,z)\bigr).
\]
Initialize at
\((z_0,\pi_0,2h+2)\)
and assign priority \(c\).
The priority stream is
the initial color followed
by the successive outputs,
so its recurring parity
is exactly the one
just proved.
There are
\[
 |Z|\,h!\,(2h+3)
\]
states in this full
product.
Convert its parity
condition to Rabin by
using, for each even
\(d\in\{2,4,\ldots,2h+2\}\),
the pair
\[
 A_d=\{v:p(v)=d\},\qquad
 B_d=\{v:p(v)<d\}.
\]
Infinite visits to
\(A_d\) and finitely
many to \(B_d\)
hold exactly when the
least recurring priority
is that even \(d\).
Priority zero cannot
recur in this monitor.
Thus the same states
with \(h+1\) Rabin
pairs recognize its
Streett language.

\paragraph{Finite memory before collecting threads.}
Let \(Q=OR\), with
\(|Q|=n\), in the
total input automaton.
Its Streett condition
defines the Muller
family
\[
 \mathcal F=\{I\subseteq Q:
    \forall j\ (I\cap G_j=\varnothing
                         \text{ or }I\cap R_j\ne\varnothing)\}.
\]
Apply the full-state
permutation record of
Section~\ref{sec:p09-muller-parity-record}
to this original
\(n\)-element set.
Use augmented controls
\[
 \widehat Q=
 Q\times\operatorname{Perm}(Q)\times\{0,\ldots,2n+1\},
 \qquad
 a_0=(q_0,\pi_0,0).
\]
Their priorities are
their last coordinates.
There are
\[
 S=n\,n!\,(2n+2)
     =2^{O(n\log(n+1))}
\]
augmented controls.

For \(a=(q,\pi,c)\)
and an original successor
\(r\), let \(\vartheta(a,r)\)
be the augmented successor
obtained by reading \(r\)
in that record.
Explicitly, if \(r\)
has old position \(b\),
move it to the end
to get \(\pi'\), and
let \(I_b\) be the
old suffix set. Then
\[
 \vartheta(a,r)=
 \begin{cases}
  (r,\pi',2b),&I_b\in\mathcal F,\\
  (r,\pi',2b+1),&I_b\notin\mathcal F.
 \end{cases}
\]
At \((x,a)\), Player~I
chooses an original enabled
\(d\in DAND\) for \(q\)
with label \(t(x)\).
Player~II chooses direction
\(i\), then any
\[
 r\in\operatorname{Succ}(d,i)
       :=\{r:\chi(\delta(d,i),r)\},
\]
and play continues at
\((xi,\vartheta(a,r))\).
The original \(DAND\)
choice may be tagged
by \(a\) to give
the formal three-layer
tuple.

Every original history
has a unique augmentation,
and projection recovers
that history.
On every infinite branch,
the record's parity
is equivalent to the
original Streett condition.
Strategies transfer in
both directions by retaining
or updating this finite
record.
On a fixed accepted
tree, positional parity
determinacy then supplies
one winning choice
\[
 \lambda(x,a)\in DAND.
\]
The intervening universal
layers receive larger
default priorities.
Totality allows arbitrary
enabled extensions outside
the winning region.
This gives the finite
memory strategy before
any histories are merged.
The same construction
works for any specified
Muller family on \(Q\).

\paragraph{The deterministic monitor for all augmented threads.}
Let
\(\mathcal R=\mathcal P(\widehat Q\times\widehat Q)\).
A trace of a relation
word begins at \(a_0\)
and follows its successive
relation edges.
Use the bad-trace
B\"uchi construction of
Section~\ref{sec:p09-mu-trace-monitor}
with the augmented priorities.
Its search mode may
commit to an odd \(j\);
after commitment it
allows only priorities
at least \(j\) and
requires infinitely many
equal to \(j\).
It recognizes exactly
the words having a
bad infinite trace.
Its state count satisfies
\[
 B=S+\sum_{\substack{0\leq j\leq2n+1\\j\text{ odd}}}
       |\{a\in\widehat Q:p(a)\geq j\}|
       \leq S(n+2).
\]

Safra determinization gives
a total deterministic
Rabin word automaton
with
\[
 M=2^{O(B\log(B+1))}
\]
states and \(B\)
pairs for these bad
words. Complementing
the pairs' predicate
on its unique run
gives a Streett condition
for the good words:
all traces from \(a_0\)
have even least
recurring priority.
Apply the pair-response
record just proved,
and its Rabin conversion.
The resulting deterministic
good-word monitor
\[
 \mathcal G=(Z_G,\mathcal R,\zeta,z_G^0,\mathcal P_G)
\]
has at most
\[
 N_G=M\,B!\,(2B+3)
\]
states and \(B+1\)
Rabin pairs.
The permutation here
is over the \(B\)
pair names. Hence
\[
 \log N_G=O(B\log(B+1)).
\]
All of its state
updates and acceptance
tests are explicitly
given by the preceding
constructions.

\paragraph{The nondeterministic state and transition.}
Take
\[
 OR'=\{(D,z):\varnothing\ne D\subseteq\widehat Q,\ z\in Z_G\},
 \qquad s'_0=(\{a_0\},z_G^0).
\]
For a source \((D,z)\),
choose a map
\[
 f:D\longrightarrow DAND
\]
such that \(f(a)\)
is enabled at the
original control projected
from \(a\), and all
these choices have the
same input label \(\ell\).
For direction \(i\),
form the whole relation
\[
 \rho_i^{D,f}
   =\{(a,\vartheta(a,r)):
               a\in D,\ r\in\operatorname{Succ}(f(a),i)\}.
\]
Put
\[
 D_i=\operatorname{ran}\rho_i^{D,f},
 \qquad z_i=\zeta(z,\rho_i^{D,f}).
\]
Totality makes \(D_i\)
nonempty.
Create one nondeterministic
\(AND'\) choice
\(d'=((D,z),f)\), with
\[
 (D,z)\,\psi'\,d',\qquad
 L'(d')=\ell,\qquad
 \delta'(d',i)=(D_i,z_i).
\]
Every allowed augmented
successor is included in
\(\rho_i^{D,f}\).
The chosen map fixes
both child states.
Player~II's remaining
choice is the input
direction, exactly as
required for a nondeterministic
tree automaton.

Lift each Rabin pair
\((A,C)\) of \(\mathcal G\)
to
\[
 \bigl(\{(D,z):z\in A\},\{(D,z):z\in C\}\bigr).
\]
A run branch accepts
precisely when its
good-word monitor does.
Different input branches
may use different
witness pairs.

\paragraph{Every ithread is represented.}
Fix a run, with
chosen maps \(f_x\)
and states \((D_x,z_x)\)
at its input nodes.
On a branch
\(x_0=\epsilon,x_1,\ldots\),
use the relation for
each chosen direction.
Induction on the branch
prefix gives
\[
 D_{x_j}=
 \{a_j:\text{some finite trace }
   a_0,a_1,\ldots,a_j
   \text{ follows the first }j\text{ relations}\}.
\]
Initially the set is
\(\{a_0\}\); the
induction step is
exactly taking the
relation image.
The monitor coordinate
is the unique state
after that same relation
prefix.

Use \(f_x(a)\) as
the augmented alternating
choice at \((x,a)\)
when \(a\in D_x\).
A finite thread generated
by these choices is
exactly a finite relation
trace from \(a_0\).
The same holds for
infinite threads: each
relation edge is an
allowed universal successor,
and every such successor
is present.
The source's condition
on all finite prefixes
of an ithread therefore
agrees with the monitor's
trace condition.
The choice at a node
and augmented control
is one fixed value,
so merging endpoints
does not merge incompatible
strategy decisions.

If the original alternating
automaton accepts, take
the augmented positional
strategy \(\lambda\).
At each nondeterministic
node choose
\(f_x(a)=\lambda(x,a)\)
on its current \(D_x\).
On every input branch,
all resulting relation
traces are plays of
that winning strategy.
The word is good,
so its monitor and
the run branch accept.
This constructs an
accepting nondeterministic
run.

Conversely, an accepting
nondeterministic run supplies
the maps \(f_x\).
Use them in the
augmented alternating game,
and extend by enabled
choices elsewhere.
Every compatible play
follows a branch of
the fixed run and
is a trace of its
relation word.
That word is good,
so the play accepts.
Update the finite record
along the original
control history to
implement this strategy
in the original automaton.
The record equivalence
makes every original
ithread Streett accepting.
This proves the two
tree languages equal.

\paragraph{Counting the construction.}
The control count is
\[
 |OR'|\leq(2^S-1)\,M\,B!\,(2B+3).
\]
For each nonempty \(D\),
there are at most
\(m^{|D|}\leq m^S\)
maps into the original
\(DAND\) set.
Enabledness and the
common-label requirement
only reduce this count.
Since \(D\ne\varnothing\),
the common letter is
already determined by
any value of the map.
Therefore
\[
 |AND'|\leq |OR'|\,m^S.
\]
No further relation or
monitor successor is
chosen after that map.
Totality also implies
\(|\Sigma|\leq m\),
by fixing one original
control and selecting
an enabled choice for
each distinct letter.

The bounds on \(S,B,M\)
give
\[
\begin{aligned}
 |OR'|&\leq2^{\,2^{O(n\log(n+1))}},\\
 |AND'|&\leq2^{\,2^{O(n\log(n+1))}}\,
                         m^{\,2^{O(n\log(n+1))}},\\
 \#\Omega&\leq B+1
                  =2^{O(n\log(n+1))}.
\end{aligned}
\]
For \(k\geq1\), these
imply the stated bounds
with the additional
\(k^2\) and \(k\)
factors:
\[
\begin{aligned}
 |OR'|&\leq2^{\,k^2\,2^{O(n\log(n+1))}},\\
 |AND'|&\leq2^{\,k^2\,2^{O(n\log(n+1))}}\,
                         m^{\,2^{O(n\log(n+1))}},\\
 \#\Omega&\leq k\,2^{O(n\log(n+1))}.
\end{aligned}
\]
For \(n\geq2\),
\(n\log(n+1)=\Theta(n\log n)\),
giving the asymptotic
form in Theorem~2.2
\cite[\S2.1]{EmersonJutla1991}.
For a partial input,
substitute the completed
values bounded above.
In particular the
directional-choice factor
uses
\(m_0+(n_0+2)|\Sigma|\).
The construction and
counts cover all
represented ithreads and
both input directions.

\paragraph{Binary-tree semantics and formula closure.}
Definition~3.0 has proposition letters \(P,Q,\ldots\), proposition variables
\(X,Y,\ldots\), Boolean connectives, the child modalities \(EX,AX\), and the
fixed-point forms
\[
 \mu X.f(X),\qquad \nu X.f(X),
\]
where \(f\) is syntactically monotone in \(X\).  A sentence has no free
proposition variables.  Unlike the general transition-system semantics below,
Definition~3.1 deliberately takes a Kripke tree to be
\[
 t:\{0,1\}^*\longrightarrow Prop
\]
and puts \(x\models P\) exactly when \(P=t(x)\).  Thus each node has one
proposition-symbol label.  The paper notes the usual
\(2^{Prop}\)-labelling convention but uses this single-label convention to
match its tree-automaton alphabet.  The two presentations must not be silently
identified.

Put \(U=\{0,1\}^*\).  The modal clauses are
\[
\begin{aligned}
 x\models EXp
 &\Longleftrightarrow
   \exists i\in\{0,1\}\;(xi\models p),\\
 x\models AXp
 &\Longleftrightarrow
   \forall i\in\{0,1\}\;(xi\models p).
\end{aligned}
\]
The fixed-point clauses are
\[
\begin{aligned}
 x\models\mu X.f(X)
 &\Longleftrightarrow
 x\in\bigcap\left\{
   S\subseteq U:
   \{y\in U:y\models f(X)\text{ under }X:=S\}\subseteq S
   \right\},\\
 x\models\nu X.f(X)
 &\Longleftrightarrow
 x\in\bigcup\left\{
   S\subseteq U:
   S\subseteq
   \{y\in U:y\models f(X)\text{ under }X:=S\}
   \right\}.
\end{aligned}
\]
The first clause takes the least fixed point through
pre-fixed sets; the second takes the greatest through
post-fixed sets.

Every formula is next put in positive normal form.  The source writes
\(\operatorname{not}(q)\) for the positive normal form of \(\neg q\).
Definition~3.3 makes \(FL(p)\) the least set satisfying
\[
\begin{array}{rcl}
 p&\in&FL(p),\\
 q\in FL(p)&\Longrightarrow&\operatorname{not}(q)\in FL(p),\\
 q\lor r\in FL(p)&\Longrightarrow&q,r\in FL(p),\\
 q\land r\in FL(p)&\Longrightarrow&q,r\in FL(p),\\
 EXq\in FL(p)&\Longrightarrow&q\in FL(p),\\
 AXq\in FL(p)&\Longrightarrow&q\in FL(p),\\
 \lambda X.f(X)\in FL(p)&\Longrightarrow&
 f(\lambda X.f(X))\in FL(p),
 \quad\lambda\in\{\mu,\nu\}.
\end{array}
\]
This Fischer--Ladner closure is a finite statement family of subformulas,
positive complements, and one-step unfoldings used as the state material of
Theorem~3.1. It is not a semantic closure of the input tree.
Both modal clauses retain the operand of the formula
being decomposed.

\paragraph{Local unfolding and the formula automaton.}
\label{sec:p09-emerson-jutla-formula-automaton}
For local closure labels, write \(x\Vdash q\).
Their proposition literals agree with the
input tree, and their other conditions include
\[
\begin{aligned}
 x\Vdash p
 &\Longleftrightarrow x\not\Vdash\operatorname{not}(p),\\
 x\Vdash p\lor q
 &\Longleftrightarrow x\Vdash p\text{ or }x\Vdash q,\\
 x\Vdash EXp
 &\Longleftrightarrow
   \exists i\in\{0,1\}\;(xi\Vdash p),\\
 x\Vdash\mu X.f(X)
 &\Longleftrightarrow x\Vdash f(\mu X.f(X)).
\end{aligned}
\]
The conjunction and greatest-fixed-point
conditions follow by duality.
These are the local pre-model equations
of Definition~3.4
\cite[\S3]{EmersonJutla1991}.
An unfolding equation admits every
fixed point of its body operator.
The automaton must also enforce
the extremal choice of least
and greatest fixed points.

We prove the following construction.
For every closed formula \(f_0\)
there is a finite Streett
alternating tree automaton
accepting precisely its models
on the source's proposition-labelled
full binary trees. Every
accepted tree has an accepting
strategy depending only on
its current node and automaton
state.

Let \(P_0\) be the finite
set of proposition symbols
occurring in \(f_0\).
Choose a finite alphabet
\(\Sigma\) of proposition symbols
containing \(P_0\).
The source choice is
\(\Sigma=P_0\).
The empty-alphabet case
has no full labelled
input trees and the
equivalence is vacuous.
For that case take
\(OR=\{q\}\), initial
state \(q\), empty
\(DAND,AAND\) sets and
empty transition maps;
any acceptance condition
has the required empty
input language.
For the remaining construction,
assume \(\Sigma\ne\varnothing\).
At a node labelled \(P\),
the observed valuation on
the used propositions is
\[
 \operatorname{obs}(P)=\{Q\in P_0:Q=P\}.
\]
For the source alphabet this
is the singleton \(\{P\}\);
an unused ambient symbol
has the empty observation.
Translate the child operators
by
\[
\begin{aligned}
 EX\,A&\longmapsto
          \langle0\rangle A\lor\langle1\rangle A,\\
 AX\,A&\longmapsto
          [0]A\land[1]A.
\end{aligned}
\]
The two directional relations
are the unique left and
right successors. Structural
induction preserves denotations,
including the fixed points;
copied binders are renamed
freshly when the notation
is expanded.

Apply the formula-to-parity-automaton
construction of
Section~\ref{sec:p09-mu-automaton-construction}
to the resulting closed
formula. It already handles
unguarded fixed points and
internal syntax moves.
Let \(Q,q_0,p,\Delta\)
be its finite controls,
initial control, priorities,
and one-step transition map.
Read each original symbol
\(P\in\Sigma\) through
\(\operatorname{obs}(P)\) in
its proposition-letter table.
On a full binary input,
either modality of direction
\(i\) has the same single
successor, so write its
transition atom as \((i,r)\).
The resulting positive
Boolean transition is
\(\Delta(q,P)\).

\paragraph{Terms and the three transition layers.}
For a positive Boolean
expression \(B\), form its
finite family \(\mathcal U(B)\)
of conjunctive terms:
\[
\begin{aligned}
 \mathcal U(\top)&=\{\varnothing\},&
 \mathcal U(\bot)&=\varnothing,\\
 \mathcal U(\alpha)&=\{\{\alpha\}\}
       &&\text{for an atom }\alpha,\\
 \mathcal U(B\lor C)&=\mathcal U(B)\cup\mathcal U(C),\\
 \mathcal U(B\land C)&=
       \{U\cup V:U\in\mathcal U(B),\
                    V\in\mathcal U(C)\}.
\end{aligned}
\]
Thus \(B\) is equivalent
to the disjunction of
the conjunctions of its
terms. Each term retains
a derivation from the
Boolean expression; this
records its local existential
choices.

Add controls \(q_\top,q_\bot\).
Choose an even integer
\(e>\max p(Q)\), and
give them priorities \(e\)
and \(e+1\), respectively.
The extended control set
is \(OR=Q\cup\{q_\top,q_\bot\}\).
For \(q\in Q\), let
\[
 \mathcal U_{q,P}=
 \begin{cases}
  \mathcal U(\Delta(q,P)),
             &\mathcal U(\Delta(q,P))\ne\varnothing,\\
  \{\{(0,q_\bot),(1,q_\bot)\}\},
             &\mathcal U(\Delta(q,P))=\varnothing.
 \end{cases}
\]
For the two new controls,
use their unique terms
\[
 \{(0,q_\top),(1,q_\top)\},
 \qquad
 \{(0,q_\bot),(1,q_\bot)\}
\]
at every input letter.
There is now at least
one term for every
control and letter.

Define
\[
\begin{aligned}
 DAND&=\{(q,P,U):q\in OR,\ P\in\Sigma,\
                              U\in\mathcal U_{q,P}\},\\
 AAND&=DAND\times\{0,1\},\\
 L(q,P,U)&=P,\qquad
 \psi(q,(q,P,U)),\\
 \delta(d,i)&=(d,i).
\end{aligned}
\]
For \(d=(q,P,U)\), put
\[
 U_i=\{r:(i,r)\in U\},\qquad
 V_i=
 \begin{cases}
  U_i,&U_i\ne\varnothing,\\
  \{q_\top\},&U_i=\varnothing,
 \end{cases}
\]
and define
\[
 \chi((d,i),r)\quad\Longleftrightarrow\quad r\in V_i.
\]
The initial state is
\(s_0=q_0\).
These are exactly the
required types of the
\(OR,DAND,AAND\) presentation.
Every directional target
has a successor.
The accepting control
returns to itself in
both directions at every
letter, and so does
the rejecting control.

An empty true term
sends both directions
to \(q_\top\).
An unrequired direction
also leads to \(q_\top\).
A false transition sends
both directions to
\(q_\bot\).
Thus every state and
input letter has a legal
transition choice, including
those occurring outside
an accepting play.

\paragraph{Equivalence of the layered and Boolean games.}
Suppose a strategy wins
the original parity automaton
on the fixed input tree.
Retain a virtual history
of that game.
At \((u,q)\), its choices
in the finite Boolean
transition, retaining every
universal branch, determine
a term \(U\).
Choose the corresponding
\(DAND\) node.
If Player~II selects an
atom \((i,r)\in U\),
choose one finite compatible
Boolean path realizing that
atom and extend the virtual
history through it.
Repeated occurrences with
the same atom label pose
no problem: choose a
compatible occurrence and
keep that choice in
the virtual history.
An unrequired direction
enters \(q_\top\) and
wins from then on.

Every infinite play that
avoids the added accepting
control has the same
successive original controls
as one virtual winning
play. Its parity agrees.
A winning original strategy
cannot reach a false
transition, so the rejecting
fallback is not encountered.
This constructs a winning
layered strategy.

Conversely, retain a
virtual winning layered
history while playing
the original Boolean
automaton. At a control,
the selected term makes
the Boolean expression
true when precisely its
atoms are true.
The term's derivation
therefore supplies a
local strategy whose
reached atoms all lie
in that term and whose
constant terminals are
true.
A false fallback cannot
be selected on a winning
layered history, since
it forces \(q_\bot\)
in either direction.
For each reached atom
\((i,r)\), let the
virtual Player~II choose
that direction and control.
The resulting history
is legal. An infinite
original play then has
the same control sequence
as its infinite layered
play; a true terminal
already wins.
This proves equivalence
in both directions.

\paragraph{Streett pairs and the strategy theorem.}
For every odd priority
\(j\) occurring in \(OR\),
set
\[
 GREEN_j=\{q:p(q)=j\},\qquad
 RED_j=\{q:p(q)<j,\ p(q)\text{ even}\}.
\]
Give the layered automaton
the Streett condition
\[
 \bigwedge_{j\text{ odd}}
 \bigl(GREEN_j\text{ infinitely often}
       \Longrightarrow RED_j\text{ infinitely often}\bigr).
\]
On any infinite control
sequence, let \(m\)
be its least recurring
priority. If \(m\)
is even, it witnesses
the red response for
every recurring odd \(j\),
since \(m<j\).
If \(m\) is odd,
its green set recurs
and no lower even
priority does, violating
that pair. Thus this
Streett condition is
equivalent to parity
on every control sequence.

An auxiliary priority height
has the same comparison.
Choose \(H\) larger
than every
\(\lfloor p(q)/2\rfloor\)
and set
\[
 h(q)=H-\lfloor p(q)/2\rfloor.
\]
For an odd priority
at height \(i\), its
red set consists of
even controls at heights
at least \(i\).
The summary controls
carry these priorities
across input transitions.
On the original variable
diagram, enclosing binders
have greater height than
their descendants. This
is the height comparison
in the source's green
and red conditions.
The pre-model trail and
ordinal-signature arguments
in Sections~\ref{sec:p09-mu-premodel}
and~\ref{sec:p09-mu-premodel-signatures}
explain why the condition
enforces the required
fixed-point recurrence.

Finally fix the input
tree and consider the
layered parity arena.
Player~I makes choices
only at positions
\((u,q)\) with \(q\in OR\).
The intervening directional
and universal-control choices
belong to Player~II.
Give those intermediate
positions a default priority
above all \(OR\) priorities.
Every infinite play visits
\(OR\) infinitely often,
so the parity winner
is unchanged.
Positional determinacy
therefore supplies, on
every accepted tree,
a winning choice
\[
 f(u,q)\in DAND.
\]
Extend it outside its
winning region by any
enabled choice with label
\(t(u)\); such a choice
exists by construction.
The resulting total
strategy satisfies the
source's legality requirement
at every state history.
Its allowed dependence
is exactly the input
node and current state.
The equivalent Streett
condition makes it an
accepting history-free
strategy as well.
Together with the two
translations, this proves
the formula-automaton and
strategy conclusions of
Theorem~3.1
\cite[\S3]{EmersonJutla1991}.

\paragraph{The reverse translation retains the tree directions.}
The forward construction and
the nondeterministic conversion
give tree automata for the
source's \(EX,AX\) formulas.
For general ordered binary-tree
automata, the reverse translation
uses the directional modal
language with separate actions
\(0\) and \(1\).

Here is the reduction for
a nondeterministic Rabin tree
automaton. Its finite control
set \(Q\) and Rabin pairs
define a Muller family
\[
 \mathcal F=
 \{I\subseteq Q:\text{some pair }(G,B)
               \text{ has }I\cap G\ne\varnothing,\ I\cap B=\varnothing\}.
\]
Apply the permutation record
of Section~\ref{sec:p09-muller-parity-record}
to each control sequence
in a run. An augmented
control stores the original
control, its permutation
record, and the last
emitted priority.
For an original transition
with child controls \(r_0,r_1\),
update the current record
separately with \(r_0\)
and \(r_1\), giving
the two augmented child
controls.
Each original run has a
unique such augmentation
from the chosen initial
record, and every augmented
run projects to an
original run.
The proved record property
makes every branch parity
accepting exactly when its
original infinity set
belongs to \(\mathcal F\).
This produces a finite
nondeterministic parity
tree automaton for the
same language.

Encode each finite input
symbol \(a\) by a
fresh proposition \(P_a\),
with observed valuation
\(\{P_a\}\).
Use the original transition
on these singleton observations
and the fixed transition
\(\bot\) on every other
powerset observation.
Express each of its
transition choices as
\(\langle0\rangle r_0
 \land\langle1\rangle r_1\),
and take the disjunction
over enabled choices at
the current letter.
On a full binary tree
these are exactly the
two required child moves.
The reverse construction
of Section~\ref{sec:p09-mu-automaton-to-formula}
now yields a closed
formula in the directional
language.
Its semantic dual defines
the complementary language;
the forward construction
starts directly with that
directional formula over the
full alphabet of the \(P_a\).
Rename \(P_a\) back to
the original symbol \(a\)
in the resulting tree automaton.
The subsequent nondeterministic
conversion therefore gives
the tree-automaton complementation
construction.

The role of directions
can be seen in two
\(\{P,Q\}\)-labelled trees.
Let \(t_L\) have its
only \(P\) at node
\(0\), and \(t_R\)
have its only \(P\)
at node \(1\); all
other nodes carry \(Q\).
Swapping the two root
subtrees is an isomorphism
for the unindexed successor
relation. Every closed \(EX,AX\)
fixed-point formula therefore
agrees at their roots.
An ordered tree automaton
can check that child
\(0\) carries \(P\),
so it accepts \(t_L\)
and rejects \(t_R\).
The directional formula
\(\langle0\rangle P\)
expresses that distinction.
General ordered-tree equivalence
thus uses the two
directional modalities.

\paragraph{Parity fixed points and determinacy.}
\label{sec:p09-emerson-jutla-parity-formulas}
Fix \(n\geq0\) and a
set-sized total bipartite arena
\[
 V=O\mathbin{\dot\cup}A,\qquad
 E\subseteq(O\times A)\cup(A\times O).
\]
Player~I chooses at \(O\),
and Player~II at \(A\).
Every vertex has a successor.
The source case is a
game tree rooted in \(O\).
Its color map is
\[
 c:V\longrightarrow\{0,\ldots,n\},\qquad
 E(o,a)\Longrightarrow c(a)=c(o).
\]
Thus a play
\[
 o_0,a_0,o_1,a_1,\ldots
\]
has the same recurring
colors as its round
sequence \(o_0,o_1,\ldots\).
Write \(even_n\) and
\(odd_n\) for evenness
and oddness of the
largest recurring color.
The finite color range
and infinite play make
these complementary conditions
\cite[def.~4.0]{EmersonJutla1991}.

For \(U\subseteq O\),
the two round predecessor
operations are
\[
\begin{aligned}
 P_I(U)&=\{o\in O:\exists a\,[E(o,a)\land
                   \forall o'\,(E(a,o')\Rightarrow o'\in U)]\},\\
 P_{II}(U)&=\{o\in O:\forall a\,[E(o,a)\Rightarrow
                   \exists o'\,(E(a,o')\land o'\in U)]\}.
\end{aligned}
\]
They are represented at
\(O\)-vertices by \(EX\,AX\)
and \(AX\,EX\), respectively.
The order records that
Player~II observes Player~I's
first move before responding.
Both operations are monotone,
and De Morgan duality gives
\[
 O\setminus P_I(O\setminus U)=P_{II}(U).
\]

\paragraph{Target parameters and the nested formulas.}
First let the target
\(T\) be a subset of
\(O\), tested at round
boundaries, including the
initial vertex. Define
\[
\begin{aligned}
 I_n(T)&=\{o:\text{I can force }even_n\lor F T\},\\
 II_n(T)&=\{o:\text{II can force }odd_n\lor F T\}.
\end{aligned}
\]
Here \(F T\) means
that the round sequence
eventually enters \(T\).
For a formula parameter
\(t\), fix its external
environment and take
\(T=O\cap\llbracket t\rrbracket\).
The variables \(x_i\)
below are fresh for
that parameter.
During the proof, \(t\)
can be replaced by a
fresh predicate for \(T\);
semantic substitution then
restores the formula.

For every \(i\geq0\), set
\[
 \lambda_i=
 \begin{cases}\nu,&i\text{ even},\\\mu,&i\text{ odd},\end{cases}
 \qquad
 \bar\lambda_i=
 \begin{cases}\mu,&i\text{ even},\\\nu,&i\text{ odd}.\end{cases}
\]
Let \(C_i\) be the
predicate for color \(i\).
The two characterizing
formulas are
\[
\begin{aligned}
 \Phi_n^I(t)
 &=\lambda_nx_n.\cdots\lambda_0x_0.
      \left(t\lor\bigvee_{i=0}^n(C_i\land EX\,AX\,x_i)\right),\\
 \Phi_n^{II}(t)
 &=\bar\lambda_nx_n.\cdots\bar\lambda_0x_0.
      \left(t\lor\bigvee_{i=0}^n(C_i\land AX\,EX\,x_i)\right).
\end{aligned}
\]
Their denotations at
\(O\)-vertices are
\(I_n(T)\) and \(II_n(T)\).
Every recursive variable
is reached after two
arena edges. Therefore
the body value at \(O\)
depends only on variable
values at \(O\).
Restriction to \(O\)
commutes with unions and
intersections, so transfinite
induction on approximants
allows all these fixed
points to be computed
in \(\mathcal P(O)\).

For \(n=0\), the prefixes
are explicitly \(\nu x_0\)
for Player~I and
\(\mu x_0\) for Player~II.
The source's boundary
convention \(I_{-1}(T)=T\)
is the target test
with no continuing
color case.

\paragraph{Proof of the two characterizations.}
In the evaluation game
of \(\Phi_n^I(t)\),
the existential player
represents Player~I.
At a target it can
finish at the true
target predicate.
Otherwise a winning
strategy chooses the
actual color disjunct:
choosing another color
would expose a false
literal to its opponent.
After the true color
guard, \(EX\) chooses
an \(A\)-successor and
\(AX\) allows the
opponent's \(O\)-successor.
These are precisely the
two choices in one
arena round.
At its end, \(x_i\)
returns to its binder
and through the remaining
finite prefix to the
body at the next
\(O\)-vertex.

The variable-return sequence
therefore records the colors
of the source vertices
\(o_0,o_1,\ldots\).
At \(x_i\), the
evaluation priority is
\[
 2(n-i)+(i\bmod2).
\]
It decreases as the
color index increases
and has that color's
parity. Equivalently,
choose even \(M\geq n\)
and use \(M-i\);
then
\[
 \min\operatorname{Inf}(M-c)
       =M-\max\operatorname{Inf}(c).
\]
All other recurring
evaluation priorities are
larger defaults.
An infinite continuing
evaluation play thus wins
exactly when the arena
play has \(even_n\).

A strategy forcing
\(even_n\lor FT\)
gives an evaluation
strategy by choosing
the target at its
first occurrence and
otherwise copying the
arena moves.
True guard challenges
win immediately.
Conversely, keep a
virtual winning evaluation
history during an arena
play, continue past
true color guards,
and copy its existential
move while inserting the
opponent's actual reply.
Before a target, it
cannot select a false
target or color.
If a target is reached
the objective is met;
otherwise every infinite
play expands to a
winning evaluation play.
This proves the first
characterization.

For \(\Phi_n^{II}(t)\),
the existential evaluation
player represents Player~II.
The \(AX\) choice is
now made by the universal
evaluation player, representing
Player~I's original move;
the following \(EX\)
choice is Player~II's
response. The variable
priorities are
\[
 2(n-i)+(1-(i\bmod2)),
\]
equivalent to \(M-i+1\).
They reverse the color
order and reverse its
parity. Thus an infinite
evaluation play wins
exactly for \(odd_n\).
The same history construction
handles targets and both
strategy directions.
The evaluation-game theorem
proves the two formulas
for every \(n\geq0\)
\cite[thms.~4.1--4.2]{EmersonJutla1991}.

\paragraph{Duality and pure parity determinacy.}
For a general target,
the complement of Player~I's
objective is odd parity
together with avoidance
of the target.
For a predicate \(u\) denoting
\(U\subseteq O\), the corresponding
formula is
\[
 \Psi_n^{II}(u)=
 \bar\lambda_nx_n.\cdots\bar\lambda_0x_0.
 \left(u\land
      \bigvee_{i=0}^n(C_i\land AX\,EX\,x_i)\right).
\]
Write \(\llbracket\cdot\rrbracket_O\)
for restriction of a denotation
to \(O\), and use a set
as a predicate parameter by
assigning it that denotation.
The formula enforces \(u\) at
every round and odd
parity on infinite
continuations. Complement
duality of fixed points,
the predecessor identity,
and the fact that
the color predicates
partition \(O\) give
\[
 O\setminus I_n(T)
     =\llbracket\Psi_n^{II}(O\setminus T)\rrbracket_O.
\]
In particular, at the
empty target,
\[
 I_n(\varnothing)=O\setminus II_n(\varnothing).
\]
This gives a winner
at every root of the
pure parity game.
An \(A\)-root can
be given a fresh
forced \(O\)-parent
of the same color;
that finite prefix
preserves the winner.
These are the determinacy
conclusions of Corollary~4.3.

For \(T=O\), both
reach-target objectives
hold immediately, illustrating
why the general parameter
uses the reachability and
avoidance duality above.

\paragraph{Targets at every vertex.}
If the intended target
is a subset of all
of \(V\), including
intermediate \(A\)-vertices,
keep the same prefixes
and use the bodies
\[
\begin{aligned}
 H_I(\mathbf x)
   &=t\lor\bigvee_{i=0}^n
         \bigl(C_i\land EX(t\lor AX\,x_i)\bigr),\\
 H_{II}(\mathbf x)
   &=t\lor\bigvee_{i=0}^n
         \bigl(C_i\land AX(t\lor EX\,x_i)\bigr).
\end{aligned}
\]
The added test occurs
at the intermediate vertex
before its next move.
The same game correspondence
now detects the first
target on the full play.
This supplies the full-vertex
target version while the
previous formulas compute
the round-target version.

\paragraph{Identical remaining game trees.}
\label{sec:p09-emerson-jutla-residual-strategies}
For the strategy conclusion
of Theorem~4.4, present
the game tree by a
prefix-closed set \(D\)
of finite move sequences.
Its node labels specify
the owner and color.
The remaining game at
node \(u\) has relative
domain
\[
 D_u=\{w:uw\in D\}
\]
and the inherited labels
at \(uw\).
Call two occurrences
equivalent when these
remaining labelled trees
are equal, including
their available move names.
Then an available move
\(a\) at either occurrence
has equivalent continuations:
\[
 u\sim v\quad\Longrightarrow\quad ua\sim va.
\]

Form the quotient arena
with states \([u]\)
and named edges
\[
 [u]\xrightarrow{a}[ua].
\]
The equality just stated
makes its edges well
defined. Ownership, colors,
and totality are preserved.
The quotient is a
set of equivalence classes
of the given tree's
nodes and may be
finite or infinite.

Every tree play projects
to a quotient play
with the same colors.
After a starting occurrence
is fixed, every legal
sequence of named quotient
moves lifts recursively
to a unique tree play.
A quotient strategy
therefore lifts to the
tree. To transfer a
tree strategy in the
other direction, retain
its virtual tree history,
copy its own moves,
and extend that history
with each actual opponent
move in the quotient.
The lifted play preserves
the winner. Winning
strategies thus exist
at exactly the corresponding
starting positions.

The transferred winning strategy
places the starting class in
its player's quotient winning
region. Apply the least-recurring
version of positional determinacy
to this quotient with priorities
\(M-c\), choosing even \(M\geq n\).
On that winning region, choose
the player's positional strategy.
For each vertex owned by that
player in the region, fix a
move name realizing its selected
successor. Lift those
fixed choices back to
the tree.
At any two occurrences
with identical remaining
games, the strategy now
chooses the same move
name. Every resulting
play projects to a
winning quotient play.
Arbitrary fixed legal
choices on the other
quotient states extend
the strategy to all
histories.
This proves the source's
history-free conclusion
for either player
\cite[thm.~4.4]{EmersonJutla1991}.
For a target objective,
include target membership
in the remaining-tree
labels and stop at
the prescribed target
outcome before taking
another move.

\paragraph{The two on-line-game applications.}
\label{sec:p09-emerson-jutla-fsigma}
On a total bipartite game
tree, keep the round
predecessors \(P_I,P_{II}\).
Predicates and fixed-point
variables below range over
its \(O\)-vertices.
For a countable family
of state predicates,
\(\exists i\) denotes
the union of their
denotations and \(\forall i\)
their intersection.
The resulting monotone
operators have their usual
powerset least and greatest
fixed points. This gives
the semantics of the
infinitary schemata in
Section~5 of
\cite{EmersonJutla1991}.

\paragraph{Constrained reachability and safety with an exit.}
For \(B,C\subseteq O\),
put
\[
 R_{II}(B,C)=
   \mu X.\bigl(C\cap(B\cup P_{II}(X))\bigr).
\]
This is exactly the
region where Player~II
can force a finite
visit to \(B\), with
every round boundary
up to and including
that visit in \(C\).
An initial state in
\(B\cap C\) already
meets the objective.

For the forward proof,
assign each member its
first entry stage
\(\alpha+1\) in the
least approximants.
Before reaching \(B\),
the approximant equation
places the current vertex
in \(C\cap P_{II}(R_\alpha)\).
After every first move
by Player~I, Player~II
can choose a reply in
\(R_\alpha\).
The next entry stage
is strictly smaller,
and the next boundary
still belongs to \(C\).
There cannot be infinitely
many rounds of strict
ordinal descent, so a
permitted target is
reached.

Conversely, outside this
fixed point \(R\), a
vertex outside \(C\)
already fails the constrained
objective. At a vertex
of \(C\setminus R\),
the target is absent
and \(P_{II}(R)\) fails.
Predecessor duality gives
Player~I a first move
all of whose replies
stay outside \(R\).
Repeating this choice
either leaves \(C\)
or stays outside the
target forever.
This proves the converse
by an explicit opposing
strategy.

For \(T,B\subseteq O\),
put
\[
 S_I(T,B)=
   \nu X.\bigl((T\cap P_I(X))\cup B\bigr).
\]
It is exactly the
region where Player~I
can stay in \(T\)
until reaching \(B\),
or stay in \(T\)
forever if that exit
never occurs.
At a vertex in
\(S_I(T,B)\setminus B\),
the fixed-point equation
supplies an I-move
whose replies remain
in \(S_I(T,B)\), and
the vertex lies in
\(T\). This gives
the required strategy.
Its complement is
\[
 O\setminus S_I(T,B)
   =R_{II}(O\setminus T,O\setminus B).
\]
The preceding lemma
therefore gives the
opposing strategy outside
the stated region:
leave \(T\) before
any permitted exit
into \(B\).

\paragraph{The complementary nested regions.}
Given any countable family
\((T_i)_{i\in\mathbb N}\)
of subsets of \(O\),
define
\[
\begin{aligned}
 \mathcal F(Y)&=
       \bigcup_{i\in\mathbb N}S_I(T_i,Y),\\
 \mathcal G(Z)&=
       \bigcap_{i\in\mathbb N}R_{II}(O\setminus T_i,Z),\\
 I^*&=\mu Y.\mathcal F(Y),\qquad
 J^*=\nu Z.\mathcal G(Z).
\end{aligned}
\]
Each inner fixed point
is monotone in its
external parameter.
The complement identity
for \(S_I\), countable
De Morgan duality, and
fixed-point duality give
\[
 O\setminus\mathcal F(O\setminus Z)=\mathcal G(Z),
 \qquad
 J^*=O\setminus I^*.
\]
Expanding the definitions,
these are the schemata
\[
\begin{aligned}
 I^*&=\mu y\,\exists i\,\nu x
          \bigl((T_i\land EX\,AX\,x)\lor y\bigr),\\
 J^*&=\nu y\,\forall i\,\mu x
          \bigl((\neg T_i\lor AX\,EX\,x)\land y\bigr).
\end{aligned}
\]

\paragraph{Player I's decreasing sequence of phases.}
Let \(A_\alpha\) be the
least approximants of
\(\mathcal F\).
At a vertex of \(I^*\)
with first entry stage
\(\alpha+1\), choose
an index \(i\) such
that the vertex belongs
to \(S_I(T_i,A_\alpha)\).
It is outside \(A_\alpha\),
by minimality of its
entry stage.
Keep this index and
parameter fixed during
the phase and follow
the safety-with-exit
strategy.

If the exit never
occurs, every subsequent
boundary remains in
\(T_i\).
If the phase reaches
\(o'\in A_\alpha\),
restart there using
its first entry stage,
which is strictly
smaller than \(\alpha+1\).
Infinitely many restarts
would give an infinite
descending ordinal sequence.
Consequently every compatible
play eventually stays
in one component:
\[
 \forall\pi\text{ compatible}\quad
 \exists i\ \exists N\ \forall k\geq N
                   \quad o_k\in T_i.
\]
The final component
and the start of
that tail may depend
on the play.

\paragraph{Closed winning sets as prefix predicates.}
Let \(\mathcal P\) be
the legal infinite plays
from the game root,
with the topology
generated by finite
play prefixes.
Suppose
\[
 A=\bigcup_{i\in\mathbb N}F_i,
 \qquad F_i\subseteq\mathcal P\text{ closed}.
\]
For each component let
\(\widehat T_i\) be
the finite legal prefixes
that extend to a
play in \(F_i\).
It is prefix closed,
and
\[
 \pi\in F_i
 \quad\Longleftrightarrow\quad
 \text{every finite prefix of }\pi
                         \text{ belongs to }\widehat T_i.
\]
If a play is outside
the closed set, its
open complement contains
a finite-prefix cylinder.
That prefix has no
extension in \(F_i\),
which proves the
nontrivial direction.

At a round vertex
\(o\), define
\[
 o\in T_i
 \quad\Longleftrightarrow\quad
 \operatorname{hist}(o)\in\widehat T_i,
\]
where \(\operatorname{hist}(o)\)
is the entire original
root-to-\(o\) play prefix.
The round prefixes are
cofinal among all finite
prefixes. A violation
at an intermediate vertex
persists to the next
round boundary.
Once a boundary leaves
\(T_i\), every later
boundary remains outside
it.
Conversely, a tail that
stays in \(T_i\)
puts every earlier
prefix in \(\widehat T_i\)
as well, by prefix
closure.
Such a play belongs
to \(F_i\).

Player~I's phase strategy
from \(I^*\) eventually
stays in some \(T_i\).
It therefore produces
a play in \(F_i\),
and hence in \(A\),
against every opponent.

\paragraph{Player II's sequence of exclusions.}
Since \(J^*=\mathcal G(J^*)\),
for every \(i\),
\[
 J^*\subseteq R_{II}(O\setminus T_i,J^*).
\]
At a current boundary
in \(J^*\), record
the components whose
prefix trees have already
been left.
If all have been
left, every continuation
is outside \(A\);
choose any legal
continuation.
Otherwise choose the
least index \(i\)
not yet excluded.
The current vertex
lies in \(T_i\).
The constrained-reachability
lemma gives a strategy
forcing a finite visit
outside \(T_i\), with
all boundaries through
that visit in \(J^*\).
This phase contains
at least one round,
since its target is
initially false.

At the end, every
earlier exclusion persists,
and \(i\) is excluded
as well. Repeat.
For each fixed \(k\),
after at most \(k+1\)
completed phases all
indices through \(k\)
have been excluded;
an index can be
skipped only when its
violation has already
occurred.
Every phase that starts
finishes after finitely
many rounds.
Thus the resulting
infinite play belongs
to none of the
closed components and
lies outside \(A\).

The regions \(I^*,J^*\)
partition \(O\), and
their strategies win
complementary objectives.
Two opposing winning
strategies at one vertex
would produce a play
both inside and outside
\(A\), which is
impossible.
Hence these are exactly
the two winning regions.
This proves the fixed-point
characterization and determinacy
for \(F_\sigma\) winning
sets in Theorem~5.2
\cite[\S5]{EmersonJutla1991}.

\paragraph{Bounded and unbounded natural colors.}
For Theorem~5.1, give
each \(O\)-vertex a
natural-number color
\(c(o)\), and take
\[
 T_i=\{o:c(o)<i\}.
\]
The preceding algebraic
complement identity still
holds. The phase strategy
for Player~I eventually
stays below one bound
on each compatible play.
Its preceding finite
prefix has bounded
colors.
Increasing the final
bound to include that
prefix gives
\[
 \exists b\in\mathbb N\ \forall k\in\mathbb N
                          \quad c(o_k)<b.
\]
This is the source's
bounded-color objective.

For Player~II, at each
phase start let \(M\)
be the largest color
observed on the whole
finite play so far,
including the current
boundary. Choose
\(i=M+1\).
The fixed-point equation
for \(J^*\) supplies
constrained reachability
of color at least
\(i\), while remaining
in \(J^*\).
Its target is initially
false, so the phase
takes a positive finite
number of rounds.
At its end the
historical maximum has
strictly increased.
Repeating produces an
unbounded color sequence
on every compatible
play.
Using the historical
maximum is essential
here: a later color
may fall below a
previous threshold.

The complementary regions
therefore characterize bounded
colors for Player~I and
unbounded colors for
Player~II. Writing the
color predicates out
recovers
\[
\begin{aligned}
 I&=\mu y\,\exists i\,\nu x
       \left(\left(EX\,AX\,x\land
                    \bigvee_{r<i}color_r\right)\lor y\right),\\
 II&=\nu y\,\forall i\,\mu x
       \left(\left(AX\,EX\,x\lor
                    \bigvee_{r\geq i}color_r\right)\land y\right).
\end{aligned}
\]
The color play in
Section~5 is the sequence
of \(O\)-vertices.
Under the earlier paired-color
convention the same
boundedness condition also
holds on the full
alternating vertex sequence.
This proves Theorem~5.1
\cite[\S5]{EmersonJutla1991}.

\subsubsection{Bradfield--Stirling's semantic and alternation synthesis}
\label{sec:p09-n-c-018-bradfield-stirling}

The accessible preprint is dated 15 November 2005.
Section and theorem locators here refer to that text;
the bibliographic entry identifies the 2007 published
chapter \cite{BradfieldStirling2007}.

\paragraph{Labelled systems and the modal language.}
Let \(\mathsf{Prop}\), \(\mathsf{Var}\), and
\(\mathsf{Act}\) be the proposition, variable, and
action-label sets. A labelled transition system is
\[
 \mathcal T=(S,(R_a)_{a\in\mathsf{Act}},V_{\mathsf{Prop}}),
 \quad R_a\subseteq S\times S,\quad
 V_{\mathsf{Prop}}:\mathsf{Prop}\to\mathcal P(S).
\]
An environment is
\(\eta:\mathsf{Var}\to\mathcal P(S)\).
Neither the system nor its successor sets must be finite.
The chapter's parsimonious syntax takes box and
greatest fixed point as primitive:
\[
 \varphi::=P\mid Z\mid\varphi\land\varphi
       \mid\neg\varphi\mid[a]\varphi\mid\nu Z.\varphi,
\]
with \(Z\) positive in its binding body.
Boolean constants are used with their usual meanings.
Disjunction, diamond, and least fixed point are derived
by complement duality; in positive normal form both
modalities and both binders may therefore be used
\cite[\S\S3.1--3.4]{BradfieldStirling2007}.

The Boolean and modal interpretations are the powerset
operations established above, now with \(R_a\) and
\(V_{\mathsf{Prop}}\) in place of \(I(a)\) and \(I(P)\).
In particular,
\[
\begin{aligned}
 \llbracket Z\rrbracket_\eta&=\eta(Z),&
 \llbracket P\rrbracket_\eta&=V_{\mathsf{Prop}}(P),\\
 \llbracket[a]\varphi\rrbracket_\eta
  &=\{s:\forall t\,(R_a(s,t)\Rightarrow
                         t\in\llbracket\varphi\rrbracket_\eta)\},\\
 \llbracket\nu Z.\varphi\rrbracket_\eta
  &=\bigcup\{U\subseteq S:
              U\subseteq\llbracket\varphi\rrbracket_{\eta[Z:=U]}\}.
\end{aligned}
\]
Thus the choice of primitive \(\nu\) rather than
\(\mu\) changes the syntax, not the fixed-point order.
For example,
\(\nu Z.(P\land[a]Z)\) means that every finite
\(a\)-reachable state satisfies \(P\), while
\(\mu Y.(P\lor(Q\land\langle a\rangle Y))\) requires
a finite path to \(P\) with \(Q\) at the earlier
positions. Dead ends retain the usual box-vacuity
and diamond-failure conventions.

\paragraph{Bisimulation preserves fixed-point truth.}
\label{sec:p09-mu-bisimulation-invariance}
For systems
\(\mathcal T_i=(S_i,(R_a^i)_{a\in\mathsf{Act}},V_i)\)
over the
same proposition and action sets, a bisimulation is
a relation \(B\subseteq S_1\times S_2\) such that,
whenever \(sBt\),
\[
\begin{aligned}
 s\in V_1(P)&\iff t\in V_2(P)
                         &&(P\in\mathsf{Prop}),\\
 R_a^1(s,s')&\Rightarrow
       \exists t'\,[R_a^2(t,t')\land s'Bt']
                         &&(a\in\mathsf{Act}),\\
 R_a^2(t,t')&\Rightarrow
       \exists s'\,[R_a^1(s,s')\land s'Bt']
                         &&(a\in\mathsf{Act}).
\end{aligned}
\]
States are bisimilar when some such relation connects them.
For formulas with free variables, call environments
\(\eta_1,\eta_2\) coherent along \(B\) when
\[
 sBt\ \Longrightarrow\
   (s\in\eta_1(Z)\iff t\in\eta_2(Z))
       \qquad(Z\in\mathsf{Var}).
\]
Then, for every modal fixed-point formula \(\varphi\),
\[
 sBt\quad\Longrightarrow\quad
 \bigl(s\in\llbracket\varphi\rrbracket_{\eta_1}
       \iff t\in\llbracket\varphi\rrbracket_{\eta_2}\bigr).
\]

Prove this by structural induction. Atomic propositions
use label agreement and variables use coherence.
Boolean operations preserve the equivalence. For a
diamond witness, transfer its edge by the appropriate
bisimulation clause and use the operand induction;
the converse direction uses the other clause.
Boxes follow either directly or by duality.
At a least binder, start both approximant chains at
the empty set. If the current approximants are
coherent, the body induction makes their images
coherent. Limit unions preserve coherence because
membership on either side is witnessed at the same
ordinal stage. Transfinite induction therefore
matches every stage. Choose a common ordinal beyond
both stabilization stages; their fixed points then
match. The greatest case uses descending chains
and intersections, or complement duality.
This proves the claim on arbitrary state sets.
For closed formulas no environment condition remains
\cite[\S4.2, thms.~4.1--4.2]{BradfieldStirling2007}.

\paragraph{The image-finite converse.}
A system is image-finite when
\(\{t:R_a(s,t)\}\) is finite for every state \(s\)
and action \(a\). Let \(\mathrm{HML}\) be the
variable-free, fixed-point-free modal language over
\(\mathsf{Prop}\) and \(\mathsf{Act}\), including Boolean
constants, conjunction, negation, and modalities.
In image-finite systems, two states satisfy the
same HML formulas exactly when they are bisimilar.

For the nontrivial direction, relate states exactly
when they agree on all HML formulas. Proposition
labels agree. Suppose \(s,t\) are related and
\(R_a^1(s,s')\). List the finitely many
\(a\)-successors of \(t\) as \(t_1,\ldots,t_m\).
If none is related to \(s'\), for each \(i\)
choose an HML formula \(\psi_i\) true at \(s'\)
and false at \(t_i\); negate a distinguishing
formula if its truth values have the reverse order.
Put \(\psi=\bigwedge_{i=1}^{m}\psi_i\), with the
empty conjunction equal to true. Then
\(\langle a\rangle\psi\) is true at \(s\)
and false at \(t\), a contradiction.
Thus the edge has a related match.
The symmetric argument supplies the back clause,
so the relation is a bisimulation.
The converse follows from the preceding invariance
proof even without image-finiteness.
Since HML is a fragment of the modal \(\mu\)-calculus,
agreement on all closed fixed-point formulas also
characterizes bisimilarity in image-finite systems
\cite[\S4.2, thm.~4.3]{BradfieldStirling2007}.

The finiteness hypothesis is used at exactly the
conjunction of distinguishing formulas. An infinite
successor family would require an infinitary
conjunction, which is not an HML formula.

\paragraph{Unravelling gives a tree model.}
Fix a state \(s_*\) in a labelled system and put \(s_0=s_*\).
The unravelled states are finite labelled paths
\[
 \pi=(s_*,a_1,s_1,\ldots,a_n,s_n),
 \qquad R_{a_i}(s_{i-1},s_i),
\]
including the length-zero path \((s_*)\).
Let \(\operatorname{last}(\pi)=s_n\).
There is an \(a\)-edge from \(\pi\) to
\(\pi\cdot(a,t)\) precisely when
\(R_a(\operatorname{last}(\pi),t)\).
Give each path the proposition labels of its last
state. Every nonroot path has one predecessor
obtained by deleting its last edge and endpoint,
so this is a rooted tree.

The relation
\[
 \{(\pi,\operatorname{last}(\pi)):\pi
                        \text{ is an unravelled state}\}
\]
is a bisimulation with the original system:
labels were copied, each tree edge projects to
its original edge, and every original successor
can be appended to the current path.
Bisimulation invariance therefore gives the same
closed formulas at \((s_*)\) and \(s_*\).
For an open formula, use the pulled-back environment
\[
 \eta_{\mathrm{tree}}(Z)
   =\{\pi:\operatorname{last}(\pi)\in\eta_{\mathrm{original}}(Z)\}.
\]
The two environments are coherent along the displayed
bisimulation, so the same conclusion holds for
open formulas with these environments.
In particular, every satisfiable modal fixed-point
formula has a tree model. This argument claims no
finite branching bound when the original successor
sets are infinite \cite[\S4.2, prop.~4.4]{BradfieldStirling2007}.

\paragraph{The complete model-checking game.}
\label{sec:p09-mu-evaluation-game}
Fix a formula \(\varphi\) in positive normal form,
with negation only at proposition constants.
Rename its bound variables distinctly, also avoiding
the names of free variables. A position is
\((s,\psi)\), where \(s\in S\) and \(\psi\)
is an original subformula occurrence of \(\varphi\)
with its binder links. The start is \((s_*,\varphi)\).
There are two players, \(\exists\) and \(\forall\).
The moves are
\[
\begin{array}{lll}
 \text{Position}&\text{Choice owner}&\text{Successors}\\
 (s,\alpha\lor\beta)&\exists&(s,\alpha),(s,\beta)\\
 (s,\alpha\land\beta)&\forall&(s,\alpha),(s,\beta)\\
 (s,\langle a\rangle\alpha)&\exists&
                  (t,\alpha)\text{ for }R_a(s,t)\\
 (s,[a]\alpha)&\forall&
                  (t,\alpha)\text{ for }R_a(s,t)\\
 (s,\sigma X.\alpha)&\text{forced}&(s,\alpha)\\
 (s,X),\ X\text{ bound}&\text{forced}&
                  (s,\sigma X.\alpha)\text{ at its binder}.
\end{array}
\]
Assign forced positions to \(\forall\); their
unique successor makes this ownership immaterial.
At a proposition literal or Boolean constant,
play ends and \(\exists\) wins exactly when
that literal or constant is true at \(s\).
For the semantic induction, permit free variables
as additional terminal positions: \((s,Z)\) is
won by \(\exists\) exactly when \(s\in\eta(Z)\).
A diamond with no successor is lost by \(\exists\);
a box with no successor is won by \(\exists\).
These cases give the ordinary modal semantics
at dead ends \cite[\S4.4]{BradfieldStirling2007}.

An infinite play must visit some bound-variable
position infinitely often. Otherwise, after
the last variable return, every move would
descend a finite syntax tree. Among the
infinitely recurring variables there is an
outermost binder. Indeed, moving infinitely
often between incomparable binding subtrees
would require infinitely many returns through
an enclosing variable; that variable would
then recur farther outside.
Player \(\exists\) wins the infinite play
exactly when this outermost recurring variable
is a \(\nu\)-variable.

This is a parity condition. If \(d_X\) is the
number of proper enclosing binders of \(X\),
give each bound-variable position priority
\[
 \Omega(s,X)=
 \begin{cases}
  2d_X+1,&X\text{ is bound by }\mu,\\
  2d_X,&X\text{ is bound by }\nu.
 \end{cases}
\]
Give other nonterminal positions priority
\(2b+2\), where \(b\) is the number of binders
in \(\varphi\). Outer binders have lower
priority than inner ones, and every infinite
play has a recurring variable priority below
the default. Thus \(\exists\) wins exactly
when the least priority recurring infinitely
often is even. This assignment is sufficient;
it does not claim to minimize the number of
priorities or define the alternation-depth
measures considered later.

A strategy specifies a legal choice after
each finite play ending at a nonterminal
position owned by its player. It is winning if every maximal
play following it is won by that player.
A positional strategy depends only on the
current position, not on the preceding play.

\paragraph{Truth and winning strategies.}
\label{sec:p09-mu-game-truth-proof}
For every state and environment,
\[
 s\in\llbracket\varphi\rrbracket_\eta
 \quad\Longleftrightarrow\quad
 \exists\text{ has a winning strategy from }(s,\varphi);
\]
if the left side is false, \(\forall\) has a
winning strategy. We prove both assertions
simultaneously by structural induction, allowing
the free-variable exits just defined.

At constants, literals, and free variables the
claim is the terminal rule. A true disjunction
has a true disjunct for \(\exists\) to choose;
if both are false, \(\forall\) uses the induction
strategy in whichever one is chosen.
A true conjunction lets \(\exists\) answer
either choice by \(\forall\); a false one
lets \(\forall\) choose a false conjunct.
The modal cases are the same arguments over
successor states: an existential witness gives
the diamond choice, and a failed universal
condition gives the box refutation. Empty
successor sets are already covered.

For a binder \(\mu X.\psi\), let
\[
 F(A)=\llbracket\psi\rrbracket_{\eta[X:=A]},
 \qquad L=\mu F,
 \qquad A_0=\varnothing,\quad
 A_{\alpha+1}=F(A_\alpha),\quad
 A_\lambda=\bigcup_{\beta<\lambda}A_\beta.
\]
If \(s\in L\), its first membership stage
is a successor \(\alpha+1\). In the body
game with \(X\) interpreted by \(A_\alpha\),
the formula \(\psi\) is true at \(s\), so
the induction hypothesis gives an \(\exists\)
winning strategy. Use it until either play
continues forever inside the body or a free
\(X\)-exit is reached.
Such an exit can be winning only at
\(t\in A_\alpha\). In the actual bound game,
return to the binder at \(t\) and restart
the same construction using \(t\)'s first
membership stage. This rank is strictly
smaller than \(\alpha+1\).
There cannot be infinitely many such returns,
since they would strictly descend through
ordinals. After the final return, the body
strategy wins, whether by a terminal outcome
or its own infinite-play condition.
Changing finitely many initial priorities
cannot change a parity outcome.

If \(s\notin L=F(L)\), the body is false
under \(X:=L\). Use its induction strategy
for \(\forall\). Any free \(X\)-exit reached
while following that strategy lies outside
\(L\); restart there. With finitely many
restarts the final body strategy wins.
With infinitely many, \(X\) is the outermost
recurring variable and has odd priority,
so \(\forall\) wins again.

For a binder \(\nu X.\psi\), write \(U=\nu F\).
If \(s\in U=F(U)\), use the body's
\(\exists\)-strategy under \(X:=U\).
Every reached \(X\)-exit stays in \(U\),
so it can be restarted. Finitely many
restarts leave a winning body play;
infinitely many recur through the outer
even-priority variable and are also won
by \(\exists\).

Finally suppose \(s\notin U\). Use the
descending approximants
\[
 B_0=S,\qquad B_{\alpha+1}=F(B_\alpha),
 \qquad B_\lambda=\bigcap_{\beta<\lambda}B_\beta.
\]
The first stage excluding \(s\) is a
successor \(\alpha+1\). At \(s\), the body
is false under \(X:=B_\alpha\), giving a
\(\forall\)-strategy by induction.
Any \(X\)-exit it reaches lies outside
\(B_\alpha\), and hence has a strictly
smaller first-exclusion rank. Restart
using that rank. Infinite restarting
is impossible, and the final body play
is won by \(\forall\).

These four binder cases complete the
induction. When a body is considered as
an open formula, removing its outer binder
shifts all inner binder priorities by an
even amount and leaves their order unchanged.
Thus the induction uses the same winner
on body plays that never return through
the newly bound variable. The constructed
strategies may use their current ordinal
rank; positionality is a separate game
theorem, not an assumption in this proof
\cite[\S4.4, thm.~4.6]{BradfieldStirling2007}.

\paragraph{Finite arenas and two basic examples.}
If \(S\) is finite, the arena has at most
\(|S|\,|\operatorname{Sub}(\varphi)|\)
ordinary positions. Add two absorbing
vertices of priorities \(0\) and \(1\),
and redirect each terminal outcome to
the absorber for its winner, assigning
the former terminal the default priority.
These added moves are forced. This gives
a finite total parity game with two
additional vertices and the same winner.
The priorities record fixed-point
dependencies, not the number of
approximant iterations used by an
implementation.

At a one-state \(a\)-loop where \(P\)
is false, the game for
\(\mu X.(P\lor\langle a\rangle X)\)
is lost by \(\exists\). Choosing \(P\)
loses immediately; continuing the loop
visits the odd \(X\)-priority forever.
In contrast, \(\nu X.\langle a\rangle X\)
is won there by repeatedly choosing the
loop, whose recurring \(X\)-priority is
even.

For a nested example, consider
\[
 \nu X.\mu Y.\bigl(
     (P\land\langle a\rangle X)
       \lor(\neg P\land\langle a\rangle Y)\bigr).
\]
On a two-state \(a\)-cycle, with \(P\)
true exactly at one state, \(\exists\)
chooses the disjunct matching the current
label. Any literal challenge by \(\forall\)
ends in a true literal. A play that always
continues visits \(X\) and \(Y\) repeatedly,
with priorities \(0\) and \(3\), and is
won by \(\exists\). On a one-state loop
where \(P\) is always false, continuing
instead visits only the inner odd
\(Y\)-priority infinitely often, and
\(\forall\) wins. This distinguishes
repeated completion of an inner least
obligation from postponing it forever.

\paragraph{Positional determinacy, including infinite arenas.}
\label{sec:p09-parity-positional-proof}
We prove the game theorem needed to remove
history dependence from the semantic strategies.
Work in ZFC, allowing the choices of successors
and strategies below. The arena is a set
\[
 G=(V,V_0,V_1,E,c),\qquad
 V=V_0\mathbin{\dot\cup}V_1,\quad
 E\subseteq V^2,\quad c:V\to\{0,\ldots,d\},
\]
where every \(E(v)=\{w:(v,w)\in E\}\) is nonempty
and \(d\) is a nonnegative integer. No finiteness
of \(V\) or of successor sets is required.
For this proof use the largest-recurring
convention: player \(i\) wins when the largest
priority seen infinitely often has parity \(i\).
There is always such a priority, even if no
individual vertex recurs.

A positional \(i\)-dominion is a subset \(D\subseteq V\)
with a choice function on \(D\cap V_i\)
selecting successors in \(D\), such that
every opponent edge from \(D\) also stays
in \(D\), and every play following that
function is won by \(i\).
We prove the stronger statement that \(V\)
partitions into two positional dominions.
Thus one strategy works simultaneously
from every vertex of each winning region.

\emph{Attraction and ranks.}
In a total induced subarena \(G[U]\), the
\(i\)-attractor of \(X\subseteq U\) is
computed by
\[
\begin{aligned}
 B_0&=X,\\
 B_{\xi+1}&=B_\xi
  \cup\{v\in U\cap V_i:E(v)\cap U\cap B_\xi\ne\varnothing\}\\
 &\hspace{12mm}\cup
    \{v\in U\cap V_{1-i}:E(v)\cap U\subseteq B_\xi\},\\
 B_\lambda&=\bigcup_{\xi<\lambda}B_\xi
                 \qquad(\lambda\text{ a nonzero limit ordinal}).
\end{aligned}
\]
There must be a stationary stage: otherwise
the disjoint nonempty increments would give
an injection of an arbitrarily large ordinal
into \(U\). The stationary set
\(\operatorname{Attr}_i^{G[U]}(X)\) is
the least superset of \(X\) closed under
the two predecessor clauses. This makes
attraction monotone and idempotent.
Totality gives \(\operatorname{Attr}_i(\varnothing)
=\varnothing\).

A vertex outside \(X\) first enters at a
successor stage. At an \(i\)-owned vertex,
choose a successor of smaller entry rank;
at an opposing vertex, every successor
in the subarena has smaller rank.
These ranks strictly decrease until \(X\)
is reached, which must happen after
finitely many moves. No uniform finite
time bound is asserted.
If \(A\) is this attractor, its complement
satisfies
\[
\begin{aligned}
 v\in(U\setminus A)\cap V_i
     &\ \Longrightarrow\ E(v)\cap U\subseteq U\setminus A,\\
 v\in(U\setminus A)\cap V_{1-i}
     &\ \Longrightarrow\ E(v)\cap(U\setminus A)\ne\varnothing.
\end{aligned}
\]
Consequently \(G[U\setminus A]\) is total.
Attractor target vertices themselves need
not have all successors in the attractor;
their exits will be checked separately.

\emph{Removal of lower-priority dominions.}
Induct on \(d\). Empty arenas are trivial.
At \(d=0\), every play is won by player
\(0\), so \(V\) itself is that player's
dominion with any legal choices.
For \(d>0\), put \(p=d\bmod2\) and
\(q=1-p\).
Construct increasing \(q\)-attractor-closed
sets \(R_\alpha\) in the original arena.
Start at \(R_0=\varnothing\), and define
\[
\begin{aligned}
 H_\alpha&=V\setminus R_\alpha,\\
 T_\alpha&=\{v\in H_\alpha:c(v)=d\},\\
 A_\alpha&=\operatorname{Attr}_p^{G[H_\alpha]}(T_\alpha),\\
 K_\alpha&=H_\alpha\setminus A_\alpha.
\end{aligned}
\]
Both residual arenas are total by the
complement facts above. The arena on
\(K_\alpha\) has fewer priorities.
By induction it partitions into positional
\(p\)- and \(q\)-dominions. Let \(W_\alpha\)
be its \(q\)-dominion, with strategy
\(\tau_\alpha\).
If \(W_\alpha=\varnothing\), stop.
Otherwise set
\[
 R_{\alpha+1}
   =\operatorname{Attr}_q^G(R_\alpha\cup W_\alpha).
\]
At a nonzero limit stage use
\[
 U_\lambda=\bigcup_{\beta<\lambda}R_\beta,
 \qquad
 R_\lambda=\operatorname{Attr}_q^G(U_\lambda).
\]
The second attraction at limits is essential.
For example, a \(p\)-vertex may have
successors \(w_0,w_1,\ldots\), each a
\(q\)-owned self-loop. Each finite set
\(\{w_0,\ldots,w_{n-1}\}\) is \(q\)-attractor-closed,
but their union must additionally attract
the \(p\)-vertex. Removing just the union
would leave it with no successor.
The displayed limit rule preserves totality
without assuming finite branching.

Each non-stopping successor stage adds
the nonempty seed \(W_\alpha\), disjoint
from \(R_\alpha\). Choosing one vertex
from every such seed would inject the
stage indices into \(V\).
Thus the construction stops at some
set ordinal \(\theta\), before an ordinal
that cannot inject into \(V\).
All strategy selections consequently
belong to a set-indexed family.

The seed's exits obey
\[
\begin{aligned}
 v\in W_\alpha\cap V_q
     &\ \Longrightarrow\ \tau_\alpha(v)\in W_\alpha,\\
 v\in W_\alpha\cap V_p
     &\ \Longrightarrow\ E(v)\subseteq R_\alpha\cup W_\alpha.
\end{aligned}
\]
For the second line, a successor outside
\(H_\alpha\) is already in \(R_\alpha\).
A successor inside \(H_\alpha\) cannot
enter \(A_\alpha\), because this is a
\(p\)-owned vertex outside the
\(p\)-attractor. It lies in \(K_\alpha\),
and the inductive dominion property
then puts it in \(W_\alpha\).

\emph{One strategy on the removed region.}
For \(v\in R_\theta\), let
\[
 e(v)=\min\{\gamma:v\in R_\gamma\}.
\]
This stage is positive and may be a
limit ordinal. At entry stage
\(\gamma=\alpha+1\), use
\(\tau_\alpha\) on the fresh seed
\(W_\alpha\), and elsewhere use a
fixed rank-decreasing \(q\)-attraction
move toward \(R_\alpha\cup W_\alpha\).
At a limit entry stage \(\lambda\),
use the rank-decreasing attraction
move toward \(U_\lambda\).
The latter vertex is not in the union
target, or it would have entered
earlier, so its rank is positive.
These choices depend only on \(v\)
and define one positional strategy.

Every move following it, including
every opponent move, stays inside
\(R_{e(v)}\). On a seed this follows
from the exit conditions; off the
seed it follows from the attraction
rank. Hence entry stages along any
compatible play are nonincreasing.
They are eventually constant:
the set of stages occurring on the
play has a least member, and after
its first occurrence none can be
smaller.

A constant limit-stage tail is
impossible, because its attraction
ranks would strictly decrease
forever outside the union target.
On a constant successor-stage
tail, say \(\alpha+1\), rank
descent forces arrival at
\(W_\alpha\). Any later departure
would enter \(R_\alpha\) and
lower the entry stage, so the
tail stays in \(W_\alpha\) and
follows its inductively winning
strategy. Its finite prefix does
not change parity. Therefore
\(R_\theta\) is a positional
\(q\)-dominion.

\emph{One strategy on the remainder.}
At stopping, \(W_\theta\) is empty,
so the whole arena \(K_\theta\)
is an inductive \(p\)-dominion.
On \(H_\theta\), player \(p\)
uses its strategy in \(K_\theta\),
the \(p\)-attraction strategy
on \(A_\theta\setminus T_\theta\),
and any fixed residual successor
at its vertices of \(T_\theta\).
All those choices stay in
\(H_\theta\). Opponent vertices
there also have every successor
in \(H_\theta\), since
\(R_\theta\) is \(q\)-attractor-closed.

If a play visits \(T_\theta\)
infinitely often, its largest
recurring priority is \(d\),
so \(p\) wins. Otherwise,
after its last such visit it
cannot enter \(A_\theta\):
attraction would force another
visit. Its tail stays in
\(K_\theta\) and is won by the
inductive strategy. Thus
\(H_\theta\) is a positional
\(p\)-dominion, completing the
partition and induction.

These dominions are the actual
winning regions against arbitrary
opponent strategies. Otherwise,
playing two allegedly winning
strategies against each other
would produce a play required
to have both winning parities.
This proves positional determinacy.

\paragraph{Returning to the least-recurring convention.}
For any finite-range priority map \(\Omega\),
choose an even integer \(M\geq\max\Omega(V)\)
when \(V\ne\varnothing\), and replace it by
\(c(v)=M-\Omega(v)\). Writing \(\operatorname{Inf}\)
for the set of priorities occurring infinitely often,
one has on every infinite play
\[
 \max\operatorname{Inf}_c
      =M-\min\operatorname{Inf}_\Omega.
\]
Evenness of \(M\) preserves the
winning parity; an odd reversal would
exchange the players.
The empty arena needs no choices.
Prescribed finite terminal outcomes
are handled by the absorbing vertices
already defined.

Apply the positional theorem to the
totalized model-checking arena, whose
priority range is finite even when
the state set is infinite. Restrict
its strategies back to nonterminal
positions of the original game.
Together with the semantic proof,
this gives
\[
 s\in\llbracket\varphi\rrbracket_\eta
 \quad\Longleftrightarrow\quad
 \exists\text{ has a positional winning strategy
 from }(s,\varphi).
\]
This supplies the history-free conclusion
of Theorem~4.6, not merely determinacy
among history-dependent strategies
\cite[\S4.4, thm.~4.6]{BradfieldStirling2007}.
The even reversal also relates this
convention to the largest-recurring
one used by Emerson--Jutla
\cite[\S4.2]{EmersonJutla1991}.

\paragraph{Binding contexts for annotated subformulas.}
\label{sec:p09-mu-contextual-truth}
An annotation of a raw subformula must
retain its original binding context.
Fix a finite well-named positive-normal
formula \(\varphi\), a model \(\mathcal T\),
and an environment \(\eta\) for its
genuinely free variables.
Bound names are globally distinct
and disjoint from genuinely free names.
Let \(C\) be the original subformula
occurrences, with their binder links.
An occurrence below a binder may be
open when read in isolation even
though its variables are bound
in \(\varphi\).

The contextual environment is explicit.
List the binding occurrences
\[
 b_i=\sigma_iX_i.\beta_i\quad(1\leq i\leq b)
\]
with every enclosing binder before
its proper descendants.
Starting from \(\eta\), successively
assign \(X_i\) the denotation of
\(b_i\) with its already assigned
outer parameters.
Initially unassigned bound names
may have arbitrary values: any
free bound name of \(b_i\)
belongs to an enclosing binder,
whereas its descendant binders
bind their own names locally.
Later assignments therefore do
not change \(b_i\)'s denotation.
Call the resulting environment
\(\xi_{\varphi,\eta}\), or \(\xi\).
In particular,
\[
 \xi(X_i)=\llbracket b_i\rrbracket_\xi
         =\llbracket\beta_i\rrbracket_\xi.
\]
The second equality is the
fixed-point equation.
For the whole formula,
\(\llbracket\varphi\rrbracket_\xi
 =\llbracket\varphi\rrbracket_\eta\),
since its bound-name assignments
are shadowed by their binders.

For every original occurrence
\(\psi\), the linked evaluation
game of \(\varphi\), started
at \((s,\psi)\), is won by
\(\exists\) exactly when
\[
 s\in\llbracket\psi\rrbracket_\xi.
\]
Here outer bound variables still
return to their original binders;
they are not silently made
free terminals of value \(\eta\).
We prove this contextual claim
first for binding occurrences,
by induction on their enclosing
binder depth.
Consider the ordinary standalone
evaluation game for \(b_i\),
with its free outer variables
given their values in \(\xi\).
Its winner is known from the
truth/game theorem.
Until the first such outer
variable is encountered, its
moves are those of the linked
game. The internal binder order
and parity are unchanged by
the finite shift of nesting depth.

If a play following the standalone
winning strategy reaches an outer
variable \(X_j\), its terminal
value favors that strategy's
player. By the induction
hypothesis, the same player
has a winning linked-game
strategy from the enclosing
binder \(b_j\) at that state.
Take the variable's forced
return and switch permanently
to that strategy.
A finite preceding history
does not change the outcome.
If no outer-variable exit
occurs, the standalone and
linked plays already agree
in their terminal outcome
or recurring parity.
This argument works for
both players, proving the
claim for \(b_i\) without
assuming it for its descendants.
Once binding occurrences are
covered, the same first-exit
argument applies to any
occurrence \(\psi\).
It proves the displayed claim
and its false-formula dual.

\paragraph{Locally consistent pre-model annotations.}
\label{sec:p09-mu-premodel}
A pre-model annotation is a
set \(A\subseteq S\times C\).
Write \(s@\psi\) for
\((s,\psi)\in A\).
Its local conditions are:
\begin{itemize}
\item a labelled proposition
literal is true at \(s\);
\(\bot\) is never labelled;
and a labelled genuinely free
variable \(Z\) satisfies
\(s\in\eta(Z)\);
\item \(s@(\alpha\land\beta)\)
holds exactly when both
corresponding child occurrences
are labelled at \(s\);
\(s@(\alpha\lor\beta)\) holds
exactly when at least one
of those children is labelled;
\item a labelled \([a]\psi\)
has \(t@\psi\) at every
\(a\)-successor \(t\);
a labelled \(\langle a\rangle\psi\)
has such a labelled successor;
\item a labelled binder has
its body labelled at the
same state; a labelled bound
variable has its original
binder labelled there.
\end{itemize}
The conditions are about
specified original occurrences,
not arbitrary formulas with
the same printed text.
Bound-variable labels are
checked by their unfolding
obligations, not by presupposing
their unknown semantic truth.
The contextual environment
\(\xi\) describes the theorem's
conclusion, not an additional
local test on those labels.

Choose one function \(f\)
that selects a labelled
disjunct at every labelled
disjunction and a labelled
actual successor at every
labelled diamond.
The choices depend on the
state and original occurrence.
Conjunctions retain both
dependencies and boxes retain
all their actual successor
dependencies. Binders move to
their bodies; bound variables
return to their binding
occurrences, then to those
bodies, at the same state.
These rules define a directed
dependency graph on \(A\).
Including the binder as an
administrative vertex subdivides
the direct variable-to-body
dependency without changing
the recurring variable order.
This makes every dependency
step a step of the linked
evaluation game
\cite[\S3.6]{BradfieldStirling2007}.

A trail is a maximal dependency
path, finite or infinite.
Every finite trail ends at
a true literal, a true
constant, a true genuinely
free-variable label, or a
box with no successor.
The other local rules
always supply a continuation.
Give dependency vertices the
earlier least-recurring
evaluation priorities, with
an even default above all
variable priorities.
The pre-model is well-founded
when every infinite trail
starting at any labelled pair
has even least recurring
priority, equivalently when
its outermost infinitely
recurring variable is
\(\nu\)-bound.
This terminology does not
exclude all infinite trails.
For \(\nu X.X\), a singleton
model with no action edges
may label both the binder
and \(X\); their infinite
administrative trail is accepting.
The same local annotations for
\(\mu X.X\) give an odd
trail and fail well-foundedness.

\paragraph{The pre-model truth theorem and its converse.}
In a well-founded pre-model,
use \(f\) for the existential
evaluation choices.
Every maximal compatible evaluation
play from a labelled position
stays in the annotation and
is a dependency trail.
Finite plays win by local
consistency; infinite plays
win by well-foundedness.
Thus every labelled position
has a winning linked-game
continuation, and contextual
truth gives
\[
 s@\psi\quad\Longrightarrow\quad
 s\in\llbracket\psi\rrbracket_{\xi_{\varphi,\eta}}.
\]
In particular, if \(s_0@\varphi\),
then
\(s_0\in\llbracket\varphi\rrbracket_\eta\).
This states the truth of
raw annotations with their
binding context made explicit.

Conversely, suppose
\(s_0\in\llbracket\varphi\rrbracket_\eta\).
Let \(A\) be the entire
\(\exists\)-winning region
of the linked evaluation
game. The contextual truth
lemma identifies these labels
with their denotations under
\(\xi\). Boolean semantics
gives both local Boolean
equivalences; modal semantics
gives the required successors;
the fixed-point equations
give the binder and variable
conditions.
Choose one positional strategy
winning throughout this region,
as supplied by positional
determinacy.
Its disjunction and diamond
choices define \(f\).
All opponent successors remain
in the winning region, and
every infinite path following
this one strategy wins.
This is a well-founded
pre-model on the original
structure, containing the
required label \(s_0@\varphi\).
We have proved
\[
 s_0\in\llbracket\varphi\rrbracket_\eta
 \quad\Longleftrightarrow\quad
 \text{a well-founded pre-model on \(\mathcal T\)
       contains \(s_0@\varphi\)}.
\]
This is the precise same-structure
pre-model equivalence
\cite[thms.~3.1--3.2]{BradfieldStirling2007}.

Using one coherent choice
function matters. For
\[
 \mu X.(P\lor\langle a\rangle X)
\]
on a two-state \(a\)-cycle
with \(P\) true at both
states, one winning strategy
may stop at the first
state and continue at the
second; another may make
the opposite choices.
Each wins from both states.
Taking the continuing choice
from each produces an
infinite odd \(X\)-trail.
Separate winning choices
cannot therefore be combined
without checking the resulting
global dependency condition.

\paragraph{Ordinal certificates for the trail condition.}
\label{sec:p09-mu-premodel-signatures}
The global trail condition
has an exact local certificate.
Let \(H=(V,E)\) be any
set-sized directed graph
with a finite priority map
\(p:V\to\{0,\ldots,d\}\).
Let
\[
 J=\{j_1<\cdots<j_k\}
\]
be its occurring odd priorities.
A signature assigns a tuple
of ordinals
\(\mathbf r(v)=(r_{j_1}(v),\ldots,r_{j_k}(v))\)
to every vertex.
Write \(\mathbf r_{\leq m}(v)\)
for its prefix of coordinates
indexed by odd \(j\leq m\).
The smallest priority is
the most significant
lexicographic coordinate.
The empty prefix is equal
to itself.

Every infinite path of \(H\)
has even least recurring
priority if and only if
there are signatures such
that every edge \(v\to w\)
satisfies
\[
 \mathbf r_{\leq p(v)}(v)
       \ \geq_{\mathrm{lex}}\
 \mathbf r_{\leq p(v)}(w),
\]
strictly when \(p(v)\) is odd.
Only the prefix determined
by the source priority
is compared. Later coordinates
may increase on that edge.
We prove both directions,
including graphs with
arbitrary branching and
terminal vertices.

For soundness, suppose a
path had odd least recurring
priority \(j\).
Discard a finite prefix
containing all appearances
of smaller priorities.
The signature prefix through
\(j\) is now lexicographically
nonincreasing on every edge
and strictly decreases each
time source priority \(j\)
occurs. Those strict steps
occur infinitely often.
But a nonincreasing sequence
of finite ordinal tuples
eventually stabilizes:
its first coordinate reaches
its least attained value
and is then constant;
on that tail the second
coordinate does the same,
and so on through the
finitely many coordinates.
This is a contradiction.

\paragraph{Constructing the ordinal signatures.}
Assume now that all infinite
paths have even least
recurring priority.
For each odd \(j\), define
\(v\,R_j\,w\) when there
is a finite nonempty path
\[
 v=v_0\to v_1\to\cdots\to v_\ell=w
\]
with \(p(v_i)\geq j\)
for every source \(i<\ell\),
and \(p(v_i)=j\) for
at least one such source.
The final endpoint may
have lower priority.
A path used in this
relation cannot pass through
a lower-priority vertex
as an edge source.

The relation \(R_j\) is
well-founded in the direction
of the displayed paths.
Otherwise an infinite chain
of \(R_j\)-links, with
finite witnesses chosen for
its links, could be
concatenated into an infinite
path having no priority
below \(j\) and infinitely
many occurrences of \(j\).
Every witness is finite
and nonempty, so this
is genuinely an infinite
path. A lower-priority
endpoint cannot continue
with another \(R_j\)-link.
The concatenated path
would have odd least
recurring priority, contrary
to the hypothesis.
We use the usual classical
set-theoretic background;
dependent choice suffices
to extract this chain
from a non-well-founded
set relation.

Well-founded recursion gives
\[
 r_j(v)=\sup\{r_j(w)+1:v\,R_j\,w\},
 \qquad \sup\varnothing=0.
\]
Take an original edge
\(v\to w\).
If \(p(v)=j\), that edge
itself witnesses \(vR_jw\),
so \(r_j(v)>r_j(w)\).
If \(p(v)>j\), prepend
the edge to any witness
for \(wR_ju\).
This proves
\[
 \{u:wR_ju\}\subseteq\{u:vR_ju\},
\]
and the rank equations
give \(r_j(v)\geq r_j(w)\).
For \(p(v)<j\), no
comparison is required.
Thus the constructed ranks
actually satisfy the stronger
coordinatewise conditions
\[
\begin{aligned}
 j\leq p(v)&\ \Longrightarrow\
                         r_j(v)\geq r_j(w),\\
 p(v)=j&\ \Longrightarrow\
                         r_j(v)>r_j(w).
\end{aligned}
\]
They imply the required
lexicographic prefix inequalities.
No previously chosen coordinate
or assumed bound on finite
path lengths enters this
construction.

The graph is a set and
there are finitely many
coordinates, so all these
ordinal values lie below
one common ordinal, for
example
\[
 \kappa=\sup\bigl(\{1\}\cup
       \{r_j(v)+1:j\in J,\ v\in V\}\bigr).
\]
Hence \(\mathbf r:V\to\kappa^k\)
is an ordinary set-valued
certificate. With no odd
priorities its tuple is
empty. A dead end has
rank zero in every
coordinate and imposes
no outgoing-edge comparison.

\paragraph{Resets, limit ranks, and pre-model certification.}
The accepting two-vertex cycle
\[
 u\ (p=0)\longrightarrow
 v\ (p=1)\longrightarrow u
\]
has ranks \(r_1(u)=0\)
and \(r_1(v)=1\).
The odd-priority edge
decreases the rank, while
the priority-zero edge
may reset it upward:
its compared prefix is
empty. Requiring the
whole tuple to decrease
on every dependency
would reject this valid
recurrence pattern.

Natural-number ranks do
not suffice in general.
Give a root \(r\)
priority \(1\), and let
it lead to chains of
every finite positive
length, whose vertices
also have priority \(1\).
Each chain ends at a
dead priority-zero vertex.
There is no infinite
path. The successive chain
ranks are \(1,2,\ldots\),
while the root rank is
\(\omega\), strictly above
each finite successor rank.
For a finite graph with
\(N\geq1\) vertices, by
contrast, each \(R_j\)
has rank at most \(N-1\),
because a chain cannot
repeat a vertex.

Apply the certificate theorem
to the dependency graph
of a pre-model.
Local terminal truth and
the displayed edge comparisons
then imply the full
pre-model truth theorem.
Conversely every well-founded
pre-model has such ordinal
signatures. If only a
designated root is being
certified, it is enough
to rank its reachable
dependency graph; this
does not certify unrelated
labelled components.

These ranks count selected
regeneration patterns. They
need not equal the least
semantic approximant at
which a state enters
a fixed point.
For instance, \(\mu X.P\)
can hold at its first
semantic approximant although
no return through \(X\)
occurs at all.
The signature construction
provides local ordinal
evidence for the trail
condition, with the binding
interpretation and finite
terminal checks kept explicit.

\paragraph{A bounded-branching tree model.}
Let a closed formula \(\varphi\) hold at
\(s_*\), and let \(d_\varphi\) be the
number of its diamond occurrences.
Choose one positional strategy winning
on the entire \(\exists\)-winning region
of its evaluation game.
At every model state \(s\), retain the
edge selected by this strategy for each
winning diamond position
\((s,\langle a\rangle\psi)\), and discard
all other model edges. There are at
most \(d_\varphi\) retained labelled
edges from any state.

The same game strategy still wins
from \((s_*,\varphi)\) in this
restricted model. At a winning
diamond its selected edge remains.
At a box the opponent has fewer
choices; if none remain, the box
becomes an immediate \(\exists\)-win.
Every continuing play was already
a play of the original game with
the same priorities and strategy.
Literal values and forced syntax
moves are unchanged. Thus no new
losing play is introduced.
The truth/game equivalence proves
that the restricted model still
satisfies \(\varphi\).
Unravel it as above. Each tree
node has at most \(d_\varphi\)
children, proving the bounded-branching
tree-model property, including
the zero-diamond case. The tree
need not be finite or fully
branching \cite[\S4.2, prop.~4.4]{BradfieldStirling2007}.

\paragraph{Modal alternating parity automata.}
\label{sec:p09-modal-parity-automata}
Fix a finite proposition set \(P_0\) and let
\(\Sigma=\mathcal P(P_0)\). A modal parity automaton is
\[
 \mathcal A=(Q,\Sigma,\delta,q_0,\Omega),
 \qquad q_0\in Q,\quad
 \Omega:Q\to\{0,\ldots,d\},
\]
where \(Q\) is finite and nonempty.
The transition map has type
\[
 \delta:Q\times\Sigma\to\mathcal B(Q),
\]
with the positive depth-one modal language
\[
 B::=\top\mid\bot\mid
       \langle a\rangle q\mid[a]q
       \mid B\land B\mid B\lor B
       \quad(q\in Q,\ a\in\mathsf{Act}).
\]
Every reference to a next control state occurs under
one modality. Boolean combinations may branch, but
there is no unguarded control-state atom in
\(\mathcal B(Q)\)
\cite[\S5.3]{BradfieldStirling2007}.

In a model \(\mathcal T\), let
\(\ell(s)=\{P\in P_0:s\in V_{\mathsf{Prop}}(P)\}\).
The acceptance game starts at \((s_*,q_0)\).
From \((s,q)\), take the forced step to the finite
expression \(\delta(q,\ell(s))\) at \(s\).
Player \(\exists\) chooses a disjunct and
\(\forall\) chooses a conjunct.
At \(\langle a\rangle r\), player \(\exists\)
chooses an \(a\)-successor \(t\) and play continues
at \((t,r)\); at \([a]r\), player \(\forall\)
chooses such a successor.
The constants \(\top,\bot\) are respectively
winning and losing terminals for \(\exists\).
An empty diamond loses and an empty box wins
for \(\exists\).

The expression evaluation between control
positions is finite. Therefore every infinite
play visits control positions infinitely often;
\(\exists\) wins it exactly when the least
priority recurring in their \(Q\)-sequence
is even. Equivalently, give expression positions
a default priority larger than \(d\), and use
the ordinary parity condition on all positions.
Expression occurrences may be tagged with the
current control state to keep the arena fully
specified. The automaton accepts
\((\mathcal T,s_*)\) when \(\exists\) has a
winning strategy.

\paragraph{Runs and simultaneous successor obligations.}
The game definition agrees with the run formulation
in which the automaton creates several successor
obligations at once. At a model state \(s\),
a mixture is a set
\[
 M_s\subseteq
 \{(t,r)\in S\times Q:
                   R_a(s,t)\text{ for some }a\}.
\]
Every mixture satisfies \(\top\), and none satisfies
\(\bot\). The modal clauses are
\[
\begin{aligned}
 M_s\models\langle a\rangle r
   &\iff\exists t\,[R_a(s,t)\land(t,r)\in M_s],\\
 M_s\models[a]r
   &\iff\forall t\,[R_a(s,t)\Rightarrow(t,r)\in M_s],
\end{aligned}
\]
and the usual positive Boolean clauses.
A run starts with label \((s_*,q_0)\).
A node labelled \((s,q)\) selects a mixture
satisfying \(\delta(q,\ell(s))\), and has
one child labelled \((t,r)\) for every
\((t,r)\) in that mixture.
Index children by those pairs, so no
countability assumption on successor sets
is imposed. The run is a set-sized tree
of finite histories of such pairs.
An empty satisfying mixture makes a leaf.
Acceptance requires the parity condition
on every infinite branch; finite branches
are allowed precisely through these
locally satisfied terminal obligations
\cite[\S5.3]{BradfieldStirling2007}.

For example, if
\(\delta(q,\ell(s))=\langle a\rangle q_1
                         \land[c]q_2\)
and
\[
 R_a(s,s_1),\qquad R_b(s,s_2),\qquad
 R_c(s,s_1),\qquad R_c(s,s)
\]
are all the relevant edges, a sufficient
mixture is
\[
 \{(s_1,q_1),(s_1,q_2),(s,q_2)\}.
\]
The same model state \(s_1\) is checked
with two different controls. The \(b\)-edge
creates no obligation for this expression,
while the actual \(c\)-self-loop creates
the obligation at \(s\) itself.

To prove game/run equivalence, first take
an accepting run. During each expression
evaluation, use its chosen mixture:
select a satisfied disjunct, use a
diamond witness in the mixture, and
answer every universal choice with the
corresponding mixture member. Local
Boolean induction supplies each move.
Every continuing play follows a run
branch, while a terminal play is
accepted by its local condition.
This gives a winning game strategy.

Conversely, positional determinacy gives
one strategy winning throughout the
game's winning region. At a winning
\((s,q)\), collect all successor pairs
reachable through its finite expression
evaluation when \(\exists\) follows
that strategy and \(\forall\) makes
arbitrary choices. Induction on the
positive expression shows that this
mixture satisfies \(\delta(q,\ell(s))\).
All collected pairs are winning control
positions. Repeat the construction at
them to build a run.
Every infinite branch expands to an
infinite play of the original strategy,
with exactly the same recurring control
priorities; it is accepting. Finite
branches end at expressions satisfied
by the empty mixture.
Thus the two notions of acceptance
coincide, including dead ends and
arbitrarily branching models.

\paragraph{Boolean operations on modal automata.}
Complementation is explicit. Replace
each priority \(i\) by \(i+1\), and
dualize every transition expression:
\[
\begin{gathered}
 \top\leftrightarrow\bot,\qquad
 \land\leftrightarrow\lor,\qquad
 \langle a\rangle r\leftrightarrow[a]r.
\end{gathered}
\]
The dual game exchanges the players'
choices and terminal winners.
Adding one preserves the order of
priorities and reverses the parity
of the least recurring one.
Determinacy therefore shows that the
dual accepts exactly the rooted
models rejected by the original.

For union or intersection, rename two
control sets disjointly and add a fresh
initial control \(q_*\) with any priority.
Let \(P_i\) be the proposition set of
component \(i\). Use the common set
\(P_1\cup P_2\), restricting each letter to the
propositions used by its component.
At \(q_*\), set
\[
 \delta_\cup(q_*,A)
       =\delta_1(q_1,A\cap P_1)\lor
                         \delta_2(q_2,A\cap P_2),
 \qquad
 \delta_\cap(q_*,A)
       =\delta_1(q_1,A\cap P_1)\land
                         \delta_2(q_2,A\cap P_2).
\]
Keep each component's transitions and
priorities elsewhere. No transition
returns to \(q_*\). Its single initial
priority is therefore irrelevant to
infinite acceptance.
The first expression lets \(\exists\)
choose a component; the second lets
\(\forall\) challenge either component.
They recognize exactly union and
intersection. Inlining the initial
expressions, rather than adding a
bare control-state transition, keeps
the depth-one modal grammar.

\paragraph{Two small acceptance calculations.}
Take one control \(q\), with
\[
 \delta(q,A)=
 \begin{cases}
   [a]q,&P\in A,\\
   \bot,&P\notin A.
 \end{cases}
\]
With priority \(0\), it accepts exactly
the models of \(\nu X.(P\land[a]X)\):
every reached state must satisfy \(P\),
and infinite continuing branches are
accepted. With priority \(1\), the
same transitions accept exactly
\(\mu X.(P\land[a]X)\): all states on
the run still satisfy \(P\), but an
infinite branch now loses. Arbitrarily
long finite branches are not confused
with one infinite branch.
These readings follow directly from
the evaluation-game theorem.

For another example, use controls
\(q_0,q_1\) with priorities \(0,1\),
start at \(q_0\), and set, for either
current control \(q\),
\[
 \delta(q,A)=
 \begin{cases}
  \langle a\rangle q_0,&P\in A,\\
  \langle a\rangle q_1,&P\notin A.
 \end{cases}
\]
Every successful play must make
infinitely many actual \(a\)-moves.
After the initial control, visits
to \(q_0\) correspond exactly to
the preceding model state's
satisfaction of \(P\).
The automaton therefore asks for
an infinite \(a\)-path visiting
\(P\) infinitely often.
Dualizing gives priorities \(1,2\)
and boxes in place of diamonds:
every infinite \(a\)-path must
visit \(P\) only finitely often,
while finite branches are accepted.
The two automata make both the
path quantifier and the recurring
priority condition explicit.

\paragraph{From arbitrary formulas to guarded automata.}
\label{sec:p09-mu-automaton-construction}
Every closed modal \(\mu\)-calculus formula has an
equivalent modal parity automaton of the depth-one
grammar above. In particular, the formula need not
already be guarded. We now give a finite construction
and prove its correctness, including infinite plays
that never take a model edge. This supplies the
formula-to-automaton direction of the expressive
equivalence in \cite[\S5.3, thm.~5.2]{BradfieldStirling2007}.

Put the formula \(\varphi\) in positive normal form
and rename its bound variables distinctly.
Let \(C\) be its finite set of subformula
\emph{occurrences}, retaining their original binder
links. A bound-variable occurrence returns to its
binder, and the binder moves to its body, exactly
as in the evaluation game. Let \(\omega(c)\) be
that game's priority at a control \(c\):
\[
 \omega(X)=
 \begin{cases}
  2d_X+1,&X\text{ is a }\mu\text{-variable},\\
  2d_X,&X\text{ is a }\nu\text{-variable},
 \end{cases}
 \qquad D=2b+2,
\]
where \(b\) is the number of binders and \(d_X\)
counts the proper enclosing binders of \(X\).
Set \(\omega(c)=D\) at every other control,
including terminals. Terminal priorities do not
change their prescribed finite-play outcomes.
Write \(K=\omega(C)\); it is a finite set containing
\(D\), and the minimum of any nonempty subset
of \(K\) belongs to \(K\).

Let \(\mathcal M\subseteq C\) be the modal
occurrences and set
\[
 H=\mathcal M\times K.
\]
An element \((u,k)\) of \(H\) is a \emph{colored
exit}: \(u\) identifies the actual modal occurrence,
and \(k\) records a minimum priority. Distinct
occurrences are kept distinct even when their
printed formulas coincide. Let \(P_0\) be the
finite set of proposition constants of \(\varphi\).

\paragraph{Finite local games and their Boolean summaries.}
For a letter \(A\subseteq P_0\) and a set
\(U\subseteq H\), form a finite game \(L(A,U)\)
with controls \((c,k)\in C\times K\).
It follows only syntax moves at the current
model state:
\begin{itemize}
\item Boolean choices retain their evaluation-game
owners; binder and variable returns are forced.
A syntax move \(c\to c'\) updates the register to
\(\min\{k,\omega(c')\}\).
\item Literals and constants terminate with
their truth values at the letter \(A\).
\item At a modal occurrence \(u\), stop before
taking a model edge. The terminal \((u,k)\)
is won by \(\exists\) exactly when \((u,k)\in U\).
\item An infinite local play is won by the
parity of the least infinitely recurring
\emph{actual control priority} \(\omega(c)\).
The register \(k\) is not its parity priority.
\end{itemize}
Starting at a control \(\psi\) means resetting
the register to \(D\) and immediately reading
\(\omega(\psi)\). Thus the start is
\[
 \iota(\psi)=(\psi,\omega(\psi)).
\]
This includes the first control in the recorded
minimum, even when that control is a variable.
No successor states of a model are consulted
in this local game.

Define its winning function by
\[
 f_{\psi,A}(U)=1
 \quad\Longleftrightarrow\quad
 \exists\text{ wins }L(A,U)\text{ from }\iota(\psi).
\]
Finite parity-game determinacy, with terminal
outcomes totalized if necessary, makes this
an effective finite definition. It is monotone
in \(U\): enlarging \(U\) only changes some
losing modal terminals into winning ones, so
every formerly winning strategy still wins.
Introduce one Boolean variable \(z_h\) per
\(h\in H\), and put
\[
 B_{\psi,A}(z)=
 \bigvee_{\substack{U\subseteq H\\f_{\psi,A}(U)=1}}
       \ \bigwedge_{h\in U}z_h.
\]
An empty disjunction is \(\bot\), and an empty
conjunction is \(\top\). Retain the indexing
of summands by \(U\) and of their atom occurrences
by \(h\). For any truth set \(V\subseteq H\),
\[
 B_{\psi,A}(1_V)=1\quad\Longleftrightarrow\quad
 f_{\psi,A}(V)=1.
\]
Indeed, a true summand gives a winning
\(U\subseteq V\), hence \(f_{\psi,A}(V)=1\)
by monotonicity; conversely, when that value
is one, the summand indexed by \(V\) is true.
This includes constant functions and \(H=\varnothing\).

\paragraph{The transition construction.}
Take
\[
 Q=C\times K,\qquad
 q_I=(\varphi,D),\qquad
 \Omega_Q(\psi,j)=j.
\]
For each colored exit substitute the modal atom
\[
 \rho(u,k)=
 \begin{cases}
  \langle a\rangle(\chi,k),&u=\langle a\rangle\chi,\\
  [a](\chi,k),&u=[a]\chi,
 \end{cases}
\]
where \(\chi\) denotes the particular body
occurrence of \(u\). Define
\[
 \delta((\psi,j),A)=
 B_{\psi,A}[\rho(h)/z_h:h\in H].
\]
All next-state references are now under a
single modality. The construction therefore
produces the required guarded transition
language, even for an unguarded input formula.
The incoming color \(j\) affects the current
control priority but not the transition:
it reports the minimum of the preceding
completed local episode. Each new episode
resets its register and reads \(\omega(\psi)\)
afresh. It does not carry an accumulated
minimum through the rest of the model.

\paragraph{The block-minimum lemma.}
Partition an infinite sequence over a finite
ordered priority set into infinitely many
finite, nonempty consecutive blocks. The
least infinitely recurring priority of the
sequence equals the least infinitely recurring
block minimum. To prove this, call the former
priority \(r\). Values below \(r\) occupy only
finitely many positions and hence only finitely
many blocks. Infinitely many blocks contain
\(r\), since finitely many finite blocks
contain only finitely many occurrences.
Apart from those finitely many exceptional
blocks, every block containing \(r\) has
minimum \(r\). No smaller minimum recurs
infinitely often, which proves the claim.
No uniform bound on block lengths is needed.

For an evaluation play with infinitely many
model moves, a block starts at the control
reached by the preceding model move, or at
\(\varphi\) initially, and ends at the next
modal occurrence just before its model move.
Its target body starts the following block.
If the minima are \(k_0,k_1,\ldots\), the
associated automaton control priorities are
\[
 D,k_0,k_1,\ldots.
\]
The initial \(D\) is a finite prefix.
The two plays consequently have the same
winning parity.

\paragraph{Correctness: from evaluation to acceptance.}
Suppose \(\exists\) wins the evaluation game
from \((s_*,\varphi)\), and choose a positional
winning strategy \(\sigma\). Its continuation
at every \(\sigma\)-reachable position is
winning: a losing continuation prefixed by
its finite reaching history would contradict
the choice of \(\sigma\).

At such a position \((s,\psi)\), restrict
\(\sigma\) to the local Boolean choices at
\(s\). Let \(U_{s,\psi}\) be exactly the set
of colored exits reachable from \(\iota(\psi)\)
by finite local plays following those choices.
Opponent choices at conjunctions remain
unrestricted. This set is finite, although
there may be infinitely many finite paths
realizing one of its members.
The restricted strategy wins
\(L(\ell(s),U_{s,\psi})\): every reached
literal or constant terminal is true, every reached
modal exit belongs to \(U_{s,\psi}\), and
every infinite local play is a winning
evaluation play with the same actual
control priorities. Hence
\(f_{\psi,\ell(s)}(U_{s,\psi})=1\).

In the automaton, choose that summand.
If the support is empty, it is \(\top\)
and the play wins immediately.
Otherwise \(\forall\) may select any
indexed exit \((u,k)\) in the support.
Choose a finite local path realizing
this exit under \(\sigma\).
For a diamond \(u\), use the model
successor selected by \(\sigma\);
for a box, allow any successor chosen
by \(\forall\). A reachable diamond
cannot be empty, and an empty box wins.
In every continuing case the resulting
evaluation position is again reachable
under \(\sigma\), by concatenating the
chosen local path and the model move
with its reaching history.

Repeat this prescription after every
model move. Positionality ensures that
all the chosen finite realizations are
compatible with the same \(\sigma\).
Finite acceptance plays cannot end at
a false constant or a dead diamond.
An infinite acceptance play expands,
by these realizations, to one infinite
\(\sigma\)-compatible evaluation play.
The exit colors are its block minima,
so the block-minimum lemma makes the
acceptance play winning. Winning local
divergences need not be exposed as
automaton branches: they have already
been accounted for in the local
winning function.

\paragraph{Correctness: from acceptance to evaluation.}
Conversely, let \(\tau\) be a winning
acceptance-game strategy from
\((s_*,(\varphi,D))\).
We construct an evaluation strategy
while retaining a virtual
\(\tau\)-compatible acceptance history.
At the start of an episode at
\((s,\psi)\), that history ends at
\((s,(\psi,j))\).
Follow \(\tau\)'s disjunction choices
to select a support \(U\) in
\(B_{\psi,\ell(s)}\).
A winning strategy cannot select
\(\bot\); every selected summand
satisfies \(f_{\psi,\ell(s)}(U)=1\).
Fix a winning local strategy
\(\lambda\) for that finite game
and play it during the actual episode,
resetting the register as specified.
Actual conjunction choices are its
opponent's local choices.

If the episode reaches a literal or
constant, it is true because
\(\lambda\) wins.
If it continues forever without
a modality, its least recurring
actual control priority is even;
the finite preceding history does
not change that outcome.
If it reaches a modal exit \((u,k)\),
then \((u,k)\in U\).
Extend the virtual acceptance history
by letting \(\forall\) select that
indexed conjunct.
At a diamond, use the successor
prescribed there by \(\tau\), which
exists because \(\tau\) is winning.
At a box, use the actual successor
chosen by \(\forall\) also as the
virtual choice. An empty box wins
in both games. Otherwise the histories
reach the matching positions
\[
 (t,\chi)\quad\text{and}\quad(t,(\chi,k)),
\]
and the next episode begins.
If \(U=\varnothing\), no modal exit
is reachable under the winning
\(\lambda\); the first two cases
exhaust its possible behaviors.

These instructions define a strategy
on every reachable finite evaluation
history, extended arbitrarily on
irrelevant nonterminal histories.
Every terminating compatible play
wins by its terminal rule.
An infinite play with finitely many
model moves wins in its last local
episode. A play with infinitely many
model moves yields one infinite
virtual play following \(\tau\);
its parity agrees with the actual
play by the block-minimum lemma.
The evaluation strategy therefore
wins in all cases.
The chosen local supports came from
concrete automaton summands, not from
an assumption about which successor
states are already accepted.

Combining the two strategy constructions
with the truth/game theorem gives
\[
 \mathcal T,s\models\varphi
 \quad\Longleftrightarrow\quad
 \mathcal A_\varphi\text{ accepts }(\mathcal T,s).
\]
The run/game equivalence makes this
also an equivalence with accepting
runs, with no restriction on the
cardinality of successor sets.

\paragraph{Unguarded loops and the two uses of priority.}
For \(\mu X.X\), the local loop is
losing regardless of its exit set,
so its initial transition is \(\bot\).
For \(\nu X.X\), the local loop is
winning even with no permitted exit,
so the summary is \(\top\).
These are the required interpretations
of the two unguarded formulas.

The local parity condition must use
\(\omega(c)\), not the running register.
For example, a play of
\[
 \mu X.\bigl(X\lor\nu Y.Y\bigr)
\]
may choose \(X\) once, return to the
outer binder, then enter the
\(Y\)-loop forever.
Only the even \(Y\)-priority and
the default recur infinitely often,
so that play wins. Its register
would retain the earlier odd
\(X\)-priority forever and give
the wrong answer. The register
summarizes a finite completed
episode; it cannot decide an
episode that never ends.

\paragraph{Finiteness and effective construction.}
Let
\[
 n=|C|,\quad h=|K|,\quad
 m=|\mathcal M|,\quad r=mh,\quad p=|P_0|.
\]
There are \(nh\) automaton states and
at most \(nh\) vertices per local
game, before adding up to two
absorbing terminal sinks.
Each summary has at most \(2^r\)
summands of at most \(r\) atoms.
Thus an explicitly expanded transition
table has size
\[
 O\bigl(nh\,2^p(r+1)2^r\bigr),
\]
apart from the logarithmic encoding
cost of state names. Since
\(h\le b+1\le n+1\), \(m\le n\),
and \(p\le n\), this is a crude
\(2^{O(n^2)}\) finite bound.
Identical transitions for different
incoming colors may be shared.
No optimality claim is needed.

For effectiveness, solve the finitely
many games indexed by \((A,U)\).
One elementary method enumerates
positional \(\exists\)-strategies.
For each strategy and start, retain
its choices and all \(\forall\)-edges.
It wins precisely when every
reachable terminal wins and every
reachable directed simple cycle
has even minimum priority.
An odd-minimum cycle could otherwise
be repeated. Conversely, in a losing
infinite path let \(d\) be the odd
least recurring priority. Discard
the finite prefix containing values
below \(d\). A return to a recurrent
vertex of priority \(d\), with
internal repetitions removed,
gives a reachable simple cycle
of minimum \(d\), a contradiction.
Positional determinacy guarantees
that a winning start has a strategy
found by this finite enumeration.
This establishes an effective
translation without importing an
acceptance oracle or a bound on
local episode lengths.

\paragraph{From modal automata back to formulas.}
\label{sec:p09-mu-automaton-to-formula}
We now prove the reverse direction: every modal
parity automaton \(\mathcal A\) defines a closed
modal \(\mu\)-calculus formula over its original
propositions and action labels. The construction
has two stages. First express the parity game
on a model with an explicit finite control
coordinate. Then eliminate that coordinate by
finite simultaneous fixed points and their
scalar expansion. This supplies the other
direction of \cite[\S5.3, thm.~5.2]{BradfieldStirling2007}.

\paragraph{A formula with explicit finite control.}
Fix
\(\mathcal A=(Q,\mathcal P(P_0),\delta,q_0,\Omega)\),
with \(Q\ne\varnothing\) finite and
\(\Omega:Q\to\{0,\ldots,d\}\).
For a model \(\mathcal T\) with state set \(S\),
form \(\mathcal T^Q\) with state set \(S\times Q\).
Lift an original proposition \(P\) by
\[
 (s,q)\models P\quad\Longleftrightarrow\quad
 s\models P,
\]
and introduce fresh control propositions \(E_r\)
with \((s,q)\models E_r\) exactly when \(q=r\).
For each used action \(a\) and target control \(r\),
introduce a fresh action symbol \(\widehat a_r\)
whose relation is
\[
 R_{\widehat a_r}\bigl((s,q),(t,u)\bigr)
 \quad\Longleftrightarrow\quad
 R_a(s,t)\ \land\ u=r.
\]
It records the automaton's target control,
not an additional choice of a model successor.
In particular, its successor set is in
bijection with the original \(a\)-successor
set, including when that set is empty.
Only finitely many such new symbols are
needed, since the transition table is finite.

For \(A\subseteq P_0\), let
\[
 \chi_A=
   \bigwedge_{P\in A}P\ \land\
   \bigwedge_{P\in P_0\setminus A}\neg P.
\]
Empty conjunctions are \(\top\).
Exactly one \(\chi_A\) is true at each state.
For a transition expression \(B\), define
\(B^{(i)}\) by preserving constants and
Boolean operations and replacing its
modal atoms according to
\[
 (\langle a\rangle r)^{(i)}
       =\langle\widehat a_r\rangle X_i,
 \qquad
 ([a]r)^{(i)}=[\widehat a_r]X_i.
\]
The variables \(X_0,\ldots,X_d\) are fresh.
Put
\[
 \Theta_{\mathcal A}=
 \bigvee_{\substack{q\in Q\\A\subseteq P_0}}
       \bigl(E_q\land\chi_A\land
                       \delta(q,A)^{(\Omega(q))}\bigr),
\]
and define the ordinary scalar formula
\[
 \Psi_{\mathcal A}
   =\sigma_0X_0.\sigma_1X_1.\cdots
                    \sigma_dX_d.\Theta_{\mathcal A},
 \qquad
 \sigma_i=
 \begin{cases}
  \nu,&i\text{ even},\\
  \mu,&i\text{ odd}.
 \end{cases}
\]
All bound-variable occurrences are positive
and occur under a modality. Unused priority
levels merely give vacuous binders.

We claim that, for every \(s\in S\) and \(q\in Q\),
\[
 \mathcal T^Q,(s,q)\models\Psi_{\mathcal A}
 \quad\Longleftrightarrow\quad
 \exists\text{ wins the acceptance game
                   of }\mathcal A\text{ from }(s,q).
\]
For the proof, compare the already defined
evaluation and acceptance games.
At \((s,q)\), the correct disjunct of
\(\Theta_{\mathcal A}\) has indices
\((q,\ell(s))\). Its control and letter
guards are true. Every other disjunct
has a false guard that \(\forall\)
can challenge. A winning evaluation
strategy must therefore select the
correct disjunct whenever this body
is reached.

After its guards, the remaining
Boolean choices and modal moves
are exactly those of
\(\delta(q,\ell(s))\).
A modal move with target \(r\)
goes to \((t,r)\) and encounters
\(X_{\Omega(q)}\). That variable
returns to its binder, traverses
the remaining finite binder prefix,
and reaches \(\Theta_{\mathcal A}\)
again at \((t,r)\).
True and false constants have the
same outcomes in both games.
Dead diamonds lose and dead boxes
win in both games as well.

Every infinite evaluation play
that continues through the guards
makes infinitely many actual model
moves: the transition expressions
and binder prefixes are finite,
and every variable return is
preceded by one modality.
If its acceptance controls are
\(q_0',q_1',\ldots\), its successive
variable returns have indices
\[
 \Omega(q_0'),\Omega(q_1'),\ldots.
\]
The evaluation-game priority at
\(X_i\) is
\[
 \widehat\Omega(X_i)=
 \begin{cases}
   2i,&i\text{ even},\\
   2i+1,&i\text{ odd}.
 \end{cases}
\]
This is strictly increasing in
\(i\) and preserves its parity.
All other recurring evaluation
priorities are larger defaults.
Thus the least recurring priority
has the same parity in both games.
The source control's priority is
recorded after its model move;
this timing does not change the
sequence of recurring colors.

More explicitly, a winning acceptance
strategy yields an evaluation strategy
by selecting the true guards and
copying its choices inside the
transition expression. A guard
challenge terminates with a true
literal or constant; every other play is the
corresponding acceptance play with
the finite syntax moves inserted.
Conversely, from a winning evaluation
strategy maintain one virtual
evaluation history while playing
the acceptance game. The strategy
selects the correct guarded disjunct,
as just proved. In the virtual
history always let \(\forall\)
continue into the transition body,
and copy all actual transition
choices and successors there.
After a model move, insert the
forced variable return and binder
prefix and repeat. All finite
outcomes and infinite parities
agree. This defines a winning
acceptance strategy; no positional
assumption on either translated
strategy is needed.
The truth/game theorem proves
the claimed semantic equivalence.

\paragraph{Eliminating a simultaneous fixed point.}
\label{sec:p09-mu-bekic}
We need an algebraic identity to
remove the finite control coordinate.
Let \(L,M\) be complete lattices
and let
\[
 f:L\times M\to L,\qquad
 g:L\times M\to M
\]
be monotone in both coordinates.
Fix one sign \(\sigma\in\{\mu,\nu\}\)
and define
\[
 h(x)=\sigma y.g(x,y),\qquad
 a=\sigma x.f(x,h(x)),\qquad b=h(a).
\]
Then the simultaneous fixed point is
\[
 \sigma(x,y).(f(x,y),g(x,y))=(a,b).
\]
This is the binary Beki\'c identity.
The product has the componentwise order.
Fixed-point comparison makes \(h\)
monotone, so both displayed scalar
fixed points exist.
Their equations give
\(a=f(a,b)\) and \(b=g(a,b)\).

For \(\sigma=\mu\), take any
pre-fixed pair \((u,v)\), so
\[
 f(u,v)\leq u,\qquad g(u,v)\leq v.
\]
Leastness gives \(h(u)\leq v\).
Consequently
\(f(u,h(u))\leq f(u,v)\leq u\),
and leastness again gives \(a\leq u\).
Finally \(b=h(a)\leq h(u)\leq v\).
Thus \((a,b)\) is below every
pre-fixed pair and is the least
fixed point of the product operator.

For \(\sigma=\nu\), take any
post-fixed pair \((u,v)\), so
\[
 u\leq f(u,v),\qquad v\leq g(u,v).
\]
Greatestness gives \(v\leq h(u)\),
hence \(u\leq f(u,h(u))\) and
then \(u\leq a\).
Monotonicity gives
\(v\leq h(u)\leq h(a)=b\).
Thus \((a,b)\) is above every
post-fixed pair and is the greatest
fixed point. Both proofs permit
additional fixed external parameters.

For \(k\) coordinates, group the
first lattice as \(L\) and the
remaining \(k-1\) as \(M\).
Inductively solve the latter
simultaneous \(\sigma\)-system
with the first variable \(x\)
treated as a parameter, obtaining
\[
 H_2(x),\ldots,H_k(x).
\]
If the first component function
is \(f_1\), set
\[
 a_1=\sigma x.
       f_1\bigl(x,H_2(x),\ldots,H_k(x)\bigr),
 \qquad
 a_j=H_j(a_1)\quad(2\leq j\leq k).
\]
The binary identity proves that
this tuple is the simultaneous
solution. The \(k=1\) case is
the ordinary scalar fixed point,
so induction eliminates every
finite simultaneous block.

At the formula level these are
finite capture-avoiding substitutions.
Rename internal binders freshly
before copying a formula, keep
each external parameter free
until its own surrounding binder
is applied, and substitute only
for the designated coordinate
variables. Positivity in the
coordinate variables is preserved.
This gives finite scalar formulas
for every component, not merely
an existence result about sets.
For nested blocks, eliminate the
inner one with the outer variables
held as parameters, then apply
the same construction outward.
This preserves the nested solution;
it does not commute least and
greatest fixed points or promise
to preserve the literal binder
pattern after formula duplication.
Simultaneous blocks and their
parameter interpretation are the
equation-system presentation of
\cite[\S3.7]{BradfieldStirling2007}.

\paragraph{Removing the control coordinate.}
For a positive-normal formula \(\psi\)
over the propositions \(P_0\cup\{E_r:r\in Q\}\),
the action symbols \(\widehat a_r\),
and fixed-point variables, construct a
tuple of scalar formulas
\((\tau_q(\psi))_{q\in Q}\)
in the original vocabulary.
Each free variable \(X\) receives
a fresh family \((X_r)_{r\in Q}\).
All original bound variables are
renamed distinctly before the
construction, and copied binders
are renamed as just specified.
The atomic and Boolean clauses are
\[
\begin{gathered}
 \tau_q(P)=P,\qquad
 \tau_q(\neg P)=\neg P,\qquad
 \tau_q(X)=X_q,\\
 \tau_q(E_r)=
   \begin{cases}\top,&q=r,\\\bot,&q\ne r,\end{cases}
 \qquad
 \tau_q(\neg E_r)=
   \begin{cases}\bot,&q=r,\\\top,&q\ne r.\end{cases}
\end{gathered}
\]
Constants, conjunction, and
disjunction are translated
componentwise. For modalities put
\[
 \tau_q(\langle\widehat a_r\rangle\psi)
     =\langle a\rangle\tau_r(\psi),
 \qquad
 \tau_q([\widehat a_r]\psi)
     =[a]\tau_r(\psi).
\]
The target component is \(r\),
irrespective of the current
component \(q\), exactly as in
the product relation.

For a binder \(\sigma X.\psi\),
first translate its body with
the fresh \(X_r\) as free
coordinate variables. Take the
simultaneous fixed point
\[
 \sigma(X_r)_{r\in Q}.
                 \bigl(\tau_r(\psi)\bigr)_{r\in Q}
\]
and eliminate this finite block
by the proved identity.
Its scalar \(q\)-component is,
by definition,
\(\tau_q(\sigma X.\psi)\).
This recursive construction
terminates because it translates
strict subformulas before applying
a finite block elimination.
The displayed simultaneous notation
is an intermediate instruction,
not an extra connective in the
output language.

To prove correctness, let \(\eta\)
assign subsets of \(S\times Q\)
to free variables and give the
component environment the values
\[
 \eta^\flat(X_q)
   =\{s:(s,q)\in\eta(X)\}.
\]
We prove by structural induction
that
\[
 \{s:(s,q)\in
       \llbracket\psi\rrbracket_{\mathcal T^Q,\eta}\}
   =\llbracket\tau_q(\psi)\rrbracket_
                                      {\mathcal T,\eta^\flat}.
\]
The literal and Boolean cases
are their definitions.
The modal cases use the exact
bijection between the relevant
successor sets; the universal
case includes empty sets.
For the fixed-point step use
the complete-lattice isomorphism
\[
 J:\mathcal P(S\times Q)\longrightarrow
                       \mathcal P(S)^Q,\qquad
 J(U)_q=\{s:(s,q)\in U\}.
\]
Its inverse sends \((U_q)_q\)
to \(\bigcup_q(U_q\times\{q\})\).
It preserves the order, all
unions, and all intersections.
By the induction hypothesis
the body operator is carried
by \(J\) to the operator with
components \(\tau_q(\psi)\).
Pre-fixed and post-fixed points
correspond under this isomorphism,
so their least and greatest
fixed points correspond.
Beki\'c elimination preserves
each component of precisely
that simultaneous fixed point.
This proves the binder case
and the induction.

Apply the construction to
\(\Psi_{\mathcal A}\), and set
\[
 \varphi_{\mathcal A}
       =\tau_{q_0}(\Psi_{\mathcal A}).
\]
It is closed and uses only
the original propositions and
action labels. We obtain
\[
\begin{aligned}
 \mathcal T,s\models\varphi_{\mathcal A}
 &\quad\Longleftrightarrow\quad
 \mathcal T^Q,(s,q_0)\models\Psi_{\mathcal A}\\
 &\quad\Longleftrightarrow\quad
 \mathcal A\text{ accepts }(\mathcal T,s).
\end{aligned}
\]
Together with the earlier
formula-to-automaton construction,
this proves both
expressive-equivalence directions.
The two proofs apply to arbitrary
set-sized models, without a
finite-branching or totality
assumption. The translation is
effective, but scalar substitution
may duplicate formulas; no
polynomial output-size or
alternation-depth preservation
claim is made.

\paragraph{A direct equation-system description.}
The construction can be recorded without
writing the product model's tags.
For each priority \(i\), use a
family \(\mathbf X_i=(X_{i,r})_{r\in Q}\).
Let \(B^{[i]}\) preserve the
constants and Boolean operations
of a transition expression and put
\[
 (\langle a\rangle r)^{[i]}=\langle a\rangle X_{i,r},
 \qquad
 ([a]r)^{[i]}=[a]X_{i,r}.
\]
Define the component bodies
\[
 F_q(\mathbf X_0,\ldots,\mathbf X_d)=
   \bigvee_{A\subseteq P_0}
      \bigl(\chi_A\land
                    \delta(q,A)^{[\Omega(q)]}\bigr).
\]
Then the tuple of acceptance predicates
is the nested product-lattice expression
\[
 \sigma_0\mathbf X_0.\sigma_1\mathbf X_1.
       \cdots\sigma_d\mathbf X_d.
                      (F_q)_{q\in Q}.
\]
Each binder here binds its whole
\(Q\)-tuple simultaneously with
the indicated single sign.
To verify this description, translate
\(\Theta_{\mathcal A}\) componentwise:
all false control guards disappear,
the true control guard becomes
\(\top\), and the remaining
letter tests give exactly \(F_q\).
The lattice isomorphism transports
each surrounding scalar product
binder to the corresponding tuple
binder. Thus this is the same
nested solution just proved correct,
not a new characterization assumed
from an equation system.
The \(q_0\)-component, scalarized
by Beki\'c elimination, is the
required formula. Outer tuples
remain parameters during each
inner fixed-point calculation.

\paragraph{Two calculations with fixed-point blocks.}
For a single simultaneous
\(\nu\)-block
\[
 x=P\land\langle a\rangle y,\qquad
 y=R\land[b]x,
\]
the second component is independent
of \(y\). Its greatest fixed point
with \(x\) fixed is
\(h(x)=R\land[b]x\).
The identity gives
\[
 a_*=\nu x.\bigl(P\land
              \langle a\rangle(R\land[b]x)\bigr),
 \qquad
 b_*=R\land[b]a_*.
\]
These are the two acceptance
predicates of the all-even
two-control automaton whose first
control requires \(P\) and an
\(a\)-move to the second, and whose
second requires \(R\) and sends
every \(b\)-successor to the first.
A missing \(a\)-successor loses
at the first control; a missing
\(b\)-successor satisfies the
second control's modal obligation.
The scalar formulas retain these
different dead-end outcomes.

Different fixed-point signs
cannot be interchanged.
On the two-state deterministic
cycle
\[
 s\xrightarrow{a}t\xrightarrow{a}s,
 \qquad V(P)=\{s\},
\]
consider
\[
 F(A,B)=
 \bigl(V(P)\cap\langle a\rangle A\bigr)
 \ \cup\
 \bigl((S\setminus V(P))\cap\langle a\rangle B\bigr).
\]
For fixed \(A\), the least
\(B\)-solution is \(S\) if
\(t\in A\) and \(\varnothing\)
otherwise. Its greatest outer
fixed point is therefore \(S\).
For fixed \(B\), the greatest
\(A\)-solution is \(S\) if
\(s\in B\) and \(\varnothing\)
otherwise. Its least outer
fixed point is \(\varnothing\).
Consequently
\[
\begin{aligned}
 \llbracket\nu X.\mu Y.
    ((P\land\langle a\rangle X)
       \lor(\neg P\land\langle a\rangle Y))\rrbracket
     &=S,\\
 \llbracket\mu Y.\nu X.
    ((P\land\langle a\rangle X)
       \lor(\neg P\land\langle a\rangle Y))\rrbracket
     &=\varnothing.
\end{aligned}
\]
The first formula asks for an
infinite \(a\)-path on which
\(P\) recurs infinitely often:
each \(P\)-state forces a return
through \(X\), and each non-\(P\)
state a return through \(Y\).
The cycle satisfies this condition.
The calculation exhibits directly
why elimination must respect the
nested parameter semantics.

\paragraph{Cover modalities.}
\label{sec:p09-mu-cover-modalities}
For a finite set \(\Gamma\) of formulas, define
\[
 \nabla_a\Gamma
   :=\bigwedge_{\psi\in\Gamma}\langle a\rangle\psi
       \ \land\ [a]\bigvee_{\psi\in\Gamma}\psi.
\]
At a state \(s\), every member of \(\Gamma\)
must hold at some \(a\)-successor, and every
\(a\)-successor must satisfy some member
of \(\Gamma\). One successor may witness
several members; different successors may
satisfy different members.
This is the cover constructor written
\(a\to\Gamma\) by Janin--Walukiewicz
and \((a)\Gamma\) by Bradfield--Stirling
\cite[def.~2.5 of the author version]{JaninWalukiewicz1995}.

With the usual empty conjunction and
disjunction conventions,
\[
\begin{gathered}
 \nabla_a\varnothing=[a]\bot,\qquad
 \nabla_a\{\top\}=\langle a\rangle\top,\\
 \langle a\rangle\psi
       \equiv\nabla_a\{\psi,\top\},\qquad
 [a]\psi\equiv
       \nabla_a\varnothing\lor\nabla_a\{\psi\}.
\end{gathered}
\]
The first identity says there is no
\(a\)-successor; the second says there
is at least one. For the diamond
identity its \(\psi\)-witness also
witnesses \(\top\), while the box
condition is automatic. For the
box identity, split into the cases
of an empty or nonempty successor
set. In the latter case a successor
witnesses \(\psi\) as well as every
successor satisfying it. These
arguments prove all four identities,
including dead ends.

The cover notation is a separate
one-step constructor. For example,
\[
 \nabla_a\{q,r\}
   =\langle a\rangle q\land
     \langle a\rangle r\land[a](q\lor r)
\]
contains a Boolean choice inside
the box. It is not literally a
transition expression of the earlier
grammar with only atoms \([a]q\)
and \(\langle a\rangle q\).
We therefore give cover transitions
their own game semantics and prove
the translation between the two
presentations.

\paragraph{Cover-form automata and their games.}
A cover-form parity automaton has
a finite nonempty control set,
the finite proposition alphabet
\(\mathcal P(P_0)\), an initial
control, and a finite priority map
with least-recurring-even acceptance.
Its transition at a control and
letter is a finite disjunction
of clauses
\[
 \bigwedge_{a\in J}\nabla_a\Gamma_a,
 \qquad
 \Gamma_a\subseteq Q\cup\{\top\}\text{ finite},
\]
where \(J\) is a finite set of
distinct action labels. Empty
disjunction is \(\bot\), and
the clause with \(J=\varnothing\)
is \(\top\). An omitted action
imposes no condition.

Player \(\exists\) chooses a clause
and \(\forall\) may challenge any
of its covers. At a cover, \(\forall\)
may either choose a target
\(r\in\Gamma_a\) as a diamond
obligation, or challenge the box
obligation. For the diamond,
\(\exists\) chooses an \(a\)-successor
and play continues there with \(r\).
For the box, \(\forall\) first
chooses an \(a\)-successor, and
\(\exists\) then chooses a member
of \(\Gamma_a\) at that successor.
An empty successor set wins this
box challenge; an empty
\(\Gamma_a\) loses it if a
successor exists.
The target \(\top\) wins immediately
after the model move; a control
target resumes its transition.
An empty diamond loses.
Constants and infinite control
sequences have their already
specified outcomes.

This is the evaluation game of
the displayed cover abbreviation.
In particular, the control selected
in the box branch may depend on
the chosen successor.
For example, if one successor
satisfies only \(q\) and another
only \(r\), then \([a](q\lor r)\)
holds although neither \([a]q\)
nor \([a]r\) holds. Moving that
existential choice before the
successor choice would change
the condition.

If all transitions are to mention
every action in a fixed finite
set \(A_0\), an omitted action
can be filled by the tautology
\[
 \nabla_a\varnothing\lor\nabla_a\{\top\}.
\]
Distribute the finitely many such
disjunctions across the clauses.
Using only \(\nabla_a\{\top\}\)
would instead impose the unwanted
existence of an \(a\)-successor.
The distinction between an action
alphabet and the proposition-letter
alphabet is maintained throughout.

\paragraph{A finite monitor for all obligation traces.}
\label{sec:p09-mu-trace-monitor}
Fix the controls \(Q\), initial
control \(q_0\), and priority
map \(\Omega\) of an ordinary
modal automaton. Let
\[
 \mathcal R=\mathcal P(Q\times Q).
\]
A trace of an infinite relation
word \(\rho_0\rho_1\cdots\)
is a sequence \(q_0,q_1,\ldots\)
with \((q_i,q_{i+1})\in\rho_i\).
Call the word good when every
such infinite trace has even
least recurring priority.
If it has no infinite trace,
it is good.
We construct a finite deterministic
parity word automaton for exactly
these good words.

First recognize the bad words
nondeterministically with a
B\"uchi condition. Let \(J_o\)
be the odd values in \(\Omega(Q)\).
Use search states \((q,*)\)
and committed states \((q,j)\)
with \(j\in J_o\) and
\(\Omega(q)\geq j\).
The initial state is \((q_0,*)\).
On reading \(\rho\), every move
must follow a pair \((q,r)\in\rho\).
From \((q,*)\), it may go to
\((r,*)\), or commit to
\((r,j)\) for any admissible \(j\).
From \((q,j)\), it must stay
committed to \(j\) and go to
\((r,j)\) with \(\Omega(r)\geq j\).
The accepting states are precisely
\[
 \{(q,j):j\in J_o,\ \Omega(q)=j\}.
\]
A run that never commits does
not accept.

An accepting run supplies a trace
whose priorities after commitment
are all at least its fixed odd
\(j\), with \(j\) occurring
infinitely often. That trace is
bad. Conversely, on any bad
trace choose its odd least
recurring value \(j\).
After the last smaller priority,
commit to \(j\) and follow that
trace. This is an accepting
B\"uchi run. Delaying commitment
until after the first input
letter changes no infinite
recurrence condition.
Thus the language is exactly
the bad words, including rejection
when no infinite trace exists.
There are
\[
 |Q|+\sum_{j\in J_o}
             |\{q:\Omega(q)\geq j\}|
 \leq |Q|(1+|J_o|)
\]
states.

Apply the proved B\"uchi-to-Rabin
determinization in
Section~\ref{sec:p09-safra-determinization}.
Its empty-root sink makes the
result total. Let its state set
be \(Z\), initial state \(z_0\),
transition function \(\eta\), and Rabin pairs
\((G_v,B_v)\): a word is bad
exactly when the unique run
visits some \(G_v\) infinitely
often and \(B_v\) only finitely
often. Hence the complementary
good condition is the Muller family
\[
 \mathcal F_{\mathrm{good}}=
 \{I\subseteq Z:
   \text{for every }v,\ 
       I\cap G_v=\varnothing
       \text{ or }I\cap B_v\ne\varnothing\}.
\]
This is a finite, explicitly
decidable family of subsets.

\paragraph{From the Muller condition to least parity.}
\label{sec:p09-muller-parity-record}
Put \(m=|Z|\geq1\).
Keep a permutation
\(\pi=(z_1,\ldots,z_m)\)
in least-recently-visited-first
order, initialized arbitrarily.
After a transition to \(z'\),
find its old position \(k\)
in \(\pi\), form the old suffix
\[
 I_k=\{z_k,z_{k+1},\ldots,z_m\},
\]
and move \(z'\) to the end.
The emitted priority is
\[
 c'=
 \begin{cases}
  2k,&I_k\in\mathcal F_{\mathrm{good}},\\
  2k+1,&I_k\notin\mathcal F_{\mathrm{good}}.
 \end{cases}
\]
Store this priority as part of
the next deterministic state.
Thus a monitor state is
\((z,\pi,c)\), and its priority
is \(c\). Choose an initial
permutation \(\pi_0\) and set
\(z_I=(z_0,\pi_0,0)\).
On input \(\rho\)
first compute \(z'=\eta(z,\rho)\),
then perform the update just given.
There are at most
\(m\,m!\,(2m+2)\) such states.
The record reads the deterministic
run after its initial state;
omitting that state and adding
the initial color change only
finite prefixes.

For correctness, let \(I\ne\varnothing\)
be the infinity set of the
unique \(Z\)-run.
After the last visit outside
\(I\), and after every member
of \(I\) has been visited again,
the permutation has a frozen
prefix \(Z\setminus I\) followed
by exactly \(I\).
Put \(k_*=m-|I|+1\).
No later hit has position
below \(k_*\).
Hits at \(k_*\) occur infinitely
often: otherwise its occupant
would stay the least-recent
member of \(I\) forever without
another visit, contrary to its
membership in \(I\).
At every such hit the suffix
is exactly \(I\).
All hits at larger positions
have priorities greater than
both \(2k_*\) and \(2k_*+1\).
Therefore the least recurring
priority is even exactly when
\(I\in\mathcal F_{\mathrm{good}}\).
This proves the conversion,
with no assumption that the
input has an infinite obligation
trace.

Denote the resulting total
deterministic parity monitor by
\[
 \mathcal D=(Z_D,\mathcal R,\zeta,z_I,\lambda).
\]
It accepts precisely the good
relation words. Its state and
priority after a relation are
functions of the preceding state
and that relation, so the same
finite information can be copied
to every branch of a modal run.
The construction is finite and
effective; its factorial record
bound is not an optimal-size
claim.

\paragraph{Converting an arbitrary modal automaton to cover form.}
\label{sec:p09-mu-cover-simulation}
We now use the monitor to construct
an equivalent cover-form automaton.
Fix the finite set \(A_0\) of
actions occurring in the original
transition table. A new control
will contain a set \(D\subseteq Q\)
of simultaneous obligations, a
monitor state \(z\in Z_D\), and
the relation that produced them.
The monitor is the one just
constructed with initial trace
control \(q_0\).

First put each positive Boolean
transition expression into a finite
disjunction of conjunctions of
modal atoms. This is an exact
finite operation: associate with
\(\top\) the family \(\{\varnothing\}\),
with \(\bot\) the empty family,
and with an atom \(\alpha\) the
family \(\{\{\alpha\}\}\).
Disjunction takes the union of
families; conjunction takes all
unions \(U\cup V\) of a term
from each family.
Thus each term \(U\) is a set
of atoms \(\langle a\rangle r\)
or \([a]r\).
Keep a derivation of each term
from the Boolean expression;
it specifies local existential
choices whose possible modal
exits belong to that term.
Conversely, fixing all existential
choices in this finite expression
and keeping all universal choices
produces such a term, provided
no false terminal is reached.

At a new control with obligation
set \(D\ne\varnothing\), and
current letter \(A\), choose
one term \(U_q\) of \(\delta(q,A)\)
for every \(q\in D\).
If any of these term families
is empty, there is no such
choice and no resulting clause.
For \(a\in A_0\), define two
relations in \(D\times Q\):
\[
\begin{aligned}
 B_a&=\{(q,r):q\in D,\ [a]r\in U_q\},\\
 I_a&=\{(q,r):q\in D,\ \langle a\rangle r\in U_q\}.
\end{aligned}
\]
The pairs in \(B_a\) must be
realized at every \(a\)-successor.
Each pair in \(I_a\) must be
realized at some \(a\)-successor.
Choose a finite family
\(\mathcal C_a\) of relations
satisfying
\[
 B_a\subseteq\rho\subseteq B_a\cup I_a
       \quad(\rho\in\mathcal C_a),\qquad
 I_a\subseteq\bigcup_{\rho\in\mathcal C_a}\rho.
\]
The family may be empty only
when \(I_a=\varnothing\).
The empty relation itself is
allowed as a family member.
All possibilities are finite
and can be enumerated without
consulting a model.

For the current monitor state \(z\),
give a relation \(\rho\) the
target control
\[
 k(z,\rho)=
       \bigl(\operatorname{ran}\rho,\zeta(z,\rho),\rho\bigr).
\]
Include the clause
\[
 \bigwedge_{a\in A_0}
       \nabla_a\{k(z,\rho):\rho\in\mathcal C_a\}
\]
for every admissible choice of
the terms and families.
The transition is their disjunction.
Retain the indices of the choices
on these disjunct occurrences,
even if two printed clauses coincide.
At an empty obligation set take
the transition \(\top\).
The initial control is
\[
 k_I=(\{q_0\},z_I,\star),
\]
where \(\star\notin\mathcal R\)
is an initial marker. Every
control \((D,z,\rho)\) has priority
\(\lambda(z)\).
The relation component distinguishes
two family members even if their
ranges and monitor successors
coincide; it supplies an exact
record of the chosen obligation
links. A crude control bound is
\[
 2^{|Q|}|Z_D|\bigl(1+2^{|Q|^2}\bigr).
\]
Unreachable controls can be discarded.

\paragraph{Why the cover simulation accepts every original model.}
Suppose the original automaton has
a winning strategy from \((s_0,q_0)\).
By positional determinacy choose
one positional winning strategy
\(\sigma\). Its continuations
at every reachable control
position are winning.
At a reached new position
\((s,D,z,\rho_{\mathrm{old}})\),
maintain that every \((s,q)\),
\(q\in D\), is reachable under
\(\sigma\).
For each such \(q\), the finite
expression strategy supplies a
term \(U_q\) as above.
For each diamond pair
\((q,r)\in I_a\), fix one
reachable occurrence of that
atom and its \(\sigma\)-chosen
successor \(t_{q,r}\).
These successors exist because
the strategy is winning.

For every actual \(a\)-successor
\(t\) of \(s\), set
\[
 \rho_t=B_a\cup
       \{(q,r)\in I_a:t=t_{q,r}\}.
\]
Take
\(\mathcal C_a=\{\rho_t:R_a(s,t)\}\).
Although the successor set may
be arbitrarily large, this family
is finite because \(\mathcal R\)
is finite. It satisfies the
required inclusions, and each
diamond pair appears at its
chosen witness.
At a diamond challenge to a
family member \(\rho\), choose
one successor \(t\) with
\(\rho_t=\rho\).
At a box challenge, after
\(\forall\) chooses \(t\),
choose the target \(k(z,\rho_t)\).
This satisfies each cover exactly.
When there is no successor the
family is empty and its cover
is the correct \([a]\bot\).

Every pair \((q,r)\in\rho_t\)
is realizable by a finite
\(\sigma\)-compatible expression
path and its model move.
For a pair in \(B_a\), use
a reachable box occurrence
and let \(\forall\) choose \(t\).
For a pair in \(\rho_t\setminus B_a\),
the definition makes it a
diamond pair with \(t=t_{q,r}\);
use its chosen diamond occurrence.
Hence every member of the new
obligation set
\(\operatorname{ran}\rho_t\)
is again reachable under
\(\sigma\). If that set is
empty, the next transition
accepts immediately.

Consider an infinite play of
this cover strategy. Its chosen
relations form a word
\(\rho_0\rho_1\cdots\).
Any trace of that word from
\(q_0\) expands, by the finite
realizations just described,
to one infinite play following
\(\sigma\). Positionality ensures
that realizations chosen after
different earlier histories
remain compatible with the
same strategy.
Such a trace has even least
recurring original priority.
Therefore the relation word
is good, and the deterministic
monitor gives the cover play
even least recurring priority.
Every finite play also wins
by the cover and terminal
clauses. This proves acceptance
by the constructed automaton.

\paragraph{Why a cover strategy satisfies the original obligations.}
Conversely, let \(\tau\) be a
winning strategy for the
constructed cover automaton.
While playing the original
game, maintain a virtual
\(\tau\)-compatible cover history
at the same model state, with
the current original control
\(q\) belonging to its set \(D\).
This holds initially.
The clause chosen by \(\tau\)
records a term \(U_q\) for
each \(q\in D\).
In the actual Boolean transition
at \(q\), follow the finite
existential choices supplied
by the derivation of its term.
Any reached modal atom belongs
to \(U_q\); a constant terminal
is true.

If the actual challenge is
\(\langle a\rangle r\), choose
a relation \(\rho\in\mathcal C_a\)
with \((q,r)\in\rho\).
Such a relation exists by the
diamond coverage condition.
In the virtual history let
\(\forall\) select the cover
for \(a\) and its diamond
obligation \(k(z,\rho)\).
Use the successor selected
there by \(\tau\) as the
actual diamond successor.
It exists because \(\tau\)
is winning. The actual new
control \(r\) lies in
\(\operatorname{ran}\rho\).

If the actual challenge is
\([a]r\), an empty successor
set wins immediately.
Otherwise let the actual
\(\forall\)-chosen successor
\(t\) also be its virtual
choice in the cover's box
branch. The strategy \(\tau\)
chooses some \(\rho\in\mathcal C_a\)
there. Every such relation
contains \(B_a\), so it
contains \((q,r)\), and again
\(r\in\operatorname{ran}\rho\).
An empty family with an
actual successor would lose
the virtual box challenge,
so a winning \(\tau\) cannot
encounter that case.
In either continuing case
extend both histories and repeat.

These rules preserve the
membership invariant and give
an original strategy on all
reachable histories.
They also give a trace from
\(q_0\) along the virtual
relation word.
Finite actual plays win by
the preceding terminal cases.
An infinite actual play makes
infinitely many model moves
and produces one infinite
virtual cover play following
\(\tau\). Its monitor accepts,
so every trace of that word,
including the actual control
sequence, has even least
recurring priority.
The original strategy wins.
This proves equivalence in
each original pointed model,
not merely equisatisfiability
or equivalence after adding
copies of model states.

\paragraph{The automaton normal-form grammar.}
\label{sec:p09-mu-anf-grammar}
Write \(\chi_L\) for the conjunction
of a finite set \(L\) of
proposition literals \(P,\neg P\),
with \(\chi_\varnothing=\top\).
Automaton normal form, abbreviated
ANF, is the following positive
cover grammar:
\begin{itemize}
\item constants, proposition literals,
and fixed-point variables belong
to ANF;
\item disjunctions and least or
greatest fixed-point binders
preserve ANF;
\item if \(I\) is finite,
the \(a_i\) are pairwise distinct,
and each \(\Gamma_i\) is a
finite set of ANF formulas, then
\[
 \chi_L\land\bigwedge_{i\in I}
                            \nabla_{a_i}\Gamma_i
\]
belongs to ANF, including
\(I=\varnothing\).
\end{itemize}
There is no closure clause
for arbitrary conjunction.
In particular, a bound variable
cannot be inserted into the
proposition-literal part of a
cover clause.
The guarded, well-named members
of this grammar will be the
output normal forms
\cite[def.~4.1]{JaninWalukiewicz1995}.
A binder is guarded when every
occurrence it binds is below
a modality in its own body.
All occurrences in a cover's
target formulas are below its
modalities. Guarding outside
a binder does not make its
body guarded: a cover containing
\(\nu X.X\) still contains
an unguarded binder.

ANF is closed under capture-avoiding
substitution of ANF formulas
for free fixed-point variables.
To prove this, rename bound
variables freshly and induct
on the grammar. Constants
and proposition literals are
unchanged; a designated variable
is replaced by its ANF formula.
Disjunctions and fresh binders
use the induction hypothesis.
At a cover clause only the
target formulas are substituted:
their finite image is still
ANF and the action labels
remain distinct. Repeated
targets can be removed by
Boolean idempotence.
No arbitrary conjunction is
introduced in any case.

It follows that the Beki\'c
construction preserves ANF for
a finite same-sign block of
ANF component bodies. In the
binary step its formulas are
\[
 H(x)=\sigma y.\beta(x,y),\qquad
 A=\sigma x.\alpha(x,H(x)),\qquad
 B=H(A),
\]
with any external parameters
left free. Each formula uses
only an allowed binder and
an allowed substitution.
The finite-coordinate induction
proved above applies the same
operations repeatedly.

For guarded output, keep a
stronger invariant: every
existing binder is guarded,
and every free occurrence
of each coordinate being
solved is guarded in every
component body.
In the binary step, the
new \(y\)-binder in \(H\)
is then guarded, and its
free \(x\) remains guarded.
Substituting \(H\) at the
guarded \(y\)-positions in
\(\alpha\) preserves guarding
of \(x\), so its new binder
in \(A\) is guarded.
The substitution giving \(B\)
preserves all existing and
copied guarded binders because
fresh names prevent capture.
Any external parameter guarded
throughout the original bodies
likewise remains guarded.
Grouping coordinates proves
the same invariant inductively
for every finite block.
Without the free-coordinate
condition the one-component
body \(X\) would give
\(\nu X.X\), so ANF membership
alone would not prove this
stronger conclusion.

\paragraph{From cover automata to guarded ANF.}
\label{sec:p09-mu-cover-to-anf}
Let \(\mathcal C\) be an
already given cover-form
automaton, with controls \(Q\)
and priorities \(0,\ldots,d\).
We supply its formula construction
with the successor-first box
choice preserved.
Use the product state set
\(S\times Q\), lifted propositions
and control tags \(E_r\).
For each real action \(a\)
and each control \(r\), use
fresh actions \(a^\circ,h_r\)
with
\[
\begin{aligned}
 R_{a^\circ}((s,q),(t,u))
     &\iff R_a(s,t)\land u=q,\\
 R_{h_r}((s,q),(t,u))
     &\iff t=s\land u=r.
\end{aligned}
\]
The real move preserves the
current control; the deterministic
total \(h_r\)-move changes only
that control to \(r\).
These are the action symbols
of this product model.

For a priority index \(i\),
replace a cover target \(r\in Q\)
by \(\langle h_r\rangle X_i\),
and leave the target \(\top\)
as \(\top\).
Replace its action \(a\) by
\(a^\circ\). Let
\(\delta_{\mathcal C}(q,A)^{\circ i}\)
be the resulting cover expression.
With the exact letter cubes
\(\chi_A\), form
\[
\begin{aligned}
 \Theta_{\mathcal C}
   &=\bigvee_{\substack{q\in Q\\A\subseteq P_0}}
       \bigl(E_q\land\chi_A\land
          \delta_{\mathcal C}(q,A)^{\circ\Omega(q)}\bigr),\\
 \Psi_{\mathcal C}
   &=\sigma_0X_0.\cdots\sigma_dX_d.\Theta_{\mathcal C},
\end{aligned}
\]
where \(\sigma_i=\nu\) for
even \(i\), and \(\mu\) for
odd \(i\).
Expanding covers makes this
an ordinary scalar formula
on the product model.

The guard argument is the
same exact state-and-letter
selection proved for
\(\Psi_{\mathcal A}\).
In a diamond challenge for
\(r\), its real move reaches
\((t,q)\), then the unique
\(h_r\)-move reaches \((t,r)\).
In a box challenge, \(\forall\)
first chooses \((t,q)\);
only then does \(\exists\)
choose the cover member
and its control switch.
Thus the choice of \(r\)
can depend on the successor.
A \(\top\)-target terminates
after the real move, and
control switches introduce
neither a new choice nor
a dead end.
Empty covers, empty clauses,
false transitions, and missing
real successors retain their
specified finite outcomes.

A cover strategy gives a
formula strategy by selecting
the true guards and copying
its choices, with the unique
control switches inserted.
Conversely, keep a virtual
formula history, continue past
true guard challenges into
the transition, and use its
successor-dependent target
choices in the actual cover
game. Insert each control
switch and forced variable
return in that virtual
history before the next round.
Both constructions may keep
the full preceding history.
All extra terminal challenges
are immediate true literal
or constant outcomes.

On an infinite continuing
play the variable returns
are exactly \(X_{\Omega(q)}\)
for successive source controls
\(q\). Their priorities are
the same strictly increasing,
parity-preserving encoding
\(i\mapsto2i+(i\bmod2)\)
proved above. Every return
has a preceding real move,
so no infinite administrative
play is added.
The two strategies therefore
preserve winning on every
finite or infinite play.
By the truth/game theorem,
\(\Psi_{\mathcal C}\) expresses
cover acceptance on the product.

\paragraph{ANF component bodies and scalarization.}
For a target set \(\Gamma\)
and priority \(i\), define
\[
 \Gamma_i^\flat
   =\{X_{i,r}:r\in\Gamma\cap Q\}
      \ \cup\
      \begin{cases}
        \{\top\},&\top\in\Gamma,\\
        \varnothing,&\top\notin\Gamma.
      \end{cases}
\]
Write clause \(j\) of
\(\delta_{\mathcal C}(q,A)\)
as
\(\bigwedge_{a\in J_{q,A,j}}
          \nabla_a\Gamma_{q,A,j,a}\).
Its componentwise product
translation yields the base
body
\[
 F_q^{d+1}=
 \bigvee_{A,j}
 \left(\chi_A\land
    \bigwedge_{a\in J_{q,A,j}}
      \nabla_a(\Gamma_{q,A,j,a})_{\Omega(q)}^\flat\right).
\]
Each summand is an admitted
ANF clause: its guard consists
only of proposition literals,
its action labels are distinct,
and its targets are variables
or \(\top\).
Every free coordinate variable
is guarded, and no binder is
yet present.

For completeness, the component
translation removes the two
product actions by the clauses
\[
\begin{aligned}
 \tau_q(\langle a^\circ\rangle\psi)
       &=\langle a\rangle\tau_q(\psi),&
 \tau_q([a^\circ]\psi)
       &=[a]\tau_q(\psi),\\
 \tau_q(\langle h_r\rangle\psi)
       &=\tau_r(\psi),&
 \tau_q([h_r]\psi)
       &=\tau_r(\psi).
\end{aligned}
\]
The second line follows from
the unique, always present
control-switch successor.
The tag and variable clauses
are as before. In particular,
the translated cover target
\(\langle h_r\rangle X_i\)
is \(X_{i,r}\), underneath
the surviving real modality.
Finite distribution over the
chosen clauses and removal
of false control guards
give exactly the displayed
ANF bodies. This step does
not conjoin arbitrary ANFs.

Now for \(i=d,d-1,\ldots,0\),
take the simultaneous block
\[
 \sigma_i(X_{i,r})_{r\in Q}.
                      (F_r^{i+1})_{r\in Q}
\]
and let \(F_q^i\) be its
scalar \(q\)-component produced
by Beki\'c elimination.
All lower-indexed coordinate
families remain free parameters
at this stage.
Induction downward on \(i\)
shows that these formulas
are ANF, that their binders
are guarded, and that every
remaining free coordinate
parameter is guarded.
Fresh renaming preserves those
properties and makes the
final formula well-named.

The semantic induction uses
the complete-lattice isomorphism
\(\mathcal P(S\times Q)\cong\mathcal P(S)^Q\).
At the base, the just-proved
component clauses identify
the denotation of
\(\Theta_{\mathcal C}\) for
every outer environment.
At step \(i\), varying the
product value of \(X_i\)
varies precisely the tuple
of its coordinate values.
The body induction identifies
the two monotone operators;
the isomorphism carries their
least or greatest fixed points
to each other. Beki\'c then
preserves every scalar component.
This proves the induction
uniformly in all still-free
outer parameters.
Consequently
\[
 \varphi_{\mathcal C}=F_{q_0}^{0}
\]
is closed, guarded, well-named,
and ANF over the original
propositions and actions, with
\[
 \mathcal T,s\models\varphi_{\mathcal C}
 \quad\Longleftrightarrow\quad
 \mathcal C\text{ accepts }(\mathcal T,s).
\]

\paragraph{The semantic normal-form theorem.}
\label{sec:p09-mu-semantic-normal-form}
Combining the constructions gives
the effective transformation
\[
 \varphi
 \longmapsto\mathcal A_\varphi
 \longmapsto\mathcal C_\varphi
 \longmapsto\varphi_{\mathcal C_\varphi}.
\]
The first step admits unguarded
closed input formulas; the
second uses the deterministic
all-traces monitor and the
proved cover simulation; the
third is the guarded ANF
scalarization just established.
Every step preserves truth
in the same original pointed
model. We have therefore
proved that every closed
modal \(\mu\)-calculus formula
has an effectively constructible
equivalent guarded, well-named
ANF formula
\cite[thm.~4.4]{JaninWalukiewicz1995}.
No finite-branching or
seriality hypothesis is needed.
The explicit monitor and scalar
copying establish finiteness,
not an optimal normal-form
size bound.
This is semantic equivalence;
an axiomatization argument
requires the separate statement
that the equivalences are
derivable in its calculus.

\paragraph{Cutting a guarded normal form.}
\label{sec:p09-mu-anf-cut}
Let \(\varphi\) be closed,
guarded, well-named ANF.
Define \(\operatorname{cut}\varphi\)
on its original occurrence tree:
replace each occurrence of a
\(\mu\)-bound variable by \(\bot\),
each occurrence of a \(\nu\)-bound
variable by \(\top\), and
erase the now-vacuous binders.
Keep constants, proposition literals,
disjunctions, and cover clauses,
recursively cutting their targets.
Each variable is classified by
its original binder before that
binder is erased. Thus an outer
variable occurring inside an
inner fixed point keeps its
outer binding classification.
This operation does not replace
whole fixed-point subformulas
by constants: for example,
\(\operatorname{cut}(\mu X.P)=P\).

The result is closed, fixed-point-free
ANF. We prove the exact reduction
\[
 \operatorname{Sat}(\varphi)
 \quad\Longleftrightarrow\quad
 \operatorname{Sat}(\operatorname{cut}\varphi),
\]
where \(\operatorname{Sat}\) means
truth at some state of some model.
This is the normal-form
equisatisfiability theorem
\cite[thm.~4.6]{JaninWalukiewicz1995}.
Its two models need not be
the same.

\paragraph{The finite local satisfiability calculation.}
For a closed fixed-point-free
ANF formula \(\alpha\), define
a Boolean value \(b(\alpha)\) by
\[
\begin{gathered}
 b(\top)=1,\qquad b(\bot)=0,\qquad
 b(P)=b(\neg P)=1,\\
 b(\alpha\lor\beta)
       =b(\alpha)\lor b(\beta),\\
 b\left(\chi_L\land
            \bigwedge_{i\in I}\nabla_{a_i}\Gamma_i\right)
       =\operatorname{Cons}(L)\land
           \bigwedge_{i\in I}\bigwedge_{\psi\in\Gamma_i}b(\psi).
\end{gathered}
\]
Here \(\operatorname{Cons}(L)=1\)
exactly when \(L\) contains
no pair \(P,\neg P\), and
an empty Boolean conjunction
has value \(1\).
The value \(b(P)=1\) means
that \(P\) can be satisfied
in a chosen model, not that
\(P\) is valid.

We prove by structural induction
that \(b(\alpha)=1\) exactly
when \(\alpha\) is satisfiable.
Constants and literals have
the stated one-state witnesses;
disjunction selects a satisfiable
disjunct. For a cover clause,
any model makes its literals
consistent and supplies a
witness for every target
formula, proving necessity.
Conversely, choose a pointed
model for each target and
take disjoint copies of these
models. Add a fresh root
satisfying \(L\), and for
each \(a_i\) give it exactly
the edges to the chosen
roots for members of
\(\Gamma_i\).
Each target has a witness,
and every new \(a_i\)-successor
satisfies some target of
that same cover.
The action labels are distinct,
so the demands of different
covers do not interfere.
An empty target set gives
no edge of its action.
Adding incoming edges does
not change truth within
the forward-modal witness
models. This proves sufficiency.
The induction gives finite
tree witnesses as well.

\paragraph{A finite game on normal-form occurrences.}
\label{sec:p09-mu-anf-occurrence-game}
Keep covers as primitive
constructors and distinguish
their target occurrences.
Form a finite game
\(G_\varphi\) on the original
ANF occurrence tree with
variable-to-binder return edges.
At a disjunction, \(\exists\)
chooses a child; at a binder,
move to its body; at a
variable, return to its
original binding occurrence.
At a clause with inconsistent
\(L\), \(\exists\) loses.
At a consistent clause,
\(\forall\) chooses any pair
\[
 (a_i,\psi),\qquad \psi\in\Gamma_i,
\]
and play moves to that
target occurrence. If no
such pair exists, \(\exists\)
wins. This last case includes
clauses containing only empty
covers. Standalone literals
and \(\top\) win; \(\bot\)
loses.

Give variable occurrences their
earlier evaluation-game priorities:
an outer binder has a lower
priority than a properly inner
binder, with odd parity for
\(\mu\) and even for \(\nu\).
All other nonterminal priorities
are larger defaults.
Finite terminal outcomes are
as just specified.
This is a finite parity game,
so its winning strategies
may be chosen positional.
It is a game of possible
model construction, not an
evaluation game on one
already specified model.

If \(\varphi\) has a model,
\(\exists\) wins \(G_\varphi\).
To see this, start with a
winning evaluation strategy
in that model and retain
its full virtual history,
including the current model
state. Copy its disjunction
choices and forced syntax
moves. At a cover-target
challenge \((a,\psi)\),
let the virtual \(\forall\)
challenge the mandatory
diamond \(\langle a\rangle\psi\).
The winning evaluation strategy
supplies a successor there;
use that state for the
next virtual position.
A reached literal clause
must be consistent, since
every one of its literals
can be challenged in the
evaluation game.

Every infinite resulting
occurrence play is realized
by one compatible evaluation
play with the same recurring
variable priorities. Finite
outcomes likewise win.
When cover expansion duplicates
a target formula, keep its
chosen evaluation-copy occurrence
in the virtual history and
project it to the original
primitive target occurrence.
Fresh renaming of copied
binders preserves their kinds
and nesting depths.
This handles repeated cover
targets without changing their
original binding links.
The induced strategy need
not be positional; finite
positional determinacy supplies
a positional winning strategy
on \(G_\varphi\) afterward.

\paragraph{Why satisfiability survives the cut.}
Choose such a positional
winning strategy \(\sigma\).
In the original finite syntax
tree retain its chosen child
at each reached disjunction,
the body of each binder,
and every target of each
reached cover clause.
For the moment regard all
variable occurrences as leaves,
without their return edges.
Call the resulting finite
selected tree \(T_\sigma\).

No false constant or
inconsistent clause is reached.
Suppose a reached variable
leaf \(X\) is \(\mu\)-bound.
Its original binder \(b_X\)
is an ancestor on its
root-to-leaf syntax path.
There is no earlier variable
occurrence on that path,
because variables are leaves
of this tree.
Restore the return edge
\(X\to b_X\) and repeat
the path from \(b_X\)
to \(X\).
Every repeated existential
choice is the same because
\(\sigma\) is positional;
\(\forall\) can repeat its
cover-target choices.
The only recurring variable
on this play is \(X\),
which has odd priority.
Intermediate binders have
only default priority.
This would be a losing
\(\sigma\)-compatible play,
a contradiction.

Thus every reached variable
leaf is \(\nu\)-bound.
After cutting, all such
leaves are true constants;
the selected finite tree
contains no losing leaf
or inconsistent clause.
Backward induction gives
\[
 b(\operatorname{cut}\varphi)=1.
\]
The fixed-point-free calculation
therefore gives a model of
the cut. The argument applies
also when \(X\)'s binder is
outside several alternating
inner binders: only variable
returns, not binder entries,
determine the repeated cycle.
The positional strategy used
here lives on formula occurrences
and need not keep the choices
of the original satisfying
model.

\paragraph{Restoring greatest-fixed-point returns.}
\label{sec:p09-mu-anf-cut-model}
Conversely suppose
\(b(\operatorname{cut}\varphi)=1\).
Work on the original syntax
tree, evaluating each node
with the cut value of
its subformula and retaining
one true child at each
reached disjunction.
Retain every cover target
and every binder body.
This gives a finite selected
tree \(T\) with consistent
literal clauses and no
reached \(\mu\)-variable leaf.
All reached variable leaves
are \(\nu\)-bound.
Restore their return edges
to their original ancestor
binders, obtaining a finite
directed graph \(H\).
Those ancestors already lie
in \(T\), so restoration
introduces no new syntax
occurrences.

Call a clause-to-target edge
a real modal edge and label
it by that target's action.
Other edges are administrative.
The administrative graph is
acyclic. Indeed, consider a
simple directed cycle in a
tree with ancestor return
edges, starting at its
shallowest node.
Its first return edge must
go back to that node:
a return to another ancestor
would repeat an internal
node before completing the
simple cycle.
Thus the cycle consists of
a downward syntax path from
a binder to one of its
variables and that variable's
return. Guardedness places
a real modal edge on every
such path. No cycle remains
after real modal edges
are removed.

After the disjunction choices
are fixed, each administrative
node has one successor.
Following administrative edges
from any reached occurrence
\(e\) therefore terminates
at a unique endpoint \(h(e)\):
a consistent clause, a
standalone literal, or \(\top\).
Clause endpoints include
those with no targets.
No false endpoint is
reachable.

Construct a model \(M_T\)
whose states are these
endpoint occurrences.
Its root is \(h(\varphi)\).
Regard a standalone literal
as its singleton literal
set and \(\top\) as the
empty literal set.
At endpoint \(c\), make
exactly the proposition
constants occurring positively
in its literal set \(L_c\)
true. Consistency of \(L_c\)
makes every negative literal
there true as well.
For each clause endpoint
\[
 c:\ \chi_{L_c}\land
          \bigwedge_{i\in I}\nabla_{a_i}\Gamma_i,
\]
put precisely the edges
\[
 c\xrightarrow{a_i}h(\psi)
       \quad(\psi\in\Gamma_i).
\]
Standalone literal and
\(\top\) endpoints have
no outgoing edges.
Duplicate labelled edges
may be identified.
No other edges are added.

\paragraph{Truth in the reconstructed finite model.}
We give a winning evaluation
strategy for \(\varphi\) at
the root of \(M_T\).
During an administrative
segment, follow the fixed
disjunction choices in \(T\);
binder and variable moves
are forced. The model state
stays equal to \(h(e)\)
for the current primitive
syntax occurrence \(e\),
because \(h\) is unchanged
along an administrative edge.
This includes each restored
\(\nu\)-variable return.
When the segment reaches
its endpoint \(c\), all
literal challenges win by
the chosen valuation.

For a diamond challenge
to a cover member \(\psi\),
choose the edge to \(h(\psi)\).
For a box challenge,
\(\forall\) chooses an
actual edge
\(c\xrightarrow{a_i}d\).
By construction there is
a member \(\psi\in\Gamma_i\)
with \(h(\psi)=d\).
Choose that member in
the cover's disjunction
at \(d\), and resume
the strategy at \(\psi\).
Distinct action labels in
the clause ensure that
an \(a_i\)-edge comes
from this cover, not
from another cover with
incompatible targets.
If duplicate edges were
identified, choose any
corresponding target member.
An empty cover has no
matching edges, so its
box challenge wins; all
required diamonds have
their displayed witnesses.

Every continuing evaluation
play projects to a path
in \(H\), including the
original binder links of
any expanded target copies.
No \(\mu\)-variable is
ever returned through.
An infinite play must
have a recurring variable,
since all other moves
descend finite syntax.
Its recurring variables
are therefore all
\(\nu\)-bound, and its
least recurring priority
is even.
Finite plays end only
at true literals, true
constants, or vacuous
box obligations.
The strategy wins in
all cases. The evaluation
theorem proves
\[
 M_T,h(\varphi)\models\varphi.
\]
This construction proves
cut sufficiency with a
finite model and with
no least-fixed-point
regeneration in its
selected evaluation paths.

\paragraph{Size bounds and the decision procedure.}
Let \(n\) be the size of
the original ANF occurrence
tree in cover notation,
counting its literal lists
and target occurrences.
The model states are a
subset of its endpoint
occurrences, so there are
at most \(n\) states.
There is at most one
constructed labelled edge
per original cover-target
occurrence, before duplicate
edges are identified.
Hence the edge count
is at most \(n\) as well.
These are bounds in the
given ANF representation,
not in the size of a
possibly much shorter
formula before normal-form
conversion.

The satisfiability decision
does not need to construct
a parity game or a model.
Compute one Boolean value
per original syntax node
from the leaves upward:
variables receive the
constant prescribed by
their original binder kind,
binders pass through their
body's value, and the
remaining constructors use
the displayed rules for \(b\).
No copied unfolding or
expanded cover is needed.
The result is
\[
 \operatorname{Sat}(\varphi)
 \quad\Longleftrightarrow\quad
 b(\operatorname{cut}\varphi)=1.
\]
This proves the linear
satisfiability consequence
\cite[cor.~4.7]{JaninWalukiewicz1995}
with its representation
convention explicit.

More precisely, if proposition
symbols are supplied as integer
identifiers in \(\{1,\ldots,n\}\),
and variable occurrences carry
their original binder links
or precomputed binder-kind tags,
each literal list is checked
in time proportional to its
length. Use positive and
negative marks with a
clause identifier, so a
new clause does not require
clearing an entire array.
Each constructor, literal,
and cover target is then
processed a constant number
of times, giving \(O(n)\)
time and \(O(n)\) space.
For unprocessed string names,
include their symbol-table
construction cost separately.
The input is assumed to
be the stated guarded,
well-named ANF; the claim
does not assign unit cost
to conversion of arbitrary
formulas into that syntax.

\paragraph{The finite-model and decidability consequences.}
\label{sec:p09-mu-finite-model-decision}
For an arbitrary closed
modal \(\mu\)-calculus formula
\(\theta\), first construct
its equivalent guarded ANF
\(\varphi\) by the semantic
normal-form theorem.
The Boolean cut calculation
decides its satisfiability.
When the result is true,
the finite model \(M_T\)
satisfies \(\varphi\), and
same-model equivalence gives
\(M_T,h(\varphi)\models\theta\).
Thus the full modal
\(\mu\)-calculus has the
finite-model property and
decidable satisfiability.
The linear decision bound
applies to the supplied
ANF, not automatically
to the original \(\theta\);
the preceding normal-form
construction makes no
optimal complexity claim.

\paragraph{Examples and the role of the hypotheses.}
For
\(\varphi=\nu X.\nabla_a\{X\}\),
the cut is
\(\nabla_a\{\top\}\).
Its Boolean value is \(1\),
and restoration gives one
clause state with an
\(a\)-self-loop, satisfying
\(\varphi\).
A root with one edge
to a dead end satisfies
the cut but not \(\varphi\).
Thus the theorem is
equisatisfiability, not
truth equivalence in
each fixed model.
For \(\mu X.\nabla_a\{X\}\),
the cut instead contains
the impossible target
\(\bot\), so the Boolean
value is \(0\).

Re-entering a least-fixed-point
binder is different from
returning through its variable.
For example,
\[
 \nu X.\mu Y.
       \bigl(\nabla_a\{X\}\lor\nabla_b\{Y\}\bigr)
\]
has a satisfiable cut, selecting
the \(a\)-cover.
Its reconstructed model has
one \(a\)-loop and no
\(b\)-edge. Each return
through \(X\) passes through
the inner \(\mu Y\)-binder
again, but no return through
\(Y\) occurs. The recurring
variable priority is still
the even priority of \(X\).

The restriction on bound
variables in literal
conjunctions is essential.
The formula
\[
 \nu Y.\bigl(P\land
                  \nabla_a\{Y\land\neg P\}\bigr)
\]
is outside the ANF grammar.
Any value of its greatest
fixed point lies in \(P\);
its required target would
then lie in both \(P\)
and \(\neg P\), so that
fixed point is empty.
Its putative cut,
\(P\land\nabla_a\{\neg P\}\),
is nevertheless satisfiable.

The distinct-action condition
has a separate purpose.
The formula
\[
 \nu Y.\bigl(\neg P\land
              \nabla_a\{P\}\land\nabla_a\{Y\}\bigr)
\]
uses two covers of the
same action and is also
outside ANF. A greatest
fixed-point set would
lie in \(\neg P\), but
its nonempty successor
set would have to lie
in both \(P\) and that
fixed-point set. It
is therefore unsatisfiable.
Replacing \(Y\) by \(\top\)
produces a satisfiable cut.
These examples show why
the local witness construction
cannot be applied to
arbitrary conjunctions.

\paragraph{The three alternation conventions.}
\label{sec:p09-mu-alternation-classes}
Use the formula classes of
Definition~\ref{def:p09-mu-calculus}.
In this discussion abbreviate
\(\Sigma_n^{T\mu},\Pi_n^{T\mu}\)
as \(\Sigma_n^T,\Pi_n^T\),
where \(T\in\{S,EL,N\}\).
The difference between the
three conventions is the
substitution rule: none for
the simple classes, closed
substitution for Emerson--Lei,
and capture-avoiding open
substitution for Niwi\'nski.
These are syntactically
generated classes, not
classes already closed
under semantic equivalence.

Induction on the level and
on its finite generation
derivation gives
\[
\begin{gathered}
 \Sigma_n^S\subseteq\Sigma_n^{EL}\subseteq\Sigma_n^N,
 \qquad
 \Pi_n^S\subseteq\Pi_n^{EL}\subseteq\Pi_n^N,\\
 \Sigma_n^T\cup\Pi_n^T
      \subseteq\Sigma_{n+1}^T\cap\Pi_{n+1}^T.
\end{gathered}
\]
Every simple rule is an
Emerson--Lei rule, and a
closed substitution has
no free variable to capture,
so it is a Niwi\'nski rule.
The last inclusion is
the inherited-generator
clause of each next level.

Let \(s_T(A)\) and \(p_T(A)\)
be the least levels containing
\(A\) on the two sides.
The symmetric depth defined
in the system card is
\[
 d_T(A)=
   \max\{0,\max(s_T(A),p_T(A))-1\}.
\]
In particular, fixed-point-free
formulas have depth zero,
whereas a nonempty pure
least or greatest block
has depth one.
Vacuous binders still count
syntactically: \(\mu X.P\)
has depth one although it
is equivalent to \(P\).
Such simplification concerns
semantic rather than
syntactic alternation.

\paragraph{Duality and the first level.}
The positive dual \(A^\partial\)
exchanges constants, proposition
literal polarities, conjunction
and disjunction, diamond and
box, and \(\mu\) and \(\nu\);
it retains fixed-point-variable
names. It preserves free-variable
sets and commutes with
capture-avoiding substitution,
up to fresh bound names.
Induction on class generation
therefore proves
\[
 A\in\Sigma_n^T
 \quad\Longleftrightarrow\quad
 A^\partial\in\Pi_n^T.
\]
Applying the dual twice gives
the reverse implication.
For open formulas the
semantic identity is
\[
 \llbracket A^\partial\rrbracket_\eta
    =S\setminus\llbracket A\rrbracket_{\eta^c},
 \qquad \eta^c(X)=S\setminus\eta(X).
\]
For closed formulas this
is ordinary complementation.

At level one, all three
\(\Sigma\)-classes consist
exactly of formulas whose
binders are all \(\mu\);
the three \(\Pi\)-classes
consist exactly of formulas
whose binders are all \(\nu\).
The base contains no binders,
and none of the same-level
operations can introduce
the opposite kind.
Conversely, structural induction
constructs every formula
having only the allowed
kind. This is a useful
lower bound independent of
any more refined dependency
analysis.

\paragraph{Exact simple-class indices.}
The pair of minimum simple
indices is computed structurally.
Constants, literals, and
variables have pair \((0,0)\).
Boolean operations take the
componentwise maximum of
their children's pairs;
modalities preserve the pair.
For binders,
\[
\begin{aligned}
 (s_S(\mu X.A),p_S(\mu X.A))
     &=(u,u+1),&u&=\max\{1,s_S(A)\},\\
 (s_S(\nu X.A),p_S(\nu X.A))
     &=(v+1,v),&v&=\max\{1,p_S(A)\}.
\end{aligned}
\]
To prove the lower bounds,
read the binder-sign word
on each syntax branch and
compress adjacent equal
signs. If a nonempty word
has \(r\) blocks, its
\(\Sigma\)-cost is \(r\)
when it begins with \(\mu\)
and \(r+1\) when it begins
with \(\nu\); the
\(\Pi\)-cost is dual.
The empty word has both
costs zero.
The displayed recursion
computes the maxima of
these costs over all
branches.

Every generation rule of
the simple class preserves
the appropriate bound.
Same-sign prefixing preserves
its side's cost, apart
from raising an empty
word's cost to one.
Moving from either preceding
level to the next is valid
because the two costs differ
by at most one.
Thus a larger computed cost
proves nonmembership.
Conversely, the same recursion
constructs the child formulas
at their computed levels,
adds each binder on its
own side, and uses the
next-level inclusion for
the opposite side.
This proves that the
recursion gives the actual
minimum indices.

\paragraph{Closed fragments for Emerson--Lei.}
Split a formula at its
proper closed subformulas.
In each piece replace
maximal proper closed
subformulas by distinct
fresh variable holes, and
repeat inside each removed
piece. The current piece
is not removed from itself.
This produces finitely many
fragments; each recursive
call concerns a strict
subformula.
A reference to an enclosing
fixed-point variable makes
a subformula open when
read in isolation.

For either side, membership
at Emerson--Lei level \(n\)
is equivalent to every
fragment having simple
membership on that same
side at level \(n\).
At level zero, both conditions
mean that no fixed-point
binder occurs. At positive
levels, for one direction reconstruct
the pieces from the bottom
up by permitted closed
substitutions. The simple
fragments already belong
to the Emerson--Lei class,
and each inserted piece
is closed.
For the other, induct on
Emerson--Lei generation.
A closed substitution
introduces a boundary at
each inserted piece and
may create further closed
boundaries in the context;
it cannot open a previously
closed subformula.
The remaining fragment
binder words are therefore
contiguous subwords of old
context-fragment or inserted-piece-fragment
words. Their costs do
not increase. Boolean,
modal, binder, and inherited
level rules preserve the
bound as above.
This proves the characterization,
including its nonmembership
direction.

\paragraph{A capture-stable dependency lower bound.}
For Niwi\'nski nonmembership,
use clean binder names and
put an edge \(b_X\to b_Y\)
when \(b_Y\) is properly
inside \(b_X\) and
\[
 X\in\operatorname{FV}(b_Y).
\]
Read the sign word on
any directed chain of
this finite acyclic binder
graph, including singleton
chains. Its \(\Sigma\)-
and \(\Pi\)-costs are
defined as above.
Membership in \(\Sigma_n^N\)
or \(\Pi_n^N\) respectively
bounds all these costs
by \(n\).

The proof again follows
the actual generation rules.
A new root binder can
only stand alone or prefix
its sign to a pre-existing
chain. Boolean and modal
composition join no binder
dependencies between separate
children. In a capture-safe
substitution, an inserted
free variable cannot acquire
a context binder.
Consequently no new dependency
edge crosses from a context
binder into an inserted
binding occurrence.
An inserted binder cannot
enclose a context binder,
and no new dependency
between old context binders
is created by such a
captured occurrence either.
Every resulting chain
therefore stays in the
context or in one inserted
copy. The premise bounds
give the result.
This necessary condition
suffices for the lower
bounds below; no complete
Niwi\'nski classification
algorithm is inferred from
this one invariant.

\paragraph{Three exact class calculations.}
Define
\[
\begin{aligned}
 E&=\mu Z.(P\lor\langle a\rangle Z),\\
 A_1&=\nu Y.(E\land\langle a\rangle Y),\\
 A_2&=\nu Y.\mu Z.
          ((P\lor\langle a\rangle Z)\land\langle a\rangle Y),\\
 A_3&=\mu X.\nu Y.
          ([-]Y\land\mu Z.[-](X\lor Z)).
\end{aligned}
\]
Here \([-]\) abbreviates
the box over a fixed
finite action alphabet containing \(a\);
expanding it adds no
fixed-point binder.
These are the three
illustrative patterns in
\cite[\S7]{BradfieldStirling2007}.

For \(A_1\), the simple
recursion gives \((3,2)\).
The branch \(\nu Y,\mu Z\)
has \(\Sigma\)-cost three,
so the upper bound cannot
be lowered to \(\Sigma_2^S\).
For Emerson--Lei, put
\[
 F(W)=\nu Y.(W\land\langle a\rangle Y).
\]
Then \(F(W)\in\Pi_1^{EL}
\subseteq\Sigma_2^{EL}\)
and \(E\in\Sigma_1^{EL}
\subseteq\Sigma_2^{EL}\).
Substituting the closed
\(E\) for \(W\) proves
\(A_1\in\Sigma_2^{EL}\);
its simple \(\Pi_2\)
derivation still applies.
The same holds for \(N\).
Both binder kinds occur,
so neither level-one side
contains it. Thus the
two minimum indices for
EL and \(N\) are \((2,2)\).

For \(A_2\), its inner
\(\mu Z\)-formula is
\(\Sigma_1^S\) with \(Y\)
free. Inherit it into
\(\Pi_2^S\) and add
the outer \(\nu Y\).
This gives \(\Pi_2\)
membership in all three
conventions.
But the dependency edge
\(b_Y\to b_Z\) has
word \(\nu\mu\) and
\(\Sigma\)-cost three.
Hence \(A_2\notin\Sigma_2^N\),
and therefore not in
\(\Sigma_2^{EL}\) or
\(\Sigma_2^S\).
Its two minimum indices
are \((3,2)\) throughout.
Substituting the open
inner formula into an
already formed \(\nu Y.W\)
would capture its \(Y\).
Renaming that binder freshly
avoids capture but leaves
this \(Y\) free; it
does not produce \(A_2\).

For \(A_3\), the simple
word \(\mu\nu\mu\) gives
the exact pair \((3,4)\).
It has no proper closed
subformula: both inner
fixed-point formulas have
\(X\) free.
The Emerson--Lei fragment
criterion gives the same
pair. For Niwi\'nski,
let
\[
 H(W)=\nu Y.([-]Y\land W),\qquad
 B(X)=\mu Z.[-](X\lor Z).
\]
Both belong to \(\Sigma_2^N\):
the first is already
\(\Pi_1^N\), and the
second \(\Sigma_1^N\).
The only free variable
of \(B(X)\) is \(X\),
so substitution into
\(H(W)\) does not capture
anything. Then adding
the outer \(\mu X\)
keeps \(\Sigma_2^N\).
The order matters:
perform the open substitution
before introducing the
outer binder for \(X\).
The dependency chain
\(b_X\to b_Y\) has
\(\Pi\)-cost three,
so \(A_3\notin\Pi_2^N\).
The exact pair is \((2,3)\).

The resulting minimum-index
and depth triples
\((s_T,p_T,d_T)\) are
\[
\begin{array}{c|ccc}
       &S&EL&N\\ \hline
 A_1&(3,2,2)&(2,2,1)&(2,2,1)\\
 A_2&(3,2,2)&(3,2,2)&(3,2,2)\\
 A_3&(3,4,3)&(3,4,3)&(2,3,2).
\end{array}
\]
These are proofs about
the displayed syntax.
They do not by themselves
exclude an equivalent
formula of smaller class.

The first two formulas
also have different meanings.
Take a state \(s\) with
an \(a\)-self-loop and
an \(a\)-edge to a dead
state \(t\), with \(P\)
true only at \(t\).
Then \(E\) holds at
both states, and \(A_1\)
holds at \(s\).
For any proposed value
of \(Y\), the first
least-\(Z\) approximant
in \(A_2\) is empty:
the only \(P\)-state
has no successor.
Its least fixed point
and hence \(A_2\) are
empty in this model.
This calculation separates
closed eventuality from
the dependent recurrence
requirement without confusing
it with a hierarchy
strictness proof.

\paragraph{Semantic strictness and binary-tree witnesses.}
\label{sec:p09-mu-alternation-strictness-proof}
We now prove the semantic
separation: for every \(n\geq1\)
there is a closed formula
in \(\Sigma_n^S\) equivalent
to no \(\Pi_n^N\) formula.
This stronger choice of
witness proves strictness
for all three conventions.
The proof uses binary trees,
a fixed-index translation
of the Niwi\'nski generation
rules, and a game-tree
diagonal construction
\cite[\S\S2--3]{Arnold1999BinaryHierarchy}.
The index-preservation step
is proved explicitly below.

For a finite alphabet \(B\),
let
\[
 \mathcal T_B=\{t:\{1,2\}^*\to B\}.
\]
Its members are full infinite
binary labelled trees.
At a node \(u\), action
\(i\in\{1,2\}\) has the
unique successor \(ui\).
Root propositions test the
current alphabet symbol.
Additional proposition names
are interpreted as false
and additional action relations
as empty on this tree class.
Before applying the binary
translation, replace their
literals by the corresponding
constants, their boxes by
\(\top\), and their diamonds
by \(\bot\).
This restriction preserves
each alternation class:
it commutes with binders
and capture-avoiding substitution,
and replaces erased parts
by level-zero constants.
Write \(t\equiv_k t'\)
when the two trees agree
at every node of depth
less than \(k\).
These finite-prefix agreements
define the usual prefix
topology. A convergent
sequence eventually agrees
with its limit at every
fixed finite depth.

\paragraph{The first level.}
Every \(\Pi_1^N\) formula
is \(\nu\)-only, as proved
above. Its language of
binary trees is closed.
Here is a proof that
includes arbitrary nesting
of greatest fixed points.
Use the canonical transition
model whose states are all
trees in \(\mathcal T_B\),
with actions taking left
and right subtrees.
Each subtree map is
continuous: agreement through
depth \(k+1\) gives agreement
of the subtrees through
depth \(k\).
Root-literal sets are clopen.

Induct on a \(\nu\)-only
formula, assigning closed
sets of trees to its
free variable parameters.
Finite Boolean unions and
intersections preserve closed
sets. Each modality is
a continuous preimage,
because its direction has
one successor.
For a greatest binder,
start its descending
approximants at the whole
space. The body induction
preserves closedness at
successor stages, and
arbitrary intersections
preserve it at limit
stages. Stabilization gives
a closed greatest fixed
point. This proves the
claim for every closed
\(\Pi_1^N\) formula.
The canonical state \(t\)
and the root of the
ordinary tree model of
\(t\) are bisimilar:
map each node to its
rooted subtree. Thus this
closed set is exactly
the formula's tree language.

On the alphabet recording
whether \(P\) holds, take
\[
 W_1=\mu X.
       (P\lor\langle1\rangle X\lor\langle2\rangle X).
\]
It belongs to \(\Sigma_1^S\)
and holds precisely when
some node carries \(P\).
Indeed, every pre-fixed set
for its body contains \(P\)
and all its finite directional
predecessors; those finite
predecessors themselves form
a fixed set.
Let \(t_m\) have its
only \(P\)-label at
node \(1^m\).
Every \(t_m\) satisfies
\(W_1\), but the sequence
converges to the all-\(\neg P\)
tree, which does not.
Its language is not
closed, so no \(\Pi_1^N\)
formula defines it.

\paragraph{Open recursion graphs at a fixed index.}
\label{sec:p09-mu-niwinski-index}
For \(n\geq1\), formulas
in \(\Sigma_n^N\) have
finite open game graphs
using priorities in
\[
 I_n^\Sigma=\{1,\ldots,n\},
\]
and formulas in \(\Pi_n^N\)
have such graphs using
\[
 I_n^\Pi=\{0,\ldots,n-1\}.
\]
Their games use least-recurring-even
acceptance and the same
finite terminal outcomes
as before.
An open graph has a
designated entry and finitely
many controls of the
following kinds: true or
false terminal, proposition
literal terminal, free-variable
exit, binary existential
or universal choice, forced
edge, diamond, or box.
A modal control specifies
an action and one target
control. At a model state
it makes the corresponding
modal move; a free exit
\(Z\) is true exactly
at states in \(\eta(Z)\).
No free exit has an
internal return edge.

The invariant is semantic
for every model and
every free-variable environment:
the graph's existential
winning set is the formula's
denotation, and its complement
is the universal winning
set. Prove it simultaneously
for both class sides by
induction on the level
and the finite generation
derivation.
At level zero, ordinary
fixed-point-free syntax gives
an acyclic control graph
with the required exits.
When admitted into a
positive level, its priorities
may all be that interval's
maximum, since every control
path terminates after finite
syntax descent.

For inherited generators,
the same-side preceding
interval is already included.
To embed \(\Pi_{n-1}^N\)
into \(\Sigma_n^N\), add
two to its priorities:
\(\{0,\ldots,n-2\}\)
becomes \(\{2,\ldots,n\}\).
This preserves order and
parity. The interval for
\(\Sigma_{n-1}^N\) already
fits inside that for
\(\Pi_n^N\) unchanged.
At \(n=1\), the inherited
level-zero graphs are
handled by the acyclic
case just stated.

Boolean and modal constructors
use disjoint copies of
their operand graphs and
one appropriate new entry.
Assign the interval maximum
to this finite administration.
A play enters its chosen
operand and does not return
to the new entry.
The ordinary Boolean,
modal, and terminal semantics
therefore give the invariant.

For a substitution \(A[B/Z]\),
take disjoint graphs for
\(A\) and \(B\), and
redirect only the designated
free \(Z\)-exits of the
host graph into the
entry of \(B\).
The inserted graph's
own free \(Z\)-exits
remain open, as do
its other free exits.
Capture avoidance ensures
that none acquires a
return to a host binder.
There is no internal
edge from the inserted
graph back to the host.
Evaluate the host with
\(Z\) assigned the
denotation of \(B\).
At its first \(Z\)-exit,
the winning player switches
permanently to a winning
strategy for \(B\) at
that state. If no
such exit occurs, the
host play is unchanged.
This works for both
players and proves exactly
the substitution semantics.
It also explains why
redirecting the inserted
graph's own exits again
would be incorrect.

\paragraph{Binding at the bottom priority of an interval.}
To introduce \(\mu Z.A\)
on the \(\Sigma\) side,
create a forced regeneration
control of priority \(1\),
leading to the old body
entry. Make it the new
entry and redirect all
body exits for \(Z\)
to it. For \(\nu Z.A\)
on the \(\Pi\) side,
use priority \(0\) instead.
Other free exits remain
unchanged. The new priority
is no larger than any
body priority.

Let
\(F(U)=\llbracket A\rrbracket_{\eta[Z:=U]}\).
It is monotone by positivity.
For a true least fixed
point, use the first
successor approximant in
which the current state
enters. The body has a
winning strategy when \(Z\)
is valued at the preceding
approximant. Each reached
\(Z\)-exit then has a
strictly smaller entry
ordinal, so restart the
body with that smaller
ordinal. Infinitely many
such returns are impossible.
An infinite final body
episode wins by the
body invariant.

For a state outside
the least fixed point
\(U=\mu F\), use a
universal winning body
strategy with \(Z:=U\).
Every reached \(Z\)-exit
is still outside \(U\).
If returns occur infinitely
often, their minimum priority
\(1\) makes the play
universally winning; otherwise
the final body episode
does so.
For a state in
\(U=\nu F\), the
existential body strategy
with \(Z:=U\) preserves
membership on every return.
Infinitely many returns
now carry the minimum
even priority \(0\);
finitely many leave a
winning body tail.
Finally, for a state
outside \(\nu F\),
use its first successor
exclusion stage in the
descending approximants.
A losing-for-\(\exists\)
\(Z\)-exit has a
strictly smaller exclusion
stage. Thus a universal
strategy can force only
finitely many returns,
then a winning body tail.
Finite terminal outcomes
in all four cases are
the body strategy's outcomes.
This proves both binding
rules and completes the
fixed-index graph induction.

\paragraph{Removing internal moves without adding an index.}
The local-summary construction
of Section~\ref{sec:p09-mu-automaton-construction}
applies to these finite
graphs with their assigned
priority interval \(I\).
We record the index-sensitive
details. Put \(M=\max I\).
At the start of each
local episode reset the
register to \(M\), then
read the starting control's
actual priority.
Track minima during forced
and Boolean moves; stop
at a modal control and
record that occurrence
and its minimum \(j\in I\).
Infinite local plays use
their actual control
priorities, not the
accumulated minimum.

For an open graph the local
letter records proposition values
and, for each free exit \(Z\),
whether the current state belongs
to \(\eta(Z)\).
Free exits are terminals with
these recorded outcomes.
For a closed graph there
is no extra variable coordinate.
For every such current letter,
the local winning function
on the finite set of
colored modal exits is
monotone. Its exact positive
disjunction of winning
supports is therefore
defined as before.
Replace each colored exit
by its modality with
target \((q,j)\).
Use controls \(C\times I\),
priority \(j\) at \((q,j)\),
and initial control
\((q_I,M)\).
Every output priority
remains in \(I\).
Transitions ignore the
incoming color after
it has been emitted.

The previously proved two
strategy constructions need
only these finite control
games, not an original
syntax-tree priority assignment.
A positional winning graph
strategy supplies exact
reachable supports and
finite realizations of
every selected exit.
Conversely, the support
selected by an automaton
strategy supplies an
independent winning local
strategy; a virtual acceptance
history connects successive
episodes. Local divergence,
finite terminals, and
modal dead ends retain
their proved outcomes.
On infinitely many completed
episodes, the block-minimum
lemma identifies their
recurring parity with that
of the original play.
Using \(M\) rather than
a new larger default is
sufficient: every infinite
output play visits controls
infinitely often, so extra
copies of the interval
maximum do not alter its
least recurring priority.
This proves index-preserving
elimination of internal moves.

For a closed \(\Sigma_n^N\)
formula on full binary
trees, identify either
modality of direction \(i\)
with the atom \((q,i)\):
there is exactly one
successor in that direction.
For \(n\geq2\), replace
true and false transition
constants by transitions
to fresh absorbing controls
of priorities \(2\) and
\(1\), respectively.
Each such control follows
direction \(1\) forever.
Entering it changes a
finite terminal outcome
into an infinite play
of the same winner.
All transitions are now
nonempty positive Boolean
formulas over \(Q\times\{1,2\}\),
and all priorities lie
in \(\{1,\ldots,n\}\).
Thus every closed
\(\Sigma_n^N\) formula,
for \(n\geq2\), has
a binary-tree automaton
of this exact index.
The separate first-level
proof avoids requiring
both terminal parities
inside a one-priority
interval.

\paragraph{Game formulas at the separating levels.}
For \(n\geq2\), let
\[
 B_n=\{c_1,d_1,\ldots,c_n,d_n\}.
\]
The propositions \(C_i,D_i\)
indicate exactly these
symbols on trees in
\(\mathcal T_{B_n}\).
At a \(c_i\)-node,
\(\forall\) chooses its
next direction; at a
\(d_i\)-node, \(\exists\)
chooses. The latter wins
when the least node
index occurring infinitely
often is even.
Denote its winning-tree
language by \(\mathcal W_n\).
The formula
\[
\begin{split}
 W_n={}&\mu X_1.\nu X_2.\cdots\sigma_nX_n.\\
 &\bigvee_{i=1}^n
 \left[
   \bigl(C_i\land([1]X_i\land[2]X_i)\bigr)
   \lor
   \bigl(D_i\land(\langle1\rangle X_i
                         \lor\langle2\rangle X_i)\bigr)
 \right],
\end{split}
\]
where \(\sigma_i\) is
\(\mu\) for odd \(i\)
and \(\nu\) for even
\(i\), defines
\(\mathcal W_n\).
Its matching literal
selects the actual node
symbol; the next Boolean
choice has the specified
owner; and the modal
move is followed by
return through \(X_i\).
The variable priorities
are an increasing,
parity-preserving encoding
of \(i\).
This gives the game
equivalence exactly as
in the earlier parity
formula construction.
The simple-class recurrence
puts \(W_n\) in
\(\Sigma_n^S\).

\paragraph{Encoding acceptance as a binary game tree.}
Let \(\mathcal A\) be
a finite binary-tree
automaton with priorities
\(r(q)\in\{1,\ldots,n\}\)
and no constant transitions.
View each \(\delta(q,b)\)
as a finite binary
Boolean tree with leaves
\((q',i)\).
Replace it first by
\[
 \delta(q,b)\lor\delta(q,b).
\]
Choosing between two identical
copies changes no accepted
tree. It ensures that
every transition has a
Boolean root.

Define \(G_q(t)\in\mathcal T_{B_n}\)
recursively. In the padded
transition tree for
\((q,t(\epsilon))\),
replace each conjunction
node by \(c_{r(q)}\),
each disjunction node by
\(d_{r(q)}\), and each
leaf \((q',i)\) by
\(G_{q'}(t_i)\), where
\(t_i\) is the corresponding
input subtree.
The padded root is
\(d_{r(q)}\), independently
of the current input
letter, so the construction
produces an output node
before making a recursive
call.

An acceptance strategy of
\(\mathcal A\) on \(t\)
becomes a strategy in
the game on \(G_q(t)\)
by copying its Boolean
choices. Conversely, retain
the finite unfolding
history of this construction
and copy each game-tree
choice into the corresponding
transition evaluation.
Every infinite play traverses
infinitely many finite,
nonempty transition blocks.
Within each block all
game-tree node indices
are the source priority
\(r(q)\).
Thus a priority occurs
infinitely often in the
automaton play exactly
when it occurs infinitely
often in the expanded
game-tree play.
Both strategy translations
preserve the winner, giving
\[
 t\in L(\mathcal A,q)
 \quad\Longleftrightarrow\quad
 G_q(t)\in\mathcal W_n.
\]

\paragraph{The prefix fixed point and diagonal contradiction.}
For every automaton state
\(q\), the transformation
has the prefix-gain property
\[
 t\equiv_k t'
 \quad\Longrightarrow\quad
 G_q(t)\equiv_{k+1}G_q(t').
\]
At \(k=0\), the padded
root is independent of
the input letter.
For the induction step,
the input roots agree,
so their padded transition
trees agree. Their input
subtrees agree through
one fewer level; by
the induction hypothesis,
the recursively inserted
output subtrees agree
through \(k\) levels.
Every inserted subtree
starts at depth at
least one. The whole
outputs therefore agree
through \(k+1\) levels.

When the input alphabet
is \(B_n\), \(G_q\)
is a self-map of
\(\mathcal T_{B_n}\).
Start with any tree
\(t_0\), and put
\(t_{m+1}=G_q(t_m)\).
The prefix-gain property
gives
\(t_m\equiv_m t_{m+1}\).
Every finite prefix hence
eventually stabilizes,
defining a tree \(t_*\).
For each \(k\), sufficiently
late iterates agree with
\(t_*\) through depth \(k\);
applying prefix gain shows
that \(G_q(t_*)\) agrees
with \(t_*\) through
depth \(k+1\).
Therefore
\[
 G_q(t_*)=t_*.
\]
The same prefix argument
shows uniqueness, although
existence alone is needed.

Suppose a closed
\(\Pi_n^N\) formula
defined \(\mathcal W_n\).
Its positive dual is
\(\Sigma_n^N\) and
defines the complementary
language on these trees.
By the fixed-index
translation there is
an automaton \(\mathcal A\)
with priorities \(1,\ldots,n\)
for that complement.
At the fixed tree of
its initial-state map,
the acceptance equivalence
gives
\[
\begin{aligned}
 t_*\notin\mathcal W_n
 &\quad\Longleftrightarrow\quad t_*\in L(\mathcal A,q_I)\\
 &\quad\Longleftrightarrow\quad G_{q_I}(t_*)\in\mathcal W_n\\
 &\quad\Longleftrightarrow\quad t_*\in\mathcal W_n,
\end{aligned}
\]
a contradiction.
Thus \(W_n\), already
in \(\Sigma_n^S\),
is equivalent to no
\(\Pi_n^N\) formula.
Together with the separate
\(n=1\) argument, this
proves every required
level of strictness.
Dualizing supplies the
opposite-sided witnesses.
An open formula claimed
equivalent to a closed
witness for every valuation
would give a closed competitor
by replacing all its free
variables by \(\bot\).
That substitution is allowed
at the same Niwi\'nski level,
so open competitors cannot
avoid the separation.

\paragraph{Other conventions and finite transition systems.}
The witnesses were chosen
in the simple class,
so they also belong
to the corresponding
Emerson--Lei and Niwi\'nski
\(\Sigma\)-classes.
Since every opposite-sided
simple or Emerson--Lei
formula belongs to the
opposite Niwi\'nski class,
the same semantic exclusion
applies to it.
This is the additional
witness information needed
to transfer strictness;
class inclusions alone
would not supply it.

Finally, suppose an
opposite-class closed
formula \(B\) agreed
with one of these
witnesses \(W\) on
every finite transition
system. Their difference
formula
\[
 (W\land B^\partial)
     \lor(W^\partial\land B)
\]
is satisfiable, since
the binary-tree argument
excludes equivalence on
all systems.
The full finite-model
property already proved
supplies a finite model
of this difference,
contradicting the supposed
finite-system agreement.
No alternation-class bound
on the difference formula
is needed.
Each proposed \(B\)
may require its own
finite distinguishing
model; no single finite
model for all competitors
is asserted.
This proves the finite-system
form of
Theorem~\ref{thm:p09-mu-alternation}
as well.
